% Statistiques : séries regroupées en classes — Mathématiques 1ère D — Cours signé OnBuch+
% Programme officiel MINESEC. Compiler avec : tectonic N09-statistiques-classes.tex
\newcommand{\DOCMATIERE}{Mathématiques}
\newcommand{\DOCNIVEAU}{1ère D}
\newcommand{\DOCMODULE}{Organisation et gestion des données}
\newcommand{\DOCLECON}{N09}
\newcommand{\DOCTITRE}{Statistiques : séries regroupées en classes}
% OnBuch+ — préambule « cours signés » ENRICHI (compilé avec tectonic / XeLaTeX)
% Système visuel « Pop » d'OnBuch+ : fond crème (jamais blanc), bordures encre
% épaisses, ombres « dures » décalées, titres Archivo Black en capitales.
% Version MATHÉMATIQUES : graphes (pgfplots/tikz), boîtes pédagogiques
% (définition, propriété, méthode, attention, le savais-tu, exercice résolu,
% exercice, TP, bilan), page de garde et signature OnBuch+ ancrée en bas.
%
% Macros attendues dans main.tex AVANT \input{preamble} :
%   \DOCMATIERE  \DOCNIVEAU  \DOCMODULE  \DOCLECON (numéro)  \DOCTITRE
\documentclass[11pt]{article}

\usepackage{fontspec}
\usepackage[french]{babel}
\usepackage[a4paper,margin=2cm,top=3.4cm,bottom=2.8cm,headheight=1.4cm,footskip=1.3cm]{geometry}
\usepackage{amsmath,amssymb,amsfonts}
\usepackage{mathtools} % \xmapsto, \coloneqq... souvent utilisés par les modèles
\usepackage{cancel} % simplifications barrées
% \abs{z} (module, valeur absolue) et \norm{u} (norme), utilisés par les modèles
\providecommand{\abs}{}\let\abs\relax\DeclarePairedDelimiter{\abs}{\lvert}{\rvert}
\providecommand{\norm}{}\let\norm\relax\DeclarePairedDelimiter{\norm}{\lVert}{\rVert}
\usepackage[table]{xcolor} % [table] : fournit aussi \rowcolors (alternance de lignes)
\usepackage{pagecolor}
\usepackage{tikz}
\usetikzlibrary{arrows.meta,positioning,calc,shapes.geometric,shapes.misc,shapes.arrows,decorations.pathmorphing,decorations.pathreplacing,decorations.markings,patterns,shadows,mindmap,trees,fit,backgrounds,matrix,intersections,angles,quotes}
\usepackage{pgfplots}
\usepackage{pgf-pie} % diagrammes circulaires \pie (statistiques)
\pgfplotsset{compat=1.18}
\usepgfplotslibrary{fillbetween,groupplots} % aires entre courbes (calcul intégral)
\usepackage{enumitem}
\usepackage{fancyhdr}
% [most] charge toutes les bibliothèques tcolorbox (listings, minted, raster,
% documentation...) et gonfle inutilement la mémoire TeX sur un document long
% (~300 boîtes) : on ne charge que ce qui est réellement utilisé.
\usepackage[skins,breakable]{tcolorbox}
\usepackage{graphicx}
\usepackage{colortbl}
\usepackage{booktabs}
\usepackage{tabularx}
\usepackage{diagbox} % case d'en-tête diagonale des tableaux à double entrée
% Colonnes extensibles centrée / gauche / droite (C, L, R), souvent utilisées par les modèles
\newcolumntype{C}{>{\centering\arraybackslash}X}
\newcolumntype{L}{>{\raggedright\arraybackslash}X}
\newcolumntype{R}{>{\raggedleft\arraybackslash}X}
\usepackage{array}
\usepackage{multicol}
\usepackage{siunitx}
% parse-numbers=false : accepte aussi \SI{2,5 \times 10^{-3}}{...}
\sisetup{locale=FR,output-decimal-marker={,},group-minimum-digits=4,parse-numbers=false}
\DeclareSIUnit\ohm{\ensuremath{\symup{\Omega}}} % U+2126 absent de la police
\usepackage{qrcode}
% \degree : commande fréquemment utilisée par les modèles pour un angle ou une
% température en dehors de \SI{}{} (non fournie par les paquets chargés).
\providecommand{\degree}{^{\circ}}
% \qed (fin de démonstration), utilisé par les modèles sans amsthm
\providecommand{\qed}{\hfill$\square$}
% \barre : double barre des valeurs interdites dans un tableau de variations (tkz-tab)
\providecommand{\barre}{\ensuremath{\|}}
% environnement proof (amsthm non chargé) utilisé par les modèles
\ifdefined\proof\else\newenvironment{proof}[1][Démonstration]{\par\noindent\textit{#1.}\ }{\qed\par}\fi
\usepackage{lastpage}
\usepackage{hyperref}
\hypersetup{hidelinks}

% ============================================================
% Palette « Pop » OnBuch+ — mêmes hex que la classe P (plus_ui.dart)
% ============================================================
\definecolor{poporange}{HTML}{F59321}
\definecolor{popdark}{HTML}{E07A0C}
\definecolor{popcream}{HTML}{FFF6E8}
\definecolor{popink}{HTML}{1C1714}
\definecolor{poppurple}{HTML}{7A5AE0}
\definecolor{popgreen}{HTML}{2E9E6A}
\definecolor{poppink}{HTML}{E0457B}
\definecolor{popblue}{HTML}{2F7DE1}
\definecolor{popgold}{HTML}{C98A12}
\definecolor{popmuted}{HTML}{8A7F76}
\definecolor{popline}{HTML}{EADFD2}
\definecolor{popgray}{HTML}{8A7F76}
\colorlet{popgrey}{popgray}
% Alias tolérants (le modèle nomme parfois les couleurs différemment)
\colorlet{popwhite}{white}
\colorlet{popblack}{popink}
\colorlet{popred}{poppink}
\colorlet{popyellow}{popgold}
% Teintes claires pour fonds de tableaux / zones
\colorlet{popblueL}{popblue!10!white}
\colorlet{poporangeL}{poporange!14!white}
\colorlet{popgreenL}{popgreen!12!white}
\colorlet{poppurpleL}{poppurple!12!white}
\colorlet{poppinkL}{poppink!12!white}
\colorlet{popgoldL}{popgold!14!white}

\pagecolor{popcream}

% ============================================================
% Typographie
% ============================================================
\defaultfontfeatures{Ligatures=TeX}
\setmainfont{PlusJakartaSans-Medium.ttf}[Path=../../../fonts/,BoldFont=PlusJakartaSans-Bold.ttf]
% Maths en sans-serif (Fira Math) avec chiffres et lettres droites de Plus
% Jakarta, pour que les formules se fondent dans le texte.
\usepackage[math-style=ISO,bold-style=ISO]{unicode-math}
\setmathfont{FiraMath-Regular.otf}[Path=../../../fonts/]
\setmathfont{PlusJakartaSans-Medium.ttf}[Path=../../../fonts/,range={up/num,up/{Latin,latin},\mathup}]
\usepackage{icomma} % virgule décimale sans espace en mode math
% Caractères mathématiques Unicode absents des polices de texte (Plus Jakarta,
% Archivo Black) : rendus par leur équivalent mathématique, en texte comme en
% formule — sinon ils s'affichent en carré vide (ex. « Arithmétique dans ℤ »).
\usepackage{newunicodechar}
\newunicodechar{ℝ}{\ensuremath{\mathbb{R}}}
\newunicodechar{ℤ}{\ensuremath{\mathbb{Z}}}
\newunicodechar{ℂ}{\ensuremath{\mathbb{C}}}
\newunicodechar{ℕ}{\ensuremath{\mathbb{N}}}
\newunicodechar{ℚ}{\ensuremath{\mathbb{Q}}}
\newunicodechar{𝔻}{\ensuremath{\mathbb{D}}}
\newunicodechar{α}{\ensuremath{\alpha}}
\newunicodechar{β}{\ensuremath{\beta}}
\newunicodechar{γ}{\ensuremath{\gamma}}
\newunicodechar{δ}{\ensuremath{\delta}}
\newunicodechar{ε}{\ensuremath{\varepsilon}}
\newunicodechar{θ}{\ensuremath{\theta}}
\newunicodechar{λ}{\ensuremath{\lambda}}
\newunicodechar{μ}{\ensuremath{\mu}}
\newunicodechar{σ}{\ensuremath{\sigma}}
\newunicodechar{τ}{\ensuremath{\tau}}
\newunicodechar{φ}{\ensuremath{\varphi}}
\newunicodechar{ψ}{\ensuremath{\psi}}
\newunicodechar{ω}{\ensuremath{\omega}}
\newunicodechar{Σ}{\ensuremath{\Sigma}}
% majuscules produites par \MakeUppercase dans les titres (α-aminés -> Α)
\newunicodechar{Α}{\ensuremath{\alpha}}
\newunicodechar{Β}{\ensuremath{\beta}}
\newunicodechar{Γ}{\ensuremath{\Gamma}}
\newunicodechar{Δ}{\ensuremath{\Delta}}
\newunicodechar{Ε}{\ensuremath{\varepsilon}}
\newunicodechar{Θ}{\ensuremath{\Theta}}
\newunicodechar{Λ}{\ensuremath{\Lambda}}
\newunicodechar{Μ}{\ensuremath{\mu}}
\newunicodechar{Τ}{\ensuremath{\tau}}
\newunicodechar{Φ}{\ensuremath{\Phi}}
\newunicodechar{Ψ}{\ensuremath{\Psi}}
\newunicodechar{Ω}{\ensuremath{\Omega}}
\newunicodechar{↦}{\ensuremath{\mapsto}}
\newunicodechar{∈}{\ensuremath{\in}}
\newunicodechar{∉}{\ensuremath{\notin}}
\newunicodechar{∧}{\ensuremath{\wedge}}
\newunicodechar{⇔}{\ensuremath{\Leftrightarrow}}
\newunicodechar{⟺}{\ensuremath{\Longleftrightarrow}}
\newunicodechar{⇒}{\ensuremath{\Rightarrow}}
\newunicodechar{⟹}{\ensuremath{\Longrightarrow}}
\newunicodechar{∘}{\ensuremath{\circ}}
\newunicodechar{∪}{\ensuremath{\cup}}
\newunicodechar{∩}{\ensuremath{\cap}}
\newunicodechar{≡}{\ensuremath{\equiv}}
\newunicodechar{⊂}{\ensuremath{\subset}}
\newunicodechar{⊥}{\ensuremath{\perp}}
\newunicodechar{∃}{\ensuremath{\exists}}
\newunicodechar{∀}{\ensuremath{\forall}}
\newunicodechar{∛}{\ensuremath{\sqrt[3]{\,}}}
\newunicodechar{∜}{\ensuremath{\sqrt[4]{\,}}}
\newunicodechar{⃗}{\ensuremath{{}^{\to}}}
\newunicodechar{ⁿ}{\ifmmode^{n}\else\textsuperscript{\ensuremath{n}}\fi}
\newunicodechar{ˣ}{\ifmmode^{x}\else\textsuperscript{\ensuremath{x}}\fi}
\newunicodechar{ᵇ}{\ifmmode^{b}\else\textsuperscript{\ensuremath{b}}\fi}
\newunicodechar{ʳ}{\ifmmode^{r}\else\textsuperscript{\ensuremath{r}}\fi}
\newunicodechar{ᵘ}{\ifmmode^{u}\else\textsuperscript{\ensuremath{u}}\fi}
\newunicodechar{ʸ}{\ifmmode^{y}\else\textsuperscript{\ensuremath{y}}\fi}
\newunicodechar{ᵃ}{\ifmmode^{a}\else\textsuperscript{\ensuremath{a}}\fi}
\newunicodechar{ᵉ}{\ifmmode^{e}\else\textsuperscript{\ensuremath{e}}\fi}
\newunicodechar{ᵏ}{\ifmmode^{k}\else\textsuperscript{\ensuremath{k}}\fi}
\newunicodechar{ᵖ}{\ifmmode^{p}\else\textsuperscript{\ensuremath{p}}\fi}
\newunicodechar{ᴺ}{\ifmmode^{N}\else\textsuperscript{\ensuremath{N}}\fi}
\newunicodechar{ᶜ}{\ifmmode^{c}\else\textsuperscript{\ensuremath{c}}\fi}
\newunicodechar{ʰ}{\ifmmode^{h}\else\textsuperscript{\ensuremath{h}}\fi}
\newunicodechar{ᵠ}{\ifmmode^{q}\else\textsuperscript{\ensuremath{q}}\fi}
\newunicodechar{ₙ}{\ifmmode_{n}\else\textsubscript{\ensuremath{n}}\fi}
\newunicodechar{ₐ}{\ifmmode_{a}\else\textsubscript{\ensuremath{a}}\fi}
\newunicodechar{ᵢ}{\ifmmode_{i}\else\textsubscript{\ensuremath{i}}\fi}
\newunicodechar{ᵣ}{\ifmmode_{r}\else\textsubscript{\ensuremath{r}}\fi}
\newunicodechar{ₖ}{\ifmmode_{k}\else\textsubscript{\ensuremath{k}}\fi}
\newunicodechar{ᵦ}{\ifmmode_{\beta}\else\textsubscript{\ensuremath{\beta}}\fi}
\newunicodechar{ₓ}{\ifmmode_{x}\else\textsubscript{\ensuremath{x}}\fi}
\newfontfamily\popdisplay{ArchivoBlack-Regular.ttf}[Path=../../../fonts/]
\newfontfamily\poplogo{SpaceGrotesk-Bold.ttf}[Path=../../../fonts/]
\newfontfamily\popsemi{PlusJakartaSans-SemiBold.ttf}[Path=../../../fonts/]
\newfontfamily\popxbold{PlusJakartaSans-ExtraBold.ttf}[Path=../../../fonts/]
\newfontfamily\popxboldls{PlusJakartaSans-ExtraBold.ttf}[Path=../../../fonts/,LetterSpace=3.0]

\setlist[enumerate]{leftmargin=1.7em,itemsep=5pt,topsep=4pt}
\setlist[itemize]{leftmargin=1.7em,itemsep=4pt,topsep=4pt,label={\color{poporange}\textbullet}}
\setlength{\parindent}{0pt}
\setlength{\parskip}{5pt}
\renewcommand{\arraystretch}{1.25}

% pgfplots : style Pop par défaut
\pgfplotsset{
  every axis/.append style={axis line style={popink,line width=1pt},tick style={popink},
    grid style={popline,line width=0.5pt},label style={font=\small},tick label style={font=\footnotesize}},
  popcurve/.style={poporange,line width=1.8pt,no markers,smooth},
}
% Même style disponible dans un \draw TikZ simple (hors pgfplots)
\tikzset{popcurve/.style={poporange,line width=1.8pt}}

% ============================================================
% Pill / squiggle
% ============================================================
\newcommand{\poppill}[3]{%
  \tikz[baseline=-0.55ex]\node[draw=popink,line width=1.2pt,rounded corners=2.6pt,%
    fill=#2,inner xsep=6.5pt,inner ysep=3pt]{%
    \popxboldls\fontsize{7}{7}\selectfont\color{#3}\MakeUppercase{#1}};%
}
\newcommand{\popsquiggle}[1]{%
  \tikz[baseline=-0.4ex]{
    \draw[#1,line width=2.6pt,line cap=round]
      plot [smooth] coordinates {(0,0.09) (0.34,0.01) (0.68,0.21) (1.02,0.01) (1.36,0.10)};
  }%
}
% Mot-clé surligné (marqueur orange sous le texte)
\newcommand{\cle}[1]{{\popxbold\color{popdark}#1}}

% ============================================================
% En-tête / pied
% ============================================================
\pagestyle{fancy}
\fancyhf{}
\renewcommand{\headrulewidth}{0pt}
\renewcommand{\footrulewidth}{0pt}
\lhead{\raisebox{-0.15ex}{\poplogo\fontsize{13}{13}\selectfont\color{popink}OnBuch\color{poporange}+}}
\rhead{\poppill{\DOCMATIERE\ \textbullet\ \DOCNIVEAU\ \textbullet\ Leçon \DOCLECON}{white}{popink}}
\lfoot{\popxbold\fontsize{7.6}{7.6}\selectfont\color{popmuted}\MakeUppercase{Cours signé OnBuch+}\ \textbullet\ {\popsemi\DOCTITRE}}
\rfoot{\tikz[baseline=-0.6ex]\node[rounded corners=8.5pt,fill=popink,minimum height=17pt,minimum width=17pt,inner xsep=5pt,inner sep=0]{\popxbold\fontsize{8}{8}\selectfont\color{popcream}\thepage\,/\,\pageref*{LastPage}};}

% ============================================================
% Titres
% ============================================================
\newcommand{\chaptitre}[2]{%
  \begin{tcolorbox}[enhanced,colback=poporange,colframe=popink,boxrule=2.6pt,arc=11pt,
      drop shadow=popink,left=15pt,right=15pt,top=12pt,bottom=11pt,before skip=3pt,after skip=18pt]
    \poppill{#2}{popink}{popcream}\par\vspace{9pt}
    {\raggedright\hyphenpenalty=10000\exhyphenpenalty=10000\popdisplay\fontsize{22}{24}\selectfont\color{popcream}\MakeUppercase{#1}\par}
  \end{tcolorbox}
}
% Section numérotée : pastille orange avec le numéro + titre Archivo
\newcounter{popsec}
\newcommand{\coursec}[1]{%
  \stepcounter{popsec}\setcounter{popsub}{0}%
  \par\vspace{16pt}\noindent
  \begin{minipage}{\linewidth}
  \tikz[baseline=-0.7ex]\node[draw=popink,line width=1.6pt,fill=poporange,rounded corners=4pt,minimum size=22pt,inner sep=0]{\popdisplay\fontsize{12}{12}\selectfont\color{popcream}\thepopsec};%
  \hspace{8pt}{\popdisplay\fontsize{15}{17}\selectfont\color{popink}\MakeUppercase{#1}}\par
  \vspace{4pt}\noindent{\color{poporange}\rule{36pt}{3.6pt}}
  \end{minipage}\par\nopagebreak\vspace{8pt}
}
\newcounter{popsub}
\newcommand{\courssub}[1]{%
  \stepcounter{popsub}%
  \par\vspace{9pt}\noindent{\popxbold\fontsize{11.5}{13}\selectfont\color{popink}{\color{poporange}\thepopsec.\thepopsub}\hspace{6pt}#1}\par\nopagebreak\vspace{4pt}
}

% ============================================================
% Boîtes pédagogiques — même recette : fond blanc, bordure épaisse colorée,
% ombre dure encre, badge plein. Titre optionnel [#1].
% ============================================================
\tcbset{popbase/.style={breakable,enhanced,colback=white,boxrule=2.3pt,arc=9pt,
  shadow={3pt}{-3pt}{0pt}{popink},left=14pt,right=14pt,top=11pt,bottom=11pt,before skip=11pt,after skip=15pt,
  fontupper=\popsemi\fontsize{10.6}{14.5}\selectfont\color{popink},
  fontlower=\popsemi\fontsize{10.6}{14.5}\selectfont\color{popink}}}
% Coche dessinée (le glyphe ✓ n'existe pas dans les polices du document)
\newcommand{\popcheck}[1]{\tikz[baseline=-0.5ex]\draw[#1,line width=1.6pt,line cap=round,line join=round](-0.13,0.0)--(-0.04,-0.09)--(0.13,0.1);}
\newcommand{\popboxhead}[3]{\poppill{#1}{#2}{white}\ifx\relax#3\relax\else\hspace{6pt}{\popxbold\fontsize{10.6}{12}\selectfont\color{popink}#3}\fi\par\vspace{7pt}}

\newtcolorbox{aretenir}[1][]{popbase,colframe=popgreen,before upper={\popboxhead{À retenir}{popgreen}{#1}}}
\newtcolorbox{exemplebox}[1][]{popbase,colframe=poporange,before upper={\popboxhead{Exemple}{poporange}{#1}}}
\newtcolorbox{definition}[1][]{popbase,colframe=popblue,before upper={\popboxhead{Définition}{popblue}{#1}}}
\newtcolorbox{propriete}[1][]{popbase,colframe=poppurple,before upper={\popboxhead{Propriété}{poppurple}{#1}}}
\newtcolorbox{methode}[1][]{popbase,colframe=popgold,before upper={\popboxhead{Méthode}{popgold}{#1}}}
\newtcolorbox{attention}[1][]{popbase,colframe=poppink,before upper={\popboxhead{Attention}{poppink}{#1}}}
\newtcolorbox{experience}[1][]{popbase,colframe=popblue,colback=popblueL,before upper={\popboxhead{Expérience}{popblue}{#1}}}
% « Le savais-tu ? » : carte violette pleine (contexte, culture, Cameroun)
\newtcolorbox{savaistu}[1][]{popbase,colback=poppurple,colframe=popink,
  fontupper=\popsemi\fontsize{10.4}{14.2}\selectfont\color{white},
  before upper={\poppill{Le savais-tu\,?}{white}{poppurple}\ifx\relax#1\relax\else\hspace{6pt}{\popxbold\color{white}#1}\fi\par\vspace{7pt}}}
% Exercice résolu : énoncé en haut, \tcblower puis la solution sur fond crème
\newcounter{popexo}\newcounter{popexr}
\newtcolorbox{exoresolu}[1][]{popbase,colframe=poporange,colbacklower=popcream,
  segmentation style={popink,line width=1.2pt,dashed},
  before upper={\stepcounter{popexr}\popboxhead{Exercice résolu \thepopexr}{poporange}{#1}},
  before lower={\poppill{Solution}{popink}{popcream}\par\vspace{6pt}}}
% Exercice (non résolu en place) — la correction est dans la partie Corrigés
\newtcolorbox{exercice}[1][]{popbase,colframe=popink,
  before upper={\stepcounter{popexo}\popboxhead{Exercice \thepopexo}{popink}{#1}}}
% Corrigé
\newtcolorbox{corrige}[1][]{popbase,colframe=popgreen,colback=popgreenL,
  before upper={\popboxhead{Corrigé}{popgreen}{#1}}}
% Objectifs / prérequis / bilan
\newtcolorbox{objectifs}{popbase,colframe=popink,colback=poporangeL,
  before upper={\popboxhead{Objectifs de la leçon}{popink}{}}}
\newtcolorbox{prerequis}{popbase,colframe=popink,colback=popblueL,
  before upper={\popboxhead{Prérequis}{popblue}{}}}
\newtcolorbox{bilan}{popbase,colframe=popink,colback=white,boxrule=2.8pt,
  before upper={\popboxhead{Fiche bilan}{poporange}{L'essentiel à connaître pour l'examen}}}
% Figure encadrée avec légende
\newtcolorbox{popfigure}[1][]{enhanced,breakable=false,colback=white,colframe=popline,boxrule=1.4pt,arc=8pt,
  left=10pt,right=10pt,top=10pt,bottom=8pt,before skip=10pt,after skip=12pt,halign=center,
  lower separated=false,halign lower=center,
  fontlower=\popsemi\fontsize{9}{11.5}\selectfont\color{popmuted},
  #1}
\newcommand{\legende}[1]{\par\vspace{6pt}{\popsemi\fontsize{9}{11.5}\selectfont\color{popmuted}#1}}

% ============================================================
% Signature OnBuch+
% ============================================================
\newcommand{\poplogoinline}[1]{{\poplogo\fontsize{#1}{#1}\selectfont\color{popink}OnBuch\color{poporange}+}}
\newcommand{\signatureonbuch}{%
  \begin{tcolorbox}[enhanced,colback=popink,colframe=popink,boxrule=2.2pt,arc=10pt,
      drop shadow=poporange,left=14pt,right=14pt,top=10pt,bottom=10pt,before skip=0pt,after skip=0pt]
    \tikz[baseline=-0.7ex]\node[circle,fill=popgreen,draw=popcream,line width=1.3pt,minimum size=22pt,inner sep=0]{\popcheck{white}};%
    \hspace{9pt}\begin{minipage}[c]{0.62\linewidth}
      {\popxbold\fontsize{12}{13}\selectfont\color{popcream}Signé\ \textbullet\ {\poplogo OnBuch\color{poporange}+}}\par\vspace{2pt}
      {\raggedright\popsemi\fontsize{8}{10.5}\selectfont\color{popline}\MakeUppercase{Équipe pédagogique OnBuch+}\\\MakeUppercase{Conforme au programme officiel MINESEC}\par}
    \end{minipage}\hfill
    {\poplogo\fontsize{22}{22}\selectfont\color{popcream}OnBuch\color{poporange}+}
  \end{tcolorbox}%
}

% ============================================================
% Page de garde
% #1 titre · #2 matière · #3 niveau · #4 module · #5 n° leçon · #6 durée ·
% #7 description / objectifs (texte ou itemize)
% ============================================================
\newcommand{\pagegarde}[7]{%
  \thispagestyle{empty}
  \newgeometry{a4paper,margin=1.7cm,top=1.6cm,bottom=1.4cm}%
  \begin{tikzpicture}[remember picture,overlay]
    \fill[poporange!18] ([xshift=-3.2cm,yshift=-3.6cm]current page.north east) circle (4.6cm);
    \fill[poppurple!14] ([xshift=2.4cm,yshift=7.6cm]current page.south west) circle (3.8cm);
  \end{tikzpicture}%
  {\poplogo\fontsize{30}{30}\selectfont\color{popink}OnBuch\color{poporange}+}\hfill
  \raisebox{0.6ex}{\poppill{Cours signé \textbullet\ Premium}{popink}{popcream}}\par
  \vspace{1pt}\popsquiggle{poppurple}\par
  \vspace{8pt}
  \begin{tcolorbox}[enhanced,colback=poporange,colframe=popink,boxrule=2.8pt,arc=13pt,
      drop shadow=popink,left=18pt,right=18pt,top=12pt,bottom=13pt,after skip=10pt]
    \poppill{#2 \textbullet\ #3}{popink}{popcream}\hspace{5pt}\poppill{#4}{white}{popink}\par\vspace{8pt}
    {\popdisplay\fontsize{14}{14}\selectfont\color{popink}LEÇON #5}\par\vspace{4pt}
    {\raggedright\hyphenpenalty=10000\exhyphenpenalty=10000\popdisplay\fontsize{25}{27}\selectfont\color{popcream}\MakeUppercase{#1}\par}
    \vspace{8pt}
    {\popxbold\fontsize{9.5}{9.5}\selectfont\color{popink}\MakeUppercase{Durée conseillée : #6}}
  \end{tcolorbox}
  \begin{tcolorbox}[enhanced,colback=white,colframe=popink,boxrule=2.2pt,arc=10pt,
      drop shadow=popink,left=16pt,right=16pt,top=10pt,bottom=10pt,after skip=8pt]
    \poppill{Dans cette leçon}{poppurple}{white}\par\vspace{5pt}
    \popsemi\fontsize{9.3}{12.6}\selectfont\color{popink}#7
  \end{tcolorbox}
  \noindent\begin{minipage}[t]{0.487\linewidth}
    \begin{tcolorbox}[enhanced,colback=popgreen,colframe=popink,boxrule=2.2pt,arc=10pt,
        drop shadow=popink,left=11pt,right=11pt,top=10pt,bottom=10pt,height=3.1cm,valign=center]
      \tcbox[colback=white,colframe=popink,boxrule=1.3pt,arc=5pt,boxsep=0pt,nobeforeafter,
        left=3pt,right=3pt,top=3pt,bottom=3pt,valign=center]{\qrcode[height=1.9cm]{https://play.google.com/store/apps/details?id=com.luvvix.onbuch&hl=fr}}%
      \hspace{8pt}\begin{minipage}[c]{0.5\linewidth}
        {\popdisplay\fontsize{11}{12.5}\selectfont\color{white}\MakeUppercase{Télécharge OnBuch}}\par\vspace{4pt}
        {\raggedright\popsemi\fontsize{8.4}{11}\selectfont\color{white}Gratuit \textbullet\ Google Play\\Tuteur IA, annales, concours, résultats\par}
      \end{minipage}
    \end{tcolorbox}
  \end{minipage}\hfill
  \begin{minipage}[t]{0.487\linewidth}
    \begin{tcolorbox}[enhanced,colback=poppurple,colframe=popink,boxrule=2.2pt,arc=10pt,
        drop shadow=popink,left=11pt,right=11pt,top=10pt,bottom=10pt,height=3.1cm,valign=center]
      \tikz[baseline=-0.7ex]\node[circle,fill=white,draw=popink,line width=1.3pt,minimum size=1.6cm,inner sep=0]{\popdisplay\fontsize{24}{24}\selectfont\color{poppurple}+};%
      \hspace{8pt}\begin{minipage}[c]{0.6\linewidth}
        {\popdisplay\fontsize{11}{12.5}\selectfont\color{white}\MakeUppercase{Passe à OnBuch+}}\par\vspace{4pt}
        {\raggedright\popsemi\fontsize{8.4}{11}\selectfont\color{white}Tous les cours signés, Tuteur IA illimité, fiches PDF — onglet OnBuch+ de l'app\par}
      \end{minipage}
    \end{tcolorbox}
  \end{minipage}
  \vspace{8pt}
  \popcontenu
  \vfill
  \signatureonbuch
  \restoregeometry
  \newpage
}

% Rangée « Ce cours contient » (page de garde)
\newcommand{\poptile}[3]{% #1 couleur #2 gros texte #3 légende
  \begin{tcolorbox}[enhanced,colback=#1,colframe=popink,boxrule=2pt,arc=9pt,shadow={2.5pt}{-2.5pt}{0pt}{popink},
      left=6pt,right=6pt,top=9pt,bottom=9pt,halign=center,valign=center,height=1.75cm,width=\linewidth,nobeforeafter]
    {\popdisplay\fontsize{12.5}{13}\selectfont\color{white}\MakeUppercase{#2}}\par\vspace{3pt}
    {\centering\popsemi\fontsize{7.8}{9.5}\selectfont\color{white}#3\par}
  \end{tcolorbox}}
\newcommand{\popcontenu}{%
  \par\noindent{\popxboldls\fontsize{7.5}{7.5}\selectfont\color{popmuted}CE COURS SIGNÉ CONTIENT}\par\vspace{7pt}
  \noindent\begin{minipage}[t]{0.235\linewidth}\poptile{popblue}{Cours}{complet, illustré, pas à pas}\end{minipage}\hfill
  \begin{minipage}[t]{0.235\linewidth}\poptile{poporange}{Méthodes}{et exercices résolus}\end{minipage}\hfill
  \begin{minipage}[t]{0.235\linewidth}\poptile{poppink}{Exercices}{type examen + corrigés}\end{minipage}\hfill
  \begin{minipage}[t]{0.235\linewidth}\poptile{popgreen}{Fiche bilan}{l'essentiel à retenir}\end{minipage}\par}

% Sommaire
\newtcolorbox{sommaire}{enhanced,colback=white,colframe=popink,boxrule=2.2pt,arc=10pt,
  drop shadow=popink,left=18pt,right=18pt,top=13pt,bottom=13pt,before skip=6pt,after skip=20pt,
  before upper={\poppill{Sommaire}{poppurple}{white}\par\vspace{9pt}\popsemi\fontsize{10.6}{14.5}\selectfont\color{popink}}}

% Signature de fin de cours — poussée tout en bas de la dernière page
\newcommand{\findecours}{%
  \par\vfill\nopagebreak
  \begin{center}\popsquiggle{poporange}\end{center}\vspace{4pt}
  \signatureonbuch
}
\AtBeginDocument{\def\llbracket{\mathopen{[\mkern-3.5mu[}}\def\rrbracket{\mathclose{]\mkern-3.5mu]}}} % intervalles d'entiers (glyphe ⟦ absent de la police)
% Glyphes absents de Fira Math : redéfinis à partir de glyphes disponibles
\AtBeginDocument{%
  \def\setminus{\mathbin{\backslash}}\let\smallsetminus\setminus
  \def\vdots{\mathord{\vbox{\baselineskip3.2pt\lineskiplimit0pt\kern4pt\hbox{.}\hbox{.}\hbox{.}}}}%
  \def\ddots{\mathinner{\mkern1mu\raise6.5pt\hbox{.}\mkern2mu\raise3.6pt\hbox{.}\mkern2mu\raise0.7pt\hbox{.}\mkern1mu}}%
  \def\triangle{\mathord{\scalebox{0.9}{$\symup{\Delta}$}}}%
  \def\nabla{\mathord{\raisebox{\depth}{\scalebox{1}[-1]{$\symup{\Delta}$}}}}%
  \def\checkmark{\popcheck{popdark}}%
  \def\star{\mathbin{\ast}}\def\bigstar{\mathord{\ast}}%
  \def\diamond{\mathbin{\scalebox{0.6}{$\Diamond$}}}%
  \def\bigcup{\mathop{\mathchoice{\raisebox{-0.3em}{\scalebox{1.7}{$\cup$}}}{\scalebox{1.25}{$\cup$}}{\cup}{\cup}}\displaylimits}%
  \def\bigcap{\mathop{\mathchoice{\raisebox{-0.3em}{\scalebox{1.7}{$\cap$}}}{\scalebox{1.25}{$\cap$}}{\cap}{\cap}}\displaylimits}%
  \def\lesseqqgtr{\mathrel{\lessgtr}}\def\bigoplus{\mathop{\oplus}\displaylimits}%
}
\providecommand{\ssi}{si et seulement si} % abréviation fréquente
\AtBeginDocument{\providecommand{\wideparen}{\overparen}} % arc de cercle

%% FIGFIX-BEGIN v4 — correctifs globaux des figures (docs/onbuch-generation/tools/figfix)
\usepackage{adjustbox}\usepackage{etoolbox}
% toute figure TikZ/pgfplots plus large que la ligne est réduite à la largeur disponible
\BeforeBeginEnvironment{tikzpicture}{\begin{adjustbox}{max width=\linewidth}}
\AfterEndEnvironment{tikzpicture}{\end{adjustbox}}
% légendes pgfplots lisibles par-dessus les courbes
\pgfplotsset{every axis/.append style={legend style={fill=white,fill opacity=0.92,text opacity=1,draw=black!15,inner sep=3pt}}}
% étiquettes d'arêtes (syntaxe edge["..."]) avec fond blanc
\tikzset{every edge quotes/.append style={fill=white,inner sep=1.2pt}}
%% FIGFIX-END
\begin{document}

\pagegarde{Statistiques : séries regroupées en classes}{Mathématiques}{1ère D}{Organisation et gestion des données}{N09}{5 h}{Cette leçon prolonge l'étude des séries statistiques au cas où les données sont regroupées en classes d'amplitude donnée. On y définit et calcule les paramètres de position (classe modale, mode, moyenne, médiane par interpolation linéaire) et de dispersion (écart moyen, variance, écart-type), ainsi que les représentations graphiques adaptées : histogramme et courbes cumulatives.
\vspace{4pt}
\begin{multicols}{2}\small
\begin{itemize}
\item À la fin de cette leçon, je sais organiser une série statistique en classes et calculer les effectifs et fréquences cumulés croissants ou décroissants.
\item À la fin de cette leçon, je sais calculer la moyenne d'une série regroupée en classes à l'aide des centres de classes.
\item À la fin de cette leçon, je sais déterminer la classe modale et le mode d'une série regroupée en classes.
\item À la fin de cette leçon, je sais déterminer l'intervalle médian et calculer la médiane par interpolation linéaire, par le calcul et graphiquement.
\item À la fin de cette leçon, je sais calculer l'écart moyen, la variance et l'écart-type d'une série regroupée en classes.
\item À la fin de cette leçon, je sais construire et interpréter un histogramme (y compris à classes d'amplitudes inégales).
\end{itemize}\end{multicols}}

\begin{sommaire}
\begin{multicols}{2}
\begin{enumerate}
\item Regroupement en classes, effectifs et fréquences cumulés
\item Paramètres de position : classe modale, mode, moyenne
\item Médiane : intervalle médian et interpolation linéaire
\item Paramètres de dispersion : écart moyen, variance, écart-type
\item Représentations graphiques : histogramme et courbes cumulatives
\item Synthèse et interprétation en situation concrète
\item Activité d'intégration
\item Méthodes et exercices résolus
\item Exercices
\item Corrigés des exercices
\item Fiche bilan
\end{enumerate}
\end{multicols}
\end{sommaire}

\begin{objectifs}
\begin{itemize}
\item À la fin de cette leçon, je sais organiser une série statistique en classes et calculer les effectifs et fréquences cumulés croissants ou décroissants.
\item À la fin de cette leçon, je sais calculer la moyenne d'une série regroupée en classes à l'aide des centres de classes.
\item À la fin de cette leçon, je sais déterminer la classe modale et le mode d'une série regroupée en classes.
\item À la fin de cette leçon, je sais déterminer l'intervalle médian et calculer la médiane par interpolation linéaire, par le calcul et graphiquement.
\item À la fin de cette leçon, je sais calculer l'écart moyen, la variance et l'écart-type d'une série regroupée en classes.
\item À la fin de cette leçon, je sais construire et interpréter un histogramme (y compris à classes d'amplitudes inégales).
\item À la fin de cette leçon, je sais construire la courbe des effectifs ou fréquences cumulés et l'utiliser pour lire la médiane.
\item À la fin de cette leçon, je sais interpréter les paramètres de position et de dispersion dans une situation concrète de consommation de biens et services.
\end{itemize}
\end{objectifs}

\begin{prerequis}
\begin{itemize}
\item Séries statistiques discrètes : effectifs, fréquences, moyenne, médiane, mode (classe de Seconde)
\item Effectifs et fréquences cumulés pour une série discrète
\item Proportionnalité et règle de trois (utile pour l'interpolation linéaire)
\item Lecture et construction de diagrammes (bâtons, circulaires) et repérage dans un repère orthogonal
\item Calculs sur les fractions, les puissances et la racine carrée
\end{itemize}
\end{prerequis}

\newpage

\coursec{Regroupement en classes, effectifs et fréquences cumulés}

\textbf{Situation d'accroche.} C'est la fin du mois à Bonabéri, quartier populaire de Douala. La SONEL vient de relever les factures d'électricité de $200$ abonnés. Devant le directeur commercial s'étale une liste interminable : $37$~kWh, $82$~kWh, $145$~kWh, $61$~kWh, \dots{} Deux cents nombres, tous différents ou presque ! Impossible d'y voir clair. Un stagiaire propose alors de \emph{regrouper} les abonnés par tranches de consommation : ceux qui consomment entre $0$ et $50$~kWh, ceux entre $50$ et $100$~kWh, et ainsi de suite. En quelques minutes, la liste informe devient un tableau de cinq lignes, lisible d'un coup d'œil.

Mais le directeur veut plus : la consommation \emph{moyenne}, la consommation \emph{la plus fréquente}, et celle qui \emph{partage les abonnés en deux groupes de même taille}, afin de fixer des tranches tarifaires équitables. Problème : une fois le regroupement effectué, on ne connaît plus la valeur exacte de chaque facture, seulement la tranche à laquelle elle appartient. Comment exploiter ces données regroupées en classes ? C'est tout l'objet de cette leçon : résumer, représenter et interpréter une série statistique regroupée en classes. Dans cette première section, nous apprenons à construire rigoureusement le tableau de base : classes, effectifs, fréquences et cumuls.

\courssub{Rappels : population, individu, caractère, série brute}

Avant de regrouper, rappelons le vocabulaire fondamental rencontré en classe de Seconde.

\begin{definition}[Population, individu, caractère]
\begin{itemize}
\item La \cle{population} est l'ensemble sur lequel porte l'étude statistique.
\item Un \cle{individu} (ou unité statistique) est un élément de la population.
\item Le \cle{caractère} (ou variable statistique) est la propriété que l'on observe sur chaque individu. Il est \cle{quantitatif continu} lorsqu'il peut prendre, théoriquement, n'importe quelle valeur d'un intervalle de $\mathbb{R}$ (une durée, une masse, une consommation, une taille).
\item La liste des valeurs relevées, dans l'ordre où elles ont été recueillies, s'appelle la \cle{série brute}. Le nombre total d'individus est l'\cle{effectif total}, noté $N$.
\end{itemize}
\end{definition}

Dans notre situation : la population est l'ensemble des $200$ abonnés du quartier ; un individu est un abonné ; le caractère est la consommation mensuelle en kWh, qui est un caractère quantitatif continu (une consommation peut valoir $37$~kWh, $37{,}2$~kWh, $37{,}25$~kWh\dots) ; l'effectif total est $N = 200$.

\courssub{Pourquoi regrouper en classes ?}

Imagine la série brute des $200$ factures. Deux difficultés apparaissent immédiatement :
\begin{itemize}
\item le caractère est \emph{continu} : presque toutes les valeurs sont distinctes, donc un tableau « valeur -- effectif » aurait près de $200$ lignes, chacune d'effectif $1$ ou $2$ ; il ne résumerait rien du tout ;
\item même pour un caractère discret prenant un très grand nombre de valeurs (par exemple les notes sur $80$ points possibles), le tableau valeur par valeur serait illisible.
\end{itemize}

L'idée est alors de découper l'intervalle des valeurs en quelques \emph{tranches}, appelées \cle{classes}, et de compter combien d'individus tombent dans chaque tranche. On perd une information (la valeur exacte de chaque donnée), mais on gagne une vision globale : la \emph{forme} de la répartition.

\begin{definition}[Classe, amplitude, centre]
Regrouper une série en \cle{classes}, c'est partager l'intervalle des valeurs du caractère en intervalles $[a_i\,;\,b_i[$, appelés classes, deux à deux disjoints et dont la réunion contient toutes les données.
\begin{itemize}
\item L'\cle{amplitude} de la classe $[a_i\,;\,b_i[$ est le nombre $e_i = b_i - a_i$.
\item Le \cle{centre} de la classe est le nombre
\[ c_i = \frac{a_i + b_i}{2}. \]
Le centre $c_i$ servira de \emph{représentant} de la classe dans tous les calculs (moyenne, variance\dots), puisqu'on a perdu les valeurs exactes.
\end{itemize}
\end{definition}

\begin{attention}[Convention d'appartenance : le piège de la borne commune]
Un abonné qui a consommé \emph{exactement} $50$~kWh appartient-il à la classe $[0\,;\,50[$ ou à la classe $[50\,;\,100[$ ? Si l'on n'y prend pas garde, on le compte \textbf{deux fois} (ou pas du tout) et l'effectif total est faussé. La convention universellement adoptée, et exigée au Probatoire, est :
\begin{center}
\textbf{borne inférieure incluse, borne supérieure exclue :} $[a\,;\,b[$.
\end{center}
Ainsi $50$~kWh appartient à $[50\,;\,100[$ et non à $[0\,;\,50[$. Vérifie toujours, à la fin, que la somme des effectifs redonne $N$ : c'est un test de cohérence indispensable.
\end{attention}

\begin{popfigure}
\begin{tikzpicture}[scale=1.05]
\draw[->,thick,popdark] (-0.3,0) -- (10.8,0) node[below left] {\footnotesize consommation (kWh)};
\foreach \x/\lab in {0/0, 2/50, 4/100, 6/150, 8/200, 10/250}{
  \draw[thick,popdark] (\x,-0.12) -- (\x,0.12);
  \node[below] at (\x,-0.15) {\footnotesize $\lab$};
}
\foreach \xa/\xb/\col in {0/2/popblue, 2/4/poporange, 4/6/popgreen, 6/8/poppurple, 8/10/poppink}{
  \draw[\col, very thick] (\xa,0.35) -- (\xb,0.35);
  \fill[\col] (\xa,0.35) circle (0.09);
  \draw[\col, fill=white] (\xb,0.35) circle (0.09);
}
\node[above,popblue] at (1,0.42) {\footnotesize $[0\,;\,50[$};
\node[above,poporange] at (3,0.42) {\footnotesize $[50\,;\,100[$};
\node[above,popgreen] at (5,0.42) {\footnotesize $[100\,;\,150[$};
\node[above,poppurple] at (7,0.42) {\footnotesize $[150\,;\,200[$};
\node[above,poppink] at (9,0.42) {\footnotesize $[200\,;\,250[$};
\node[align=left, text width=9.5cm] at (5,-1.1) {\footnotesize Point plein : borne \textbf{incluse}. Point creux : borne \textbf{exclue}. La valeur $50$ appartient à la classe $[50\,;\,100[$, pas à $[0\,;\,50[$.};
\end{tikzpicture}
\legende{Découpage de l'intervalle des consommations en cinq classes d'amplitude constante $50$~kWh, avec la convention $[a\,;\,b[$.}
\end{popfigure}

\textbf{Comment choisir les classes ?} Il n'existe pas de règle absolue, mais de bonnes pratiques :
\begin{itemize}
\item prendre en général entre $5$ et $10$ classes : trop peu de classes écrasent l'information, trop de classes la fragmentent ;
\item choisir, quand c'est possible, des classes de \emph{même amplitude} : la comparaison des effectifs et la construction de l'histogramme en sont simplifiées ;
\item choisir des bornes « rondes » ($0$, $50$, $100$\dots) pour faciliter la lecture.
\end{itemize}

\begin{savaistu}[La convention du BUCREP]
La convention « borne inférieure incluse, borne supérieure exclue » n'est pas qu'une habitude de manuels : c'est celle adoptée par le Bureau Central des Recensements et des Études de Population (BUCREP) lors des recensements généraux de la population du Cameroun. Lorsqu'on répartit les Camerounais par tranches d'âge $[0\,;\,5[$, $[5\,;\,10[$, \dots, un enfant ayant exactement $5$ ans révolus est compté dans la tranche $[5\,;\,10[$. Cette convention garantit que chaque personne est comptée \emph{une et une seule fois} : condition indispensable pour qu'un recensement soit fiable !
\end{savaistu}

\courssub{Effectifs et fréquences}

\begin{definition}[Effectif, fréquence]
On considère une série de $N$ individus regroupés en $k$ classes.
\begin{itemize}
\item L'\cle{effectif} $n_i$ de la classe numéro $i$ est le nombre d'individus dont la valeur appartient à cette classe. On a toujours
\[ n_1 + n_2 + \dots + n_k = N. \]
\item La \cle{fréquence} de la classe numéro $i$ est le quotient
\[ f_i = \frac{n_i}{N}. \]
C'est un nombre compris entre $0$ et $1$, et $f_1 + f_2 + \dots + f_k = 1$. Multipliée par $100$, elle donne la \cle{fréquence en pourcentage}.
\end{itemize}
\end{definition}

La fréquence répond à la question : « quelle \emph{part} de la population se trouve dans cette classe ? » Dire que la classe $[50\,;\,100[$ a une fréquence de $0{,}30$ signifie que $30\,\%$ des abonnés consomment entre $50$ et $100$~kWh.

\courssub{Effectifs et fréquences cumulés}

Voici maintenant l'outil central de cette section. Le directeur de la SONEL ne demande pas seulement « combien d'abonnés consomment entre $50$ et $100$~kWh ? », mais aussi « combien d'abonnés consomment \emph{moins de} $100$~kWh ? » ou « combien consomment \emph{au moins} $150$~kWh ? ». Pour répondre instantanément à ce type de question, on \emph{cumule} les effectifs.

\begin{definition}[Effectifs et fréquences cumulés]
On ordonne les classes dans l'ordre croissant des valeurs.
\begin{itemize}
\item L'\cle{effectif cumulé croissant} (ECC) d'une classe est la somme des effectifs de cette classe et de toutes les classes précédentes. Il répond à la question : « combien d'individus ont une valeur \textbf{strictement inférieure} à la borne supérieure $b_i$ de la classe ? »
\item L'\cle{effectif cumulé décroissant} (ECD) d'une classe est la somme des effectifs de cette classe et de toutes les classes suivantes. Il répond à la question : « combien d'individus ont une valeur \textbf{supérieure ou égale} à la borne inférieure $a_i$ de la classe ? »
\item Les \cle{fréquences cumulées croissantes} (FCC) et \cle{décroissantes} (FCD) se définissent de la même manière à partir des fréquences :
\[ \text{FCC}_i = f_1 + f_2 + \dots + f_i = \frac{\text{ECC}_i}{N}, \qquad \text{FCD}_i = f_i + f_{i+1} + \dots + f_k = \frac{\text{ECD}_i}{N}. \]
\end{itemize}
\end{definition}

\begin{aretenir}[Lecture des cumuls]
\begin{itemize}
\item $\text{ECC}$ de la classe $[a_i\,;\,b_i[$ = nombre d'individus dont la valeur est \textbf{inférieure à} $b_i$.
\item $\text{ECD}$ de la classe $[a_i\,;\,b_i[$ = nombre d'individus dont la valeur est \textbf{au moins égale à} $a_i$.
\item Le dernier ECC vaut toujours $N$ ; le premier ECD vaut toujours $N$.
\item Pour tout $i$ : $\text{ECD}_i = N - \text{ECC}_{i-1}$ (avec $\text{ECC}_0 = 0$).
\end{itemize}
\end{aretenir}

\begin{methode}[Construire pas à pas le tableau statistique complet]
\begin{enumerate}
\item \textbf{Repérer} les valeurs extrêmes de la série brute et choisir des classes $[a_i\,;\,b_i[$ (en général $5$ à $10$, d'amplitude constante si possible).
\item \textbf{Dépouiller} : parcourir la série brute une seule fois et compter, pour chaque classe, le nombre $n_i$ d'individus, en appliquant la convention « borne inférieure incluse ».
\item \textbf{Contrôler} que $n_1 + \dots + n_k = N$.
\item \textbf{Calculer} les centres $c_i = \dfrac{a_i+b_i}{2}$ et les fréquences $f_i = \dfrac{n_i}{N}$.
\item \textbf{Cumuler en montant} : le premier ECC est $n_1$ ; chaque ECC suivant s'obtient en \emph{ajoutant} l'effectif de la ligne à l'ECC précédent. Le dernier doit valoir $N$.
\item \textbf{Cumuler en descendant} : le dernier ECD est $n_k$ ; chaque ECD précédent s'obtient en ajoutant l'effectif de la ligne à l'ECD suivant. Le premier doit valoir $N$.
\item \textbf{Calculer} les fréquences cumulées : $\text{FCC}_i = \text{ECC}_i / N$ et $\text{FCD}_i = \text{ECD}_i / N$ (ou en cumulant les fréquences directement).
\end{enumerate}
\end{methode}

\begin{exemplebox}[Les 200 abonnés SONEL : tableau complet]
Le dépouillement des $200$ factures du quartier de Bonabéri donne les effectifs suivants. Calculons toutes les colonnes, avec $N = 200$.

Pour la classe $[50\,;\,100[$ par exemple :
\[ c_2 = \frac{50+100}{2} = 75 \text{ kWh}, \qquad f_2 = \frac{60}{200} = 0{,}30 \quad (30\,\%). \]

Effectifs cumulés croissants (on additionne en descendant le tableau) :
\[ 18, \quad 18+60 = 78, \quad 78+64 = 142, \quad 142+40 = 182, \quad 182+18 = 200. \checkmark \]

Effectifs cumulés décroissants (on additionne en remontant le tableau) :
\[ 18, \quad 18+40 = 58, \quad 58+64 = 122, \quad 122+60 = 182, \quad 182+18 = 200. \checkmark \]

Fréquences cumulées croissantes : $\frac{18}{200} = 0{,}09$ ; $\frac{78}{200} = 0{,}39$ ; $\frac{142}{200} = 0{,}71$ ; $\frac{182}{200} = 0{,}91$ ; $\frac{200}{200} = 1$.

\begin{center}
\small
\begin{tabularx}{\linewidth}{|c|>{\centering\arraybackslash}X|>{\centering\arraybackslash}X|>{\centering\arraybackslash}X|>{\centering\arraybackslash}X|>{\centering\arraybackslash}X|>{\centering\arraybackslash}X|>{\centering\arraybackslash}X|}
\hline
\rowcolor{popblueL} Classe & Centre $c_i$ & $n_i$ & $f_i$ & ECC & ECD & FCC & FCD \\
\hline
$[0\,;\,50[$ & $25$ & $18$ & $0{,}09$ & $18$ & $200$ & $0{,}09$ & $1$ \\
\hline
$[50\,;\,100[$ & $75$ & $60$ & $0{,}30$ & $78$ & $182$ & $0{,}39$ & $0{,}91$ \\
\hline
$[100\,;\,150[$ & $125$ & $64$ & $0{,}32$ & $142$ & $122$ & $0{,}71$ & $0{,}61$ \\
\hline
$[150\,;\,200[$ & $175$ & $40$ & $0{,}20$ & $182$ & $58$ & $0{,}91$ & $0{,}29$ \\
\hline
$[200\,;\,250[$ & $225$ & $18$ & $0{,}09$ & $200$ & $18$ & $1$ & $0{,}09$ \\
\hline
\rowcolor{popcream} Total & & $200$ & $1$ & & & & \\
\hline
\end{tabularx}
\end{center}

\textbf{Lectures immédiates :} $142$ abonnés (soit $71\,\%$) consomment \emph{moins de} $150$~kWh ; $58$ abonnés (soit $29\,\%$) consomment \emph{au moins} $150$~kWh. Remarque au passage : $142 + 58 = 200$ : les deux questions sont complémentaires, car la borne $150$ est la frontière.
\end{exemplebox}

\begin{exoresolu}[Dépouillement d'une série brute]
Un vendeur de crédit téléphonique MTN à Bafoussam a noté le montant (en centaines de FCFA) de ses $20$ dernières ventes :
\[ 12 \;;\ 35 \;;\ 8 \;;\ 47 \;;\ 22 \;;\ 15 \;;\ 38 \;;\ 5 \;;\ 28 \;;\ 41 \;;\ 18 \;;\ 32 \;;\ 9 \;;\ 25 \;;\ 44 \;;\ 11 \;;\ 36 \;;\ 21 \;;\ 48 \;;\ 30. \]
Regrouper ces données en classes d'amplitude $10$ à partir de $[0\,;\,10[$, puis donner les effectifs, fréquences, ECC et ECD.
\tcblower
\textbf{Solution.} Les valeurs s'étendent de $5$ à $48$ ; on prend les classes $[0\,;\,10[$, $[10\,;\,20[$, $[20\,;\,30[$, $[30\,;\,40[$, $[40\,;\,50[$.

Dépouillement (attention : $30$ appartient à $[30\,;\,40[$, $10$ n'apparaît pas, $20$ n'apparaît pas, $40$ n'apparaît pas) :
\begin{itemize}
\item $[0\,;\,10[$ : $8\,;\,5\,;\,9$ donc $n_1 = 3$ ;
\item $[10\,;\,20[$ : $12\,;\,15\,;\,18\,;\,11$ donc $n_2 = 4$ ;
\item $[20\,;\,30[$ : $22\,;\,28\,;\,25\,;\,21$ donc $n_3 = 4$ ;
\item $[30\,;\,40[$ : $35\,;\,38\,;\,32\,;\,36\,;\,30$ donc $n_4 = 5$ ;
\item $[40\,;\,50[$ : $47\,;\,41\,;\,44\,;\,48$ donc $n_5 = 4$.
\end{itemize}
Contrôle : $3+4+4+5+4 = 20 = N$. \checkmark

\begin{center}
\small
\begin{tabularx}{\linewidth}{|c|>{\centering\arraybackslash}X|>{\centering\arraybackslash}X|>{\centering\arraybackslash}X|>{\centering\arraybackslash}X|}
\hline
\rowcolor{popgreenL} Classe & $n_i$ & $f_i$ & ECC & ECD \\
\hline
$[0\,;\,10[$ & $3$ & $0{,}15$ & $3$ & $20$ \\
\hline
$[10\,;\,20[$ & $4$ & $0{,}20$ & $7$ & $17$ \\
\hline
$[20\,;\,30[$ & $4$ & $0{,}20$ & $11$ & $13$ \\
\hline
$[30\,;\,40[$ & $5$ & $0{,}25$ & $16$ & $9$ \\
\hline
$[40\,;\,50[$ & $4$ & $0{,}20$ & $20$ & $4$ \\
\hline
\rowcolor{popcream} Total & $20$ & $1$ & & \\
\hline
\end{tabularx}
\end{center}

\textbf{Lecture :} $11$ ventes (soit $55\,\%$) sont d'un montant inférieur à $3\,000$~FCFA ; $9$ ventes sont d'un montant d'au moins $3\,000$~FCFA.
\end{exoresolu}

\begin{attention}[Erreurs fréquentes à éviter]
\begin{itemize}
\item \textbf{Oublier la convention} $[a\,;\,b[$ et compter la borne commune dans deux classes : l'effectif total dépasse alors $N$. Toujours vérifier $\sum n_i = N$.
\item \textbf{Confondre ECC et ECD} : l'ECC se lit « moins de la borne supérieure », l'ECD se lit « au moins la borne inférieure ». Une astuce : l'ECC \emph{augmente} quand on descend le tableau, l'ECD \emph{diminue}.
\item \textbf{Se tromper de sens de cumul} : les ECC se calculent de haut en bas, les ECD de bas en haut. Contrôles : dernier ECC $= N$, premier ECD $= N$, dernière FCC $= 1$.
\item \textbf{Arrondir les fréquences trop tôt} : si tu arrondis chaque $f_i$ avant de cumuler, la dernière FCC peut ne pas valoir exactement $1$. Cumule les valeurs exactes, arrondis seulement à l'affichage.
\end{itemize}
\end{attention}

\begin{exercice}[À toi de jouer]
Les tailles (en cm) de $40$ élèves d'une classe de Première C à Yaoundé sont regroupées ainsi : $[150\,;\,160[$ : $6$ élèves ; $[160\,;\,170[$ : $14$ élèves ; $[170\,;\,180[$ : $15$ élèves ; $[180\,;\,190[$ : $5$ élèves.
\begin{enumerate}
\item Calculer les fréquences, les ECC, les ECD et les FCC.
\item Combien d'élèves mesurent moins de $170$~cm ? Au moins $170$~cm ?
\item Un élève mesure exactement $170$~cm. Dans quelle classe est-il compté ?
\end{enumerate}
\end{exercice}

\begin{corrige}[Exercice 1]
1. $N = 6+14+15+5 = 40$. Fréquences : $\frac{6}{40} = 0{,}15$ ; $\frac{14}{40} = 0{,}35$ ; $\frac{15}{40} = 0{,}375$ ; $\frac{5}{40} = 0{,}125$. ECC : $6$ ; $6+14 = 20$ ; $20+15 = 35$ ; $35+5 = 40$. ECD : $40$ ; $14+15+5 = 34$ ; $15+5 = 20$ ; $5$. FCC : $0{,}15$ ; $0{,}50$ ; $0{,}875$ ; $1$.

2. Moins de $170$~cm : c'est l'ECC de la classe $[160\,;\,170[$, soit $20$ élèves (la moitié de la classe). Au moins $170$~cm : c'est l'ECD de la classe $[170\,;\,180[$, soit $20$ élèves.

3. Par la convention « borne inférieure incluse », un élève de $170$~cm est compté dans la classe $[170\,;\,180[$.
\end{corrige}

\begin{aretenir}[L'essentiel de la section]
\begin{itemize}
\item On regroupe en \cle{classes} $[a_i\,;\,b_i[$ (borne inférieure incluse, borne supérieure exclue) lorsque le caractère est continu ou prend trop de valeurs distinctes.
\item Chaque classe est résumée par son \cle{centre} $c_i = \dfrac{a_i+b_i}{2}$ et caractérisée par son \cle{effectif} $n_i$ et sa \cle{fréquence} $f_i = \dfrac{n_i}{N}$.
\item Les \cle{effectifs cumulés croissants} répondent à « combien valent moins de $b_i$ ? » ; les \cle{effectifs cumulés décroissants} à « combien valent au moins $a_i$ ? ». Même chose pour les fréquences cumulées, en proportion.
\item Contrôles systématiques : $\sum n_i = N$, dernier ECC $= N$, premier ECD $= N$, dernière FCC $= 1$.
\end{itemize}
Ce tableau est désormais notre outil de base : c'est à partir de lui que nous calculerons, dans les sections suivantes, la moyenne, le mode, la médiane et les paramètres de dispersion de la série.
\end{aretenir}

\coursec{Paramètres de position : classe modale, mode, moyenne}

Dans la section précédente, nous avons appris à \cle{regrouper} une série statistique en \cle{classes} et à calculer les \cle{effectifs} et \cle{fréquences cumulés}. Le tableau obtenu est compact et lisible, mais il reste muet sur une question naturelle : \emph{où se situe le « centre » de la série ?} Autrement dit, quelle valeur résume le mieux l'ensemble des données ? C'est toute l'ambition des \cle{paramètres de position} : la \cle{classe modale}, le \cle{mode} et la \cle{moyenne}. Nous les découvrirons sur une situation concrète qui nous accompagnera dans toute la leçon.

\begin{exemplebox}[Situation de référence : les consommations ENEO]
La société ENEO relève les consommations mensuelles d'électricité (en kWh) de $100$ abonnés d'un quartier de Douala. Les données, regroupées en classes d'amplitude $40$, donnent :
\begin{center}
\begin{tabularx}{\linewidth}{|l|X|X|X|X|X|X|}
\hline
\rowcolor{popblueL} Classe (kWh) & $[0;40[$ & $[40;80[$ & $[80;120[$ & $[120;160[$ & $[160;200[$ & $[200;240[$ \\
\hline
Effectif $n_i$ & $6$ & $18$ & $30$ & $23$ & $15$ & $8$ \\
\hline
\end{tabularx}
\end{center}
L'effectif total est $N = 6+18+30+23+15+8 = 100$.
\end{exemplebox}

\courssub{Classe modale et mode}

\textbf{Intuition.}\par
Imagine que tu observes l'histogramme des consommations : la barre la plus haute attire immédiatement le regard. Elle correspond à la tranche de consommation la plus fréquente, celle où se concentrent le plus d'abonnés. C'est exactement l'idée de la classe modale.

\begin{definition}[Classe modale et mode]
Soit une série statistique regroupée en classes de \textbf{même amplitude}.
\begin{itemize}
\item La \cle{classe modale} est la classe dont l'\textbf{effectif est le plus grand} (ou, de façon équivalente, la fréquence la plus grande).
\item Le \cle{mode} de la série est le \textbf{centre de la classe modale} : si la classe modale est $[a;b[$, alors
\[ \text{Mode} = \frac{a+b}{2}. \]
\end{itemize}
Si deux classes ont le même effectif maximal, la série est dite \emph{bimodale} : elle possède deux classes modales.
\end{definition}

\begin{attention}[Classes d'amplitudes inégales]
La définition « plus grand effectif » n'est valable que si toutes les classes ont la \textbf{même amplitude}. Sinon, une classe très large peut mécaniquement contenir beaucoup d'individus sans être réellement « dense ». On utilise alors la \emph{densité d'effectif} (voir le « Le savais-tu » en fin de section).
\end{attention}

\begin{exemplebox}[Application à l'exemple ENEO]
Les classes ont toutes la même amplitude ($40$ kWh). Le plus grand effectif est $30$, atteint pour la classe $[80;120[$.
\begin{itemize}
\item \textbf{Classe modale} : $[80;120[$ (en kWh) ;
\item \textbf{Mode} : $\dfrac{80+120}{2} = 100$ kWh.
\end{itemize}
\emph{Interprétation} : la consommation mensuelle la plus « typique » des abonnés du quartier se situe autour de $100$ kWh ; c'est dans la tranche $[80;120[$ que l'on trouve le plus de foyers.
\end{exemplebox}

\courssub{Moyenne d'une série regroupée en classes}

\textbf{Le problème posé par le regroupement.}\par
Pour une série non regroupée de valeurs $x_1, x_2, \ldots, x_N$, la moyenne est $\bar{x} = \dfrac{1}{N}\sum x_i$. Mais après regroupement en classes, \textbf{nous avons perdu les valeurs exactes} : nous savons seulement que $30$ abonnés consomment entre $80$ et $120$ kWh, sans connaître la consommation précise de chacun. Comment calculer une moyenne ?

\textbf{L'idée fondatrice : l'approximation par les centres.}\par
On fait l'\textbf{hypothèse que les valeurs de chaque classe se répartissent de manière à peu près uniforme} dans la classe : les valeurs en dessous du centre compensent celles au-dessus. Il est alors naturel de \emph{remplacer chaque valeur de la classe par le centre de la classe}. C'est une approximation, mais c'est la meilleure possible avec l'information disponible.

\begin{popfigure}
\begin{center}
\begin{tikzpicture}[scale=1]
% Axe
\draw[->, thick, popdark] (0,0) -- (10.5,0);
% Bornes de la classe
\draw[thick, popblue] (2,0.15) -- (2,-0.15) node[below, popdark] {$80$};
\draw[thick, popblue] (8,0.15) -- (8,-0.15) node[below, popdark] {$120$};
% Classe colorée
\fill[popblueL, opacity=0.6] (2,0.05) rectangle (8,0.55);
\node[popdark] at (5,0.85) {\small classe $[80;120[$ : $30$ valeurs inconnues};
% Points représentant les valeurs inconnues
\foreach \x in {2.4,2.9,3.3,3.8,4.2,4.6,5.1,5.4,5.8,6.2,6.6,7.1,7.5}
  \fill[popmuted] (\x,0.3) circle (1.6pt);
% Centre
\draw[very thick, poporange] (5,0.15) -- (5,-0.15);
\fill[poporange] (5,-0.55) circle (2.2pt);
\node[poporange, below] at (5,-0.7) {\small centre $c = \frac{80+120}{2} = 100$};
% Flèches de remplacement
\draw[->, poporange, dashed] (3.3,0.55) .. controls (3.8,1.3) .. (4.8,1.15);
\draw[->, poporange, dashed] (6.6,0.55) .. controls (6.2,1.3) .. (5.2,1.15);
\node[poporange, align=center] at (5,1.75) {\small on remplace chaque valeur par le centre};
\end{tikzpicture}
\legende{Approximation des valeurs d'une classe par son centre : les écarts en dessous et au-dessus du centre se compensent.}
\end{center}
\end{popfigure}

\begin{propriete}[Moyenne d'une série regroupée en classes]
Soit une série regroupée en $k$ classes de centres $c_1, c_2, \ldots, c_k$, d'effectifs $n_1, n_2, \ldots, n_k$ et d'effectif total $N = n_1 + n_2 + \cdots + n_k$. La \cle{moyenne} de la série est
\[ \bar{x} = \frac{1}{N}\sum_{i=1}^{k} n_i\, c_i = \frac{n_1 c_1 + n_2 c_2 + \cdots + n_k c_k}{N}. \]
De façon équivalente, si $f_i = \dfrac{n_i}{N}$ désigne la \textbf{fréquence} de la classe $i$, alors
\[ \bar{x} = \sum_{i=1}^{k} f_i\, c_i = f_1 c_1 + f_2 c_2 + \cdots + f_k c_k. \]
La moyenne ainsi calculée est une \textbf{valeur approchée} de la moyenne exacte (inconnue) de la série brute ; elle s'exprime dans la même unité que le caractère étudié.
\end{propriete}

\begin{experience}[Pourquoi cette formule est-elle naturelle ?]
Raisonnons sur la classe $[80;120[$ d'effectif $30$. Chacune des $30$ valeurs inconnues est remplacée par le centre $c_3 = 100$. La \emph{somme approchée} des consommations de ces $30$ abonnés est donc $30 \times 100 = 3000$ kWh. En faisant de même pour chaque classe, la somme totale approchée des $N$ consommations est $\sum n_i c_i$, et en divisant par $N$ on obtient la consommation moyenne par abonné : $\bar{x} = \dfrac{1}{N}\sum n_i c_i$. La formule avec les fréquences s'en déduit immédiatement : $\dfrac{n_i c_i}{N} = f_i c_i$, donc $\bar{x} = \sum f_i c_i$.
\end{experience}

\begin{methode}[Organiser le calcul de la moyenne en tableau]
Pour calculer proprement la moyenne d'une série regroupée en classes :
\begin{enumerate}
\item Ajouter une ligne (ou colonne) des \textbf{centres} $c_i = \dfrac{\text{borne inf} + \text{borne sup}}{2}$ ;
\item Ajouter une ligne des \textbf{produits} $n_i c_i$ ;
\item Calculer la \textbf{somme} des effectifs $N$ et la somme des produits $\sum n_i c_i$ ;
\item Conclure : $\bar{x} = \dfrac{\sum n_i c_i}{N}$.
\end{enumerate}
Ce tableau sera réutilisé tel quel pour la variance (section 4) : prends l'habitude de le construire soigneusement.
\end{methode}

\begin{exoresolu}[Moyenne des consommations ENEO]
Énoncé. Calculer la moyenne des consommations mensuelles des $100$ abonnés de l'exemple ENEO, puis vérifier le résultat avec les fréquences.
\tcblower
Solution détaillée. On construit le tableau de calcul :
\begin{center}
\begin{tabularx}{\linewidth}{|l|X|X|X|X|X|X|X|}
\hline
\rowcolor{popblueL} Classe & $[0;40[$ & $[40;80[$ & $[80;120[$ & $[120;160[$ & $[160;200[$ & $[200;240[$ & Total \\
\hline
Centre $c_i$ & $20$ & $60$ & $100$ & $140$ & $180$ & $220$ & --- \\
\hline
Effectif $n_i$ & $6$ & $18$ & $30$ & $23$ & $15$ & $8$ & $N = 100$ \\
\hline
$n_i c_i$ & $120$ & $1080$ & $3000$ & $3220$ & $2700$ & $1760$ & $11880$ \\
\hline
\end{tabularx}
\end{center}
D'où :
\[ \bar{x} = \frac{11880}{100} = 118{,}8 \text{ kWh}. \]
Vérification par les fréquences : $f_1 = 0{,}06$ ; $f_2 = 0{,}18$ ; $f_3 = 0{,}30$ ; $f_4 = 0{,}23$ ; $f_5 = 0{,}15$ ; $f_6 = 0{,}08$. Alors
\begin{align*}
\bar{x} &= 0{,}06 \times 20 + 0{,}18 \times 60 + 0{,}30 \times 100 + 0{,}23 \times 140 + 0{,}15 \times 180 + 0{,}08 \times 220 \\
&= 1{,}2 + 10{,}8 + 30 + 32{,}2 + 27 + 17{,}6 = 118{,}8 \text{ kWh}.
\end{align*}
Les deux méthodes donnent le même résultat, ce qui conforte le calcul.
\emph{Interprétation} : en moyenne, un abonné du quartier consomme environ $118{,}8$ kWh par mois, soit approximativement $119$ kWh.
\end{exoresolu}

\begin{attention}[Erreur fréquente : bornes ou amplitudes à la place des centres]
Dans la formule $\bar{x} = \dfrac{1}{N}\sum n_i c_i$, le facteur $c_i$ est \textbf{toujours le centre de la classe}. Trois confusions reviennent chaque année au Probatoire :
\begin{itemize}
\item utiliser la \textbf{borne inférieure} (par exemple $80$ au lieu de $100$) : la moyenne est alors systématiquement sous-estimée ;
\item utiliser l'\textbf{amplitude} ($40$) à la place du centre : erreur grave, qui mélange deux notions différentes ;
\item oublier de \textbf{pondérer par les effectifs} : calculer $\dfrac{20+60+100+140+180+220}{6}$ revient à donner le même poids à une classe de $6$ abonnés et à une classe de $30$ !
\end{itemize}
Réflexe de survie : \emph{centre, produit, somme, division par} $N$. Autre contrôle rapide : la moyenne doit toujours appartenir à l'intervalle global des valeurs (ici $[0;240[$) ; un résultat de $720$ kWh ou de $19{,}8$ kWh signalerait immédiatement une erreur.
\end{attention}

\courssub{Linéarité de la moyenne}

\begin{propriete}[Linéarité de la moyenne (admise)]
Si l'on transforme toutes les valeurs d'une série par une fonction affine $x \mapsto ax + b$ (avec $a$ et $b$ réels), alors la moyenne subit la même transformation :
\[ \overline{ax + b} = a\,\bar{x} + b. \]
Cette propriété est \textbf{admise} ; elle se démontre facilement à partir de la définition et nous la vérifions ci-dessous sur un exemple.
\end{propriete}

\begin{exemplebox}[Vérification sur un exemple : conversion en francs CFA]
Supposons que le kWh soit facturé $a = 50$ FCFA avec un abonnement fixe $b = 1000$ FCFA. La facture d'un abonné est $y = 50x + 1000$ (en FCFA). La propriété annonce :
\[ \bar{y} = 50 \times 118{,}8 + 1000 = 5940 + 1000 = 6940 \text{ FCFA}. \]
Vérifions directement : les centres transformés sont $50c_i + 1000$, soit $2000$ ; $4000$ ; $6000$ ; $8000$ ; $10000$ ; $12000$. Alors
\[ \bar{y} = \frac{6 \times 2000 + 18 \times 4000 + 30 \times 6000 + 23 \times 8000 + 15 \times 10000 + 8 \times 12000}{100} = \frac{694000}{100} = 6940. \]
Les deux calculs coïncident : la propriété est confirmée. Elle est très utile pour changer d'unité (kWh vers FCFA, degrés Celsius vers Fahrenheit, etc.) sans refaire tout le tableau.
\end{exemplebox}

\courssub{Comparaison mode / moyenne et interprétation}

Sur l'exemple ENEO, nous avons obtenu deux « centres » différents :
\begin{itemize}
\item \textbf{Mode} : $100$ kWh (centre de la classe la plus fréquente) ;
\item \textbf{Moyenne} : $118{,}8$ kWh (consommation moyenne par abonné).
\end{itemize}
Pourquoi la moyenne est-elle \emph{supérieure} au mode ? Parce que la distribution est \textbf{étalée vers les grandes valeurs} : quelques abonnés consomment beaucoup (classes $[160;200[$ et $[200;240[$), et ces fortes consommations « tirent » la moyenne vers le haut, alors qu'elles ne déplacent pas la classe modale. C'est un phénomène général : la moyenne est \cle{sensible aux valeurs extrêmes}, le mode ne l'est pas.

\emph{En situation concrète} : pour dimensionner son réseau, ENEO raisonne sur la \textbf{moyenne} (elle donne la consommation totale du quartier : $100 \times 118{,}8 = 11880$ kWh à fournir). Pour décrire le « consommateur typique » dans une campagne publicitaire, on préférera le \textbf{mode} ($100$ kWh), moins influencé par les gros consommateurs.

\begin{aretenir}[L'essentiel de la section]
\begin{itemize}
\item \textbf{Classe modale} : classe de plus grand effectif (amplitudes égales) ; \textbf{mode} : centre de cette classe.
\item \textbf{Moyenne} : $\bar{x} = \dfrac{1}{N}\sum n_i c_i = \sum f_i c_i$, où $c_i$ est le \textbf{centre} de la classe $i$.
\item La moyenne est sensible aux valeurs extrêmes ; le mode ne l'est pas.
\item Linéarité : $\overline{ax+b} = a\bar{x} + b$.
\end{itemize}
\end{aretenir}

\begin{savaistu}[Et si les amplitudes sont inégales ?]
Lorsque les classes n'ont pas la même amplitude, comparer les effectifs bruts n'a plus de sens : une classe deux fois plus large recueille « mécaniquement » plus d'individus. On définit alors la \cle{densité d'effectif} de chaque classe par $d_i = \dfrac{n_i}{\text{amplitude}_i}$, et la \textbf{classe modale est la classe de plus grande densité}. C'est aussi la densité qui commande la \emph{hauteur} des barres d'un histogramme à classes inégales (section 5). Retiens cette idée : elle deviendra centrale en Terminale, où l'on parlera de densité de probabilité.
\end{savaistu}

\begin{exercice}[À toi de jouer]
Un opérateur téléphonique relève la durée mensuelle d'appels (en minutes) de $80$ abonnés à Yaoundé :
\begin{center}
\begin{tabularx}{\linewidth}{|l|X|X|X|X|X|}
\hline
\rowcolor{popblueL} Durée (min) & $[0;100[$ & $[100;200[$ & $[200;300[$ & $[300;400[$ & $[400;500[$ \\
\hline
Effectif & $10$ & $24$ & $26$ & $14$ & $6$ \\
\hline
\end{tabularx}
\end{center}
\begin{enumerate}
\item Déterminer la classe modale et le mode de cette série.
\item Calculer la durée moyenne d'appels à l'aide d'un tableau de calcul.
\item Comparer mode et moyenne, puis interpréter.
\end{enumerate}
\end{exercice}

\begin{corrige}[Exercice : durées d'appels]
\begin{enumerate}
\item Le plus grand effectif est $26$ : la classe modale est $[200;300[$ et le mode est $\dfrac{200+300}{2} = 250$ min.
\item Centres : $50$ ; $150$ ; $250$ ; $350$ ; $450$. Produits $n_i c_i$ : $500$ ; $3600$ ; $6500$ ; $4900$ ; $2700$. Somme : $18200$, avec $N = 80$. Donc $\bar{x} = \dfrac{18200}{80} = 227{,}5$ min.
\item Ici la moyenne ($227{,}5$) est \emph{inférieure} au mode ($250$) : la distribution est étalée vers les petites valeurs (classe $[0;100[$ non négligeable). L'abonné « typique » appelle autour de $250$ min, mais les petits consommateurs tirent la moyenne vers le bas.
\end{enumerate}
\end{corrige}

\coursec{Médiane : intervalle médian et interpolation linéaire}

Dans la section précédente, nous avons appris à résumer une série regroupée en classes par la \cle{classe modale}, le \cle{mode} et la \cle{moyenne}. Ces paramètres sont précieux, mais la moyenne possède un défaut bien connu : elle est très sensible aux valeurs extrêmes. Si, dans un quartier de Yaoundé, un seul abonné ENEO consomme énormément d'électricité (une usine, par exemple), la moyenne des consommations du quartier grimpe, alors que la consommation de la plupart des familles n'a pas changé. Il nous faut donc un autre paramètre de position, plus « robuste » : la \cle{médiane}, qui partage la population en deux groupes de même effectif. Toute cette section est consacrée à sa définition, à son calcul par \cle{interpolation linéaire} et à sa lecture graphique.

\courssub{Rappel : qu'est-ce que la médiane ?}

\begin{definition}[Médiane d'une série statistique]
La \cle{médiane} d'une série statistique, notée $Me$, est la valeur du caractère qui partage la série ordonnée en deux sous-séries de même effectif :
\begin{itemize}
\item au moins $50\,\%$ des individus ont une valeur inférieure ou égale à $Me$ ;
\item au moins $50\,\%$ des individus ont une valeur supérieure ou égale à $Me$.
\end{itemize}
\end{definition}

Pour une série discrète non regroupée, on ordonne les valeurs et on prend la valeur centrale (ou la demi-somme des deux valeurs centrales). Mais lorsque les données sont \cle{regroupées en classes}, on ne connaît plus les valeurs individuelles : on sait seulement combien d'individus tombent dans chaque intervalle. La médiane ne peut donc plus être lue directement ; il faut d'abord localiser la \emph{classe} qui la contient, puis estimer sa position à l'intérieur de cette classe. C'est toute la stratégie de cette section.

\begin{aretenir}[Idée directrice]
La médiane est la valeur $Me$ telle que l'effectif cumulé croissant atteigne exactement $\dfrac{N}{2}$, où $N$ est l'effectif total. Tout le travail consiste à résoudre, de manière exacte ou approchée, l'équation $\mathrm{ECC}(Me) = \dfrac{N}{2}$.
\end{aretenir}

\courssub{L'intervalle médian (classe médiane)}

Reprenons notre fil conducteur : l'ENEO a relevé les consommations mensuelles (en kWh) de $200$ abonnés d'un quartier de Douala et les a regroupées en classes.

\begin{center}
\begin{tabularx}{\linewidth}{|l|X|X|X|X|X|}
\hline
\rowcolor{popblueL} Classe (kWh) & $[0\,;50[$ & $[50\,;100[$ & $[100\,;150[$ & $[150\,;200[$ & $[200\,;250[$ \\
\hline
Effectif $n_i$ & $30$ & $50$ & $70$ & $35$ & $15$ \\
\hline
ECC & $30$ & $80$ & $150$ & $185$ & $200$ \\
\hline
\end{tabularx}
\end{center}

L'effectif total est $N = 200$, donc $\dfrac{N}{2} = 100$. La médiane est la consommation atteinte par le $100^{\text{e}}$ abonné lorsqu'on les classe par consommation croissante. Regardons les effectifs cumulés croissants : après la classe $[50\,;100[$, on a compté $80$ abonnés ; après la classe $[100\,;150[$, on en a compté $150$. Le $100^{\text{e}}$ abonné se trouve donc \emph{quelque part} dans la classe $[100\,;150[$.

\begin{definition}[Intervalle médian]
On appelle \cle{intervalle médian} (ou \cle{classe médiane}) d'une série regroupée en classes la \textbf{première classe} dont l'effectif cumulé croissant atteint ou dépasse $\dfrac{N}{2}$.
\end{definition}

Dans notre exemple, $80 < 100 \leq 150$ : l'intervalle médian est donc $[100\,;150[$. On sait alors que $100 \leq Me < 150$, mais cette simple localisation est insuffisante : on veut une \emph{valeur} précise de $Me$.

\begin{attention}[Classe médiane $\neq$ classe modale]
Ne confonds pas la \cle{classe médiane} et la \cle{classe modale} ! La classe modale est la classe de \textbf{plus grand effectif} ; la classe médiane est la classe où l'\textbf{ECC atteint} $\dfrac{N}{2}$. Ces deux classes peuvent être différentes. Par exemple, avec des effectifs $90$, $20$, $60$, $20$, $10$ (donc $N = 200$ et $\dfrac{N}{2} = 100$) :
\begin{itemize}
\item la classe modale est $[0\,;50[$, car c'est elle qui a le plus grand effectif ($90$) ;
\item les ECC sont $90$ ; $110$ ; $170$ ; $190$ ; $200$, et le premier ECC qui atteint $100$ est $110$ : la classe médiane est donc $[50\,;100[$.
\end{itemize}
Retiens : \textbf{les deux notions se déterminent par des critères différents} — effectif maximal pour l'une, cumul atteignant $\dfrac{N}{2}$ pour l'autre.
\end{attention}

\courssub{Interpolation linéaire : principe et formule}

\textbf{Le problème.} Nous savons que $Me \in [100\,;150[$, mais les données regroupées ne nous disent pas où exactement. Pour avancer, on fait une hypothèse naturelle et raisonnable.

\begin{definition}[Hypothèse de répartition uniforme]
À l'intérieur de la classe médiane, on suppose que les valeurs sont \cle{réparties uniformément} : les $70$ abonnés de la classe $[100\,;150[$ se répartissent de façon régulière entre $100$ et $150$ kWh. Autrement dit, l'effectif cumulé croît \emph{proportionnellement} à la valeur du caractère à l'intérieur de la classe.
\end{definition}

Sous cette hypothèse, la courbe des ECC, qui est une ligne brisée, est un \emph{segment de droite} sur la classe médiane. Trouver $Me$ revient alors à un simple problème de proportionnalité — une règle de trois — que le théorème de Thalès justifie géométriquement.

\begin{propriete}[Formule de la médiane par interpolation linéaire]
Soit une série regroupée en classes, d'effectif total $N$. Notons :
\begin{itemize}
\item $[a\,;b[$ la classe médiane, d'amplitude $b - a$ et d'effectif $n$ ;
\item $E$ l'effectif cumulé croissant de la classe \textbf{précédant} la classe médiane (avec $E = 0$ si la classe médiane est la première classe).
\end{itemize}
Alors, sous l'hypothèse de répartition uniforme dans la classe médiane :
\[ Me = a + \left(\dfrac{N}{2} - E\right) \times \dfrac{b - a}{n}. \]
Cette formule est exigible au Probatoire ; sa justification par proportionnalité fait l'objet de la démonstration guidée ci-dessous.
\end{propriete}

\begin{experience}[Démonstration guidée de la formule]
Nous allons établir la formule pas à pas, par la proportionnalité (règle de trois).

\textbf{Étape 1 — Ce qu'il manque.} Avant d'entrer dans la classe médiane $[a\,;b[$, on a déjà cumulé $E$ individus. Pour atteindre la moitié de l'effectif total, il nous en manque encore :
\[ \dfrac{N}{2} - E \quad \text{individus.} \]

\textbf{Étape 2 — Proportion dans la classe.} La classe médiane contient $n$ individus, supposés uniformément répartis sur l'intervalle $[a\,;b[$. La fraction de la classe qu'il faut « parcourir » pour rencontrer les $\dfrac{N}{2} - E$ individus manquants est donc :
\[ \dfrac{\dfrac{N}{2} - E}{n}. \]

\textbf{Étape 3 — Conversion en unités du caractère.} La classe entière a une amplitude $b - a$. Parcourir la fraction ci-dessus de la classe revient à avancer, à partir de $a$, d'une longueur :
\[ \dfrac{\dfrac{N}{2} - E}{n} \times (b - a). \]

\textbf{Étape 4 — Conclusion.} La médiane est la valeur de départ $a$ augmentée de cette avancée :
\[ Me = a + \left(\dfrac{N}{2} - E\right) \times \dfrac{b - a}{n}. \qquad \blacksquare \]

\textbf{Lecture géométrique (Thalès).} Sur le segment de la courbe des ECC joignant les points $(a\,;E)$ et $(b\,;E+n)$, le point d'ordonnée $\dfrac{N}{2}$ a pour abscisse $Me$. Les triangles rectangles formés par ce segment étant en situation de Thalès, on a $\dfrac{Me - a}{b - a} = \dfrac{\frac{N}{2} - E}{n}$, ce qui redonne exactement la formule.
\end{experience}

\begin{popfigure}
\begin{center}
\begin{tikzpicture}[scale=1.05,>=Stealth]
% Axes
\draw[->,popmuted] (0,0) -- (7.4,0) node[below left,popdark]{\small $x$ (kWh)};
\draw[->,popmuted] (0,0) -- (0,5.2) node[below left,popdark]{\small ECC};
% Bornes de la classe médiane
\draw[popdark] (1.5,0.08) -- (1.5,-0.08) node[below=4pt]{\small $a=100$};
\draw[popdark] (6,0.08) -- (6,-0.08) node[below=4pt]{\small $b=150$};
% Niveaux ECC
\draw[popdark] (0.08,1.6) -- (-0.08,1.6) node[left=4pt]{\small $E=80$};
\draw[popdark] (0.08,4.4) -- (-0.08,4.4) node[left=4pt]{\small $E+n=150$};
\draw[popdark] (0.08,3) -- (-0.08,3) node[left=4pt]{\small $\frac{N}{2}=100$};
% Segment de la courbe des ECC sur la classe médiane
\draw[very thick,popblue] (1.5,1.6) -- (6,4.4);
\fill[popblue] (1.5,1.6) circle (2.2pt);
\fill[popblue] (6,4.4) circle (2.2pt);
% Point médiane (abscisse exacte : 3.75)
\fill[poporange] (3.75,3) circle (2.6pt);
% Pointillés de lecture
\draw[dashed,poporange] (3.75,3) -- (3.75,0) node[below=4pt,poporange]{\small $Me$};
\draw[dashed,poporange] (0,3) -- (3.75,3);
% Triangles de Thalès
\draw[popgreen,thick] (1.5,1.6) -- (6,1.6) -- (6,4.4);
\draw[popgreen!70!black,thick,dashed] (1.5,1.6) -- (3.75,1.6) -- (3.75,3);
% Annotations
\node[popgreen!60!black,font=\small] at (3.9,1.25) {$b-a$};
\node[popgreen!60!black,font=\small,rotate=90] at (6.3,3) {$n$};
\node[poporange,font=\small,align=center] at (2.35,2.5) {$\frac{N}{2}-E$};
\node[popblue,font=\small,align=center] at (4.9,4.75) {courbe des ECC\\ (segment sur la classe médiane)};
\end{tikzpicture}
\end{center}
\legende{Interpolation linéaire dans la classe médiane : par Thalès, $\dfrac{Me-a}{b-a}=\dfrac{\frac{N}{2}-E}{n}$.}
\end{popfigure}

\courssub{Méthode et exemple résolu}

\begin{methode}[Déterminer la médiane en 4 étapes]
\begin{enumerate}
\item \textbf{Calculer} $\dfrac{N}{2}$ et dresser la colonne des effectifs cumulés croissants.
\item \textbf{Identifier} la classe médiane : première classe dont l'ECC atteint ou dépasse $\dfrac{N}{2}$. Noter ses bornes $a$ et $b$, son effectif $n$, et l'ECC $E$ de la classe précédente.
\item \textbf{Appliquer} la formule $Me = a + \left(\dfrac{N}{2} - E\right) \times \dfrac{b-a}{n}$.
\item \textbf{Conclure} par une phrase d'interprétation concrète avec l'unité.
\end{enumerate}
\end{methode}

\begin{exoresolu}[Médiane des consommations ENEO]
\textbf{Énoncé.} La répartition des consommations mensuelles de $200$ abonnés est donnée par le tableau ci-dessus. Déterminer la médiane de cette série par interpolation linéaire, puis interpréter le résultat.

\tcblower
\textbf{Solution détaillée.}

\textbf{Étape 1.} $N = 200$, donc $\dfrac{N}{2} = 100$. Les ECC sont : $30$ ; $80$ ; $150$ ; $185$ ; $200$.

\textbf{Étape 2.} Le premier ECC qui atteint ou dépasse $100$ est $150$ : la classe médiane est $[100\,;150[$. On relève : $a = 100$, $b = 150$, $n = 70$, et l'ECC de la classe précédente $E = 80$.

\textbf{Étape 3.} On applique la formule :
\begin{align*}
Me &= a + \left(\dfrac{N}{2} - E\right) \times \dfrac{b-a}{n} \\
&= 100 + (100 - 80) \times \dfrac{150 - 100}{70} \\
&= 100 + 20 \times \dfrac{50}{70} \\
&= 100 + \dfrac{1000}{70} \\
&\approx 100 + 14{,}29 \\
&\approx 114{,}29 \text{ kWh}.
\end{align*}

\textbf{Étape 4.} \textbf{Conclusion.} La médiane vaut environ $114{,}3$ kWh : \emph{la moitié des abonnés du quartier consomme moins de $114{,}3$ kWh par mois, et l'autre moitié consomme davantage.}

\textbf{Vérification.} On a bien $100 \leq 114{,}29 < 150$ : la valeur obtenue appartient à la classe médiane, ce qui est obligatoire. Sur la courbe des ECC ci-dessous, le point d'ordonnée $100$ a bien une abscisse comprise entre $100$ et $150$, voisine de $114$ : le calcul est cohérent avec la lecture graphique.
\end{exoresolu}

\begin{attention}[Les deux erreurs classiques]
\begin{itemize}
\item \textbf{Oublier de retrancher $E$.} L'erreur la plus fréquente est d'écrire $Me = a + \dfrac{N}{2} \times \dfrac{b-a}{n}$ au lieu de $Me = a + \left(\dfrac{N}{2} - E\right) \times \dfrac{b-a}{n}$. Ici, cela donnerait $100 + 100 \times \dfrac{50}{70} \approx 171{,}4$, une valeur \emph{hors de la classe médiane} : un tel résultat doit immédiatement te mettre la puce à l'oreille. Contrôle systématique : on doit toujours avoir $a \leq Me < b$.
\item \textbf{Se tromper de classe.} La classe médiane se repère sur les \textbf{ECC}, pas sur les effectifs $n_i$. Recalcule toujours la colonne des cumuls avant de conclure.
\end{itemize}
\end{attention}

\courssub{Détermination graphique de la médiane}

L'interpolation linéaire revient exactement à lire la médiane sur la \cle{courbe des effectifs cumulés croissants} : celle-ci est formée de segments joignant les points $(b_i\,;\mathrm{ECC}_i)$, où $b_i$ est la borne supérieure de chaque classe. La médiane est l'\textbf{abscisse du point de cette courbe dont l'ordonnée est} $\dfrac{N}{2}$.

\begin{popfigure}
\begin{center}
\begin{tikzpicture}
\begin{axis}[
width=\linewidth, height=7.2cm,
xmin=0, xmax=260, ymin=0, ymax=215,
xlabel={Consommation (kWh)}, ylabel={Effectif cumulé croissant},
xtick={0,50,100,150,200,250},
ytick={0,50,100,150,200},
grid=both, grid style={popline!40},
axis line style={popmuted},
label style={popdark},
]
% Courbe des ECC (ligne brisée)
\addplot[very thick, popblue, mark=*, mark size=1.8pt] coordinates {
(0,0) (50,30) (100,80) (150,150) (200,185) (250,200)
};
% Lecture graphique de Me
\addplot[dashed, poporange, thick] coordinates {(0,100) (114.29,100)};
\addplot[dashed, poporange, thick] coordinates {(114.29,0) (114.29,100)};
\addplot[only marks, mark=*, mark size=2.5pt, poporange] coordinates {(114.29,100)};
\node[poporange, anchor=south west, font=\small] at (axis cs:118,102) {$Me \approx 114{,}3$};
\node[popdark, anchor=south east, font=\small] at (axis cs:2,100) {$\frac{N}{2}=100$};
\end{axis}
\end{tikzpicture}
\end{center}
\legende{Courbe des effectifs cumulés croissants : la médiane est l'abscisse du point d'ordonnée $\frac{N}{2} = 100$.}
\end{popfigure}

\begin{methode}[Lecture graphique de la médiane]
\begin{enumerate}
\item Construire la courbe des ECC en plaçant les points $(b_i\,;\mathrm{ECC}_i)$ et en les reliant par des segments.
\item Tracer la droite horizontale d'ordonnée $\dfrac{N}{2}$.
\item Projeter le point d'intersection avec la courbe sur l'axe des abscisses : on lit $Me$.
\end{enumerate}
On peut aussi tracer sur le même graphique la courbe des effectifs cumulés \textbf{décroissants} : les deux courbes se coupent en un point dont l'abscisse est également la médiane (car à cet endroit, autant d'individus sont en dessous qu'au-dessus).
\end{methode}

\begin{exercice}[Consommations d'eau à Bafoussam]
La CAMWATER a relevé les consommations mensuelles (en m$^3$) de $120$ foyers : $[0\,;5[$ : $18$ foyers ; $[5\,;10[$ : $30$ ; $[10\,;15[$ : $42$ ; $[15\,;20[$ : $20$ ; $[20\,;25[$ : $10$.
\begin{enumerate}
\item Dresser le tableau des ECC et déterminer la classe médiane.
\item Calculer la médiane par interpolation linéaire.
\item Interpréter le résultat par une phrase.
\end{enumerate}
\end{exercice}

\begin{corrige}[Exercice : consommations d'eau à Bafoussam]
\begin{enumerate}
\item $N = 120$, donc $\dfrac{N}{2} = 60$. ECC : $18$ ; $48$ ; $90$ ; $110$ ; $120$. Le premier ECC atteignant $60$ est $90$ : classe médiane $[10\,;15[$, avec $a = 10$, $b = 15$, $n = 42$, $E = 48$.
\item $Me = 10 + (60 - 48) \times \dfrac{5}{42} = 10 + \dfrac{60}{42} \approx 11{,}43$ m$^3$. On vérifie : $10 \leq 11{,}43 < 15$.
\item La moitié des foyers de Bafoussam étudiés consomme moins d'environ $11{,}4$ m$^3$ d'eau par mois, et l'autre moitié consomme davantage.
\end{enumerate}
\end{corrige}

\begin{savaistu}[Pourquoi parle-t-on de « revenu médian » ?]
La médiane est \textbf{insensible aux valeurs extrêmes}, contrairement à la moyenne. Si un milliardaire s'installe dans un village, le revenu moyen du village explose, alors que le revenu médian bouge à peine : la situation de l'habitant « du milieu » n'a pas changé. C'est pourquoi les études sur la pauvreté — y compris celles de l'INS Cameroun — privilégient le \emph{revenu médian} ou le \emph{salaire médian} : il décrit le niveau de vie de la moitié de la population, sans être déformé par quelques fortunes exceptionnelles. Retiens : la \textbf{moyenne} résume les quantités totales, la \textbf{médiane} résume la position de l'individu central.
\end{savaistu}

\begin{aretenir}[L'essentiel de la section]
\begin{itemize}
\item La \cle{classe médiane} est la première classe dont l'ECC atteint ou dépasse $\dfrac{N}{2}$.
\item Sous l'hypothèse de répartition uniforme, la médiane se calcule par \cle{interpolation linéaire} :
\[ Me = a + \left(\dfrac{N}{2} - E\right) \times \dfrac{b - a}{n}, \]
où $[a\,;b[$ est la classe médiane, $n$ son effectif et $E$ l'ECC de la classe précédente.
\item Graphiquement, $Me$ est l'abscisse du point de la courbe des ECC d'ordonnée $\dfrac{N}{2}$ (ou l'abscisse de l'intersection des courbes ECC et ECD).
\item Vérification obligatoire : on doit obtenir $a \leq Me < b$.
\end{itemize}
\end{aretenir}

\coursec{Paramètres de dispersion : écart moyen, variance, écart-type}

Dans la section précédente, nous avons appris à déterminer la \cle{médiane} d'une série regroupée en classes, d'abord en repérant l'intervalle médian, puis en affinant le résultat par interpolation linéaire. Moyenne et médiane nous renseignent sur la \emph{position} des données : elles nous disent « autour de quelle valeur » se situe la série. Mais elles ne disent rien de la manière dont les valeurs se \emph{répartissent} autour de ce centre. C'est toute l'ambition de cette nouvelle section : construire des outils qui mesurent la \cle{dispersion} d'une série statistique.

\courssub{Pourquoi mesurer la dispersion ?}

\begin{experience}[Deux quartiers, une même moyenne]
Imaginons deux quartiers de Douala, Akwa et Bonabéri, où l'on relève les consommations mensuelles d'électricité (en kWh) de cinq ménages chacun :
\begin{center}
\begin{tabular}{|l|c|c|c|c|c|c|}
\hline
\rowcolor{popblueL} Quartier & Ménage 1 & Ménage 2 & Ménage 3 & Ménage 4 & Ménage 5 & Moyenne \\
\hline
Akwa & 160 & 165 & 168 & 170 & 172 & 167 \\
\hline
Bonabéri & 60 & 90 & 170 & 240 & 275 & 167 \\
\hline
\end{tabular}
\end{center}
Vérifions le calcul de la moyenne d'Akwa : $\dfrac{160+165+168+170+172}{5} = \dfrac{835}{5} = 167$~kWh ; de même pour Bonabéri : $\dfrac{60+90+170+240+275}{5} = \dfrac{835}{5} = 167$~kWh. Les deux moyennes sont identiques. Pourtant, à Akwa tous les ménages consomment presque pareil, tandis qu'à Bonabéri les écarts entre ménages sont énormes. La moyenne seule ne permet donc pas de distinguer ces deux réalités.
\end{experience}

Cette petite expérience de pensée montre qu'il faut un deuxième type de paramètre : un nombre qui mesure à quel point les valeurs \emph{s'éloignent} de la moyenne. Plus ce nombre est grand, plus la série est \emph{dispersée} ; plus il est petit, plus la série est \emph{homogène}.

\begin{popfigure}
\begin{center}
\begin{tikzpicture}[scale=0.9]
% Histogramme resserré
\begin{scope}
\draw[->, popdark] (0,0) -- (5.4,0) node[below left, font=\scriptsize]{kWh};
\draw[->, popdark] (0,0) -- (0,3.2);
\foreach \x/\h in {0.6/0.4, 1.4/1.6, 2.2/2.6, 3.0/1.6, 3.8/0.4}{
  \fill[popblue!55, draw=popblue] (\x,0) rectangle (\x+0.8,\h);}
\draw[poporange, very thick, dashed] (2.6,0) -- (2.6,3.0) node[above, font=\scriptsize]{$\bar{x}$};
\node[font=\small\bfseries, popblue] at (2.6,-0.9) {Série resserrée : petit $\sigma$};
\end{scope}
% Histogramme étalé
\begin{scope}[xshift=7.2cm]
\draw[->, popdark] (0,0) -- (5.4,0) node[below left, font=\scriptsize]{kWh};
\draw[->, popdark] (0,0) -- (0,3.2);
\foreach \x/\h in {0.2/1.3, 1.0/1.7, 1.8/2.0, 2.6/2.0, 3.4/1.7, 4.2/1.3}{
  \fill[poporange!55, draw=poporange] (\x,0) rectangle (\x+0.8,\h);}
\draw[popblue, very thick, dashed] (2.6,0) -- (2.6,3.0) node[above, font=\scriptsize]{$\bar{x}$};
\node[font=\small\bfseries, poporange] at (2.6,-0.9) {Série étalée : grand $\sigma$};
\end{scope}
\end{tikzpicture}
\end{center}
\legende{Deux séries de même moyenne $\bar{x}$ : à gauche, les valeurs se concentrent autour de la moyenne (faible dispersion) ; à droite, elles s'en éloignent fortement (forte dispersion).}
\end{popfigure}

\courssub{Les écarts à la moyenne}

Pour toute la suite, on reprend notre série de référence : les consommations mensuelles d'électricité (en kWh) de $40$ abonnés ENEO d'un quartier de Yaoundé, regroupées en classes. On a calculé dans les sections précédentes la moyenne $\bar{x} = 167{,}5$~kWh, obtenue à partir des centres $c_i$ des classes.

\begin{definition}[Écart et écart absolu à la moyenne]
Soit une série regroupée en classes de centres $c_1, c_2, \ldots, c_p$, d'effectifs $n_1, n_2, \ldots, n_p$, d'effectif total $N$ et de moyenne $\bar{x}$. Pour chaque classe :
\begin{itemize}
\item l'\cle{écart à la moyenne} est le nombre $c_i - \bar{x}$ ;
\item l'\cle{écart absolu à la moyenne} est le nombre $|c_i - \bar{x}|$.
\end{itemize}
\end{definition}

L'écart $c_i - \bar{x}$ peut être positif (la classe est au-dessus de la moyenne) ou négatif (en dessous). Si l'on faisait la somme de tous ces écarts pondérés par les effectifs, on obtiendrait \textbf{zéro} : les écarts positifs compensent exactement les écarts négatifs, c'est d'ailleurs la définition même de la moyenne. En effet :
\[ \sum_{i=1}^{p} n_i(c_i - \bar{x}) = \sum_{i=1}^{p} n_i c_i - \bar{x}\sum_{i=1}^{p} n_i = N\bar{x} - \bar{x}\,N = 0. \]
Cette somme nulle ne mesure donc aucune dispersion ! C'est pourquoi on « neutralise » les signes, de deux façons possibles :
\begin{itemize}
\item en prenant la \textbf{valeur absolue} : cela conduit à l'\emph{écart moyen} ;
\item en prenant le \textbf{carré} : cela conduit à la \emph{variance} et à l'\emph{écart-type}.
\end{itemize}

\courssub{L'écart moyen}

\begin{definition}[Écart moyen]
L'\cle{écart moyen} d'une série regroupée en classes est la moyenne des écarts absolus à la moyenne, pondérée par les effectifs :
\[ e = \frac{1}{N}\sum_{i=1}^{p} n_i\,|c_i - \bar{x}|. \]
\end{definition}

L'écart moyen s'interprète très naturellement : c'est la distance \emph{moyenne} entre les valeurs de la série et leur moyenne. Il s'exprime dans la même unité que les données (ici, en kWh), ce qui le rend facile à commenter. Son défaut : les valeurs absolues se manipulent mal dans les calculs algébriques, ce qui explique que les mathématiciens lui préfèrent la variance.

\courssub{Variance et écart-type}

\begin{definition}[Variance et écart-type]
Soit une série regroupée en classes de centres $c_i$, d'effectifs $n_i$, d'effectif total $N$ et de moyenne $\bar{x}$.
\begin{itemize}
\item La \cle{variance} est la moyenne des carrés des écarts à la moyenne :
\[ V = \frac{1}{N}\sum_{i=1}^{p} n_i\,(c_i - \bar{x})^2 = \sum_{i=1}^{p} f_i\,(c_i - \bar{x})^2, \]
où $f_i = \dfrac{n_i}{N}$ désigne la fréquence de la classe numéro $i$.
\item L'\cle{écart-type}, noté $\sigma$ (lettre grecque « sigma »), est la racine carrée de la variance :
\[ \sigma = \sqrt{V}. \]
\end{itemize}
\end{definition}

\begin{aretenir}[Interprétation de l'écart-type]
L'écart-type mesure la \emph{dispersion moyenne} des valeurs autour de la moyenne :
\begin{itemize}
\item $\sigma$ \textbf{petit} : les valeurs sont concentrées autour de $\bar{x}$, la série est \textbf{homogène} ;
\item $\sigma$ \textbf{grand} : les valeurs sont très étalées, la série est \textbf{hétérogène}.
\end{itemize}
Comme l'écart moyen, $\sigma$ s'exprime dans la même unité que les données. La variance $V$, elle, s'exprime dans l'unité \emph{au carré} (ici des kWh$^2$), ce qui la rend difficile à interpréter directement : c'est pour cela qu'on lui préfère $\sigma$ dans les commentaires.
\end{aretenir}

\courssub{La formule de König-Huygens}

Calculer $V$ avec la définition oblige à calculer chaque écart $c_i - \bar{x}$, souvent non entier, puis son carré : fastidieux et source d'erreurs. Le théorème suivant fournit un chemin de calcul beaucoup plus confortable.

\begin{propriete}[Formule de König-Huygens]
La variance d'une série statistique est égale à la \emph{moyenne des carrés} moins le \emph{carré de la moyenne} :
\[ V = \frac{1}{N}\sum_{i=1}^{p} n_i\,c_i^2 \;-\; \bar{x}^2. \]
Cette démonstration est formatrice et peut être demandée au Probatoire : sache la refaire.
\end{propriete}

\begin{experience}[Démonstration guidée de la formule de König-Huygens]
Partons de la définition et développons le carré $(c_i - \bar{x})^2 = c_i^2 - 2\bar{x}\,c_i + \bar{x}^2$ :
\begin{align*}
V &= \frac{1}{N}\sum_{i=1}^{p} n_i\,(c_i - \bar{x})^2 \\
  &= \frac{1}{N}\sum_{i=1}^{p} n_i\left(c_i^2 - 2\bar{x}\,c_i + \bar{x}^2\right) \\
  &= \frac{1}{N}\sum_{i=1}^{p} n_i c_i^2 \;-\; 2\bar{x}\cdot\frac{1}{N}\sum_{i=1}^{p} n_i c_i \;+\; \bar{x}^2\cdot\frac{1}{N}\sum_{i=1}^{p} n_i.
\end{align*}
Or, par définition de la moyenne, $\dfrac{1}{N}\displaystyle\sum_{i=1}^{p} n_i c_i = \bar{x}$, et $\dfrac{1}{N}\displaystyle\sum_{i=1}^{p} n_i = \dfrac{N}{N} = 1$. En remplaçant :
\[ V = \frac{1}{N}\sum_{i=1}^{p} n_i c_i^2 - 2\bar{x}\cdot\bar{x} + \bar{x}^2 \cdot 1 = \frac{1}{N}\sum_{i=1}^{p} n_i c_i^2 - \bar{x}^2. \]
La formule est démontrée. Retiens la phrase-clé : \emph{« moyenne des carrés moins carré de la moyenne »}.
\end{experience}

\begin{attention}[Les trois erreurs classiques]
\begin{enumerate}
\item \textbf{Oublier les effectifs.} La formule contient $n_i$ : on calcule $\sum n_i c_i^2$ et non $\sum c_i^2$. Chaque classe compte autant de fois qu'elle contient d'individus.
\item \textbf{Confondre $V$ et $\sigma$.} La variance est en kWh$^2$, l'écart-type en kWh. Si l'énoncé demande l'écart-type, n'oublie pas la racine carrée finale !
\item \textbf{Prendre la racine carrée trop tôt.} On calcule d'abord $V$ en entier, puis $\sigma = \sqrt{V}$ à la toute fin. Jamais l'inverse.
\end{enumerate}
\end{attention}

\courssub{Organisation des calculs en tableau}

\begin{methode}[Calculer $e$, $V$ et $\sigma$ à l'aide d'un tableau]
\begin{enumerate}
\item Calculer la moyenne $\bar{x} = \dfrac{1}{N}\sum n_i c_i$ (travail déjà fait à la section 2).
\item Compléter un tableau à cinq colonnes de calcul : $c_i$, $n_i c_i$, $c_i - \bar{x}$, $n_i|c_i - \bar{x}|$, $n_i(c_i - \bar{x})^2$ (ou $n_i c_i^2$ si l'on préfère König-Huygens).
\item Faire la \textbf{somme} de chaque colonne utile.
\item En déduire $e = \dfrac{1}{N}\sum n_i|c_i - \bar{x}|$ et $V = \dfrac{1}{N}\sum n_i(c_i - \bar{x})^2$.
\item Conclure par $\sigma = \sqrt{V}$, puis \textbf{rédiger une phrase d'interprétation}.
\end{enumerate}
\end{methode}

\begin{exemplebox}[Tableau de calcul complet pour la série ENEO]
Rappel des données : $\bar{x} = 167{,}5$~kWh, $N = 40$.
\begin{center}
\begin{tabularx}{\linewidth}{|c|c|c|X|c|c|}
\hline
\rowcolor{popblueL} Classe & $c_i$ & $n_i$ & \centering $c_i - \bar{x}$ & $n_i|c_i-\bar{x}|$ & $n_i(c_i-\bar{x})^2$ \\
\hline
$[50\,;\,100[$ & 75 & 6 & $-92{,}5$ & $555$ & \num{51337,5} \\
\hline
$[100\,;\,150[$ & 125 & 10 & $-42{,}5$ & $425$ & \num{18062,5} \\
\hline
$[150\,;\,200[$ & 175 & 12 & $7{,}5$ & $90$ & $675$ \\
\hline
$[200\,;\,250[$ & 225 & 8 & $57{,}5$ & $460$ & \num{26450} \\
\hline
$[250\,;\,300[$ & 275 & 4 & $107{,}5$ & $430$ & \num{46225} \\
\hline
\rowcolor{popgoldL} \textbf{Total} & & $40$ & & \num{1960} & \num{142750} \\
\hline
\end{tabularx}
\end{center}
Détaillons par exemple la première ligne : $c_1 - \bar{x} = 75 - 167{,}5 = -92{,}5$ ; puis $n_1|c_1 - \bar{x}| = 6 \times 92{,}5 = 555$ ; enfin $n_1(c_1-\bar{x})^2 = 6 \times (92{,}5)^2 = 6 \times \num{8556,25} = \num{51337,5}$. Les autres lignes se calculent de la même manière.

Vérifions la cohérence avec König-Huygens : $\sum n_i c_i^2 = \num{33750} + \num{156250} + \num{367500} + \num{405000} + \num{302500} = \num{1265000}$, donc
\[ V = \frac{\num{1265000}}{40} - (167{,}5)^2 = \num{31625} - \num{28056,25} = \num{3568,75}. \]
On retrouve bien le même résultat qu'avec la définition : $\dfrac{\num{142750}}{40} = \num{3568,75}$. Les deux méthodes concordent, c'est un excellent réflexe de vérification.
\end{exemplebox}

\begin{exoresolu}[Écart moyen, variance et écart-type des consommations ENEO]
Énoncé. Pour la série des consommations de $40$ abonnés ENEO ($\bar{x} = 167{,}5$~kWh, voir tableau ci-dessus) :
\begin{enumerate}
\item Calculer l'écart moyen $e$.
\item Calculer la variance $V$ puis l'écart-type $\sigma$ (arrondi à $0{,}1$ près).
\item Un deuxième quartier donne $\bar{x}' = 165$~kWh et $\sigma' = 22{,}4$~kWh. Comparer les deux quartiers.
\end{enumerate}
\tcblower
Solution détaillée.
\begin{enumerate}
\item D'après la colonne des écarts absolus pondérés :
\[ e = \frac{1}{N}\sum_{i=1}^{5} n_i|c_i - \bar{x}| = \frac{\num{1960}}{40} = 49. \]
En moyenne, la consommation d'un abonné s'écarte de $49$~kWh de la moyenne.
\item D'après la dernière colonne :
\[ V = \frac{1}{N}\sum_{i=1}^{5} n_i(c_i - \bar{x})^2 = \frac{\num{142750}}{40} = \num{3568,75} \ \text{(kWh}^2\text{)}, \]
\[ \sigma = \sqrt{\num{3568,75}} \approx 59{,}7 \ \text{kWh}. \]
\item Les deux quartiers ont pratiquement la même consommation moyenne ($167{,}5$ contre $165$~kWh). En revanche, le premier quartier a un écart-type de $59{,}7$~kWh, presque trois fois plus grand que celui du second ($22{,}4$~kWh). Les consommations du premier quartier sont donc beaucoup plus \textbf{dispersées} : il y a de fortes inégalités de consommation entre abonnés (petits consommateurs et gros consommateurs cohabitent). Le second quartier est plus \textbf{homogène} : les ménages y ont des profils de consommation semblables. Pour ENEO, cela signifie que la facturation moyenne y est plus prévisible.
\end{enumerate}
\end{exoresolu}

\begin{attention}[Rédaction de l'interprétation]
Une comparaison de dispersion n'a de sens que si les moyennes sont du même ordre de grandeur et si les données sont dans la même unité. Rédige toujours ta conclusion en deux temps : d'abord comparer les moyennes (niveau global), puis comparer les écarts-types (homogénéité). Et n'oublie jamais l'unité dans la réponse finale.
\end{attention}

\begin{savaistu}[{L'intervalle $\left[\bar{x} - 2\sigma \,;\, \bar{x} + 2\sigma\right]$}]
Pour une série statistique « bien répartie » (en forme de cloche, ce que tu formaliseras en Terminale avec la \emph{loi normale}), environ $95\,\%$ des valeurs se trouvent dans l'intervalle $\left[\bar{x} - 2\sigma \,;\, \bar{x} + 2\sigma\right]$. Pour notre série ENEO : $\left[167{,}5 - 2\times 59{,}7 \,;\, 167{,}5 + 2\times 59{,}7\right] \approx [48 \,;\, 287]$~kWh, qui contient effectivement la quasi-totalité de nos classes. Ce résultat, appelé \emph{intervalle de fluctuation}, est à la base des sondages d'opinion : quand un institut annonce un résultat « à plus ou moins $3\,\%$ près » avant une élection, il utilise exactement cette idée ! L'écart-type a été introduit par le statisticien britannique Karl Pearson à la fin du XIX\textsuperscript{e} siècle, et il reste aujourd'hui l'indicateur de dispersion le plus utilisé au monde, de la météorologie à la finance.
\end{savaistu}

\begin{exercice}[À toi de jouer]
Une coopérative de planteurs de cacao de la région du Centre relève les masses (en kg) de cabosses récoltées par $30$ planteurs, regroupées en classes de centres $c_i$ : $50$, $150$, $250$, $350$, avec les effectifs $4$, $12$, $9$, $5$.
\begin{enumerate}
\item Calculer la moyenne $\bar{x}$ de la série.
\item Construire le tableau de calcul complet et déterminer l'écart moyen $e$.
\item Calculer la variance $V$ par la formule de König-Huygens, puis l'écart-type $\sigma$ (arrondi à $0{,}1$).
\item Une autre coopérative a une moyenne de $210$~kg et un écart-type de $35$~kg. Laquelle des deux coopératives présente les récoltes les plus homogènes ?
\end{enumerate}
\end{exercice}

\begin{corrige}[Exercice : coopérative de planteurs]
\begin{enumerate}
\item $\sum n_i c_i = 4\times 50 + 12\times 150 + 9\times 250 + 5\times 350 = 200 + \num{1800} + \num{2250} + \num{1750} = \num{6000}$, donc $\bar{x} = \dfrac{\num{6000}}{30} = 200$~kg.
\item Écarts : $c_i - \bar{x} = -150\,;\,-50\,;\,50\,;\,150$. D'où $n_i|c_i-\bar{x}| = 600\,;\,600\,;\,450\,;\,750$, total $\num{2400}$, et $e = \dfrac{\num{2400}}{30} = 80$~kg.
\item $\sum n_i c_i^2 = 4\times \num{2500} + 12\times \num{22500} + 9\times \num{62500} + 5\times \num{122500} = \num{10000} + \num{270000} + \num{562500} + \num{612500} = \num{1455000}$.
\[ V = \frac{\num{1455000}}{30} - 200^2 = \num{48500} - \num{40000} = \num{8500}, \qquad \sigma = \sqrt{\num{8500}} \approx 92{,}2 \ \text{kg}. \]
\item La première coopérative a un écart-type d'environ $92{,}2$~kg, très supérieur à celui de la seconde ($35$~kg), pour des moyennes comparables. C'est donc la \textbf{seconde coopérative} qui présente les récoltes les plus homogènes.
\end{enumerate}
\end{corrige}

\begin{aretenir}[L'essentiel de la section]
\begin{itemize}
\item La moyenne ne suffit pas : deux séries de même moyenne peuvent être très différemment réparties.
\item Écart moyen : $e = \dfrac{1}{N}\sum n_i|c_i - \bar{x}|$ ; variance : $V = \dfrac{1}{N}\sum n_i(c_i - \bar{x})^2 = \dfrac{1}{N}\sum n_i c_i^2 - \bar{x}^2$ (König-Huygens) ; écart-type : $\sigma = \sqrt{V}$.
\item $\sigma$ s'exprime dans l'unité des données : petit $\sigma$ = série homogène, grand $\sigma$ = série dispersée.
\item Les calculs s'organisent toujours dans un tableau, et toute étude se termine par une phrase d'interprétation en contexte.
\end{itemize}
\end{aretenir}

\coursec{Représentations graphiques : histogramme et courbes cumulatives}

Dans les sections précédentes, nous avons appris à \emph{mesurer} la dispersion d'une série regroupée en classes à l'aide de nombres : écart moyen, variance, écart-type. Ces indicateurs sont précis, mais ils restent abstraits. Un chef d'entreprise, un agent de la SONEL ou un élève qui prépare le Probatoire a souvent besoin de \emph{voir} la répartition d'un coup d'œil : où se concentrent les données ? Sont-elles étalées ou resserrées ? Où se situe la médiane ? C'est exactement le rôle des représentations graphiques que nous allons construire dans cette section : l'\cle{histogramme}, qui donne une image fidèle de la répartition (respectant les aires), et les \cle{courbes cumulatives}, qui permettent de lire graphiquement la médiane et les quantiles.

Toute la section s'appuiera sur l'exemple fil conducteur suivant, déjà utilisé dans les sections précédentes.

\begin{exemplebox}[Consommations électriques SONEL]
La SONEL a relevé la consommation mensuelle (en kWh) de $90$ abonnés d'un quartier de Douala. Les résultats sont regroupés en classes :

\begin{center}
\begin{tabularx}{\linewidth}{|l|X|X|X|X|X|}\hline
\rowcolor{popblueL} Classe (kWh) & $[0;50[$ & $[50;100[$ & $[100;150[$ & $[150;200[$ & $[200;250[$ \\ \hline
Effectif $n_i$ & $12$ & $20$ & $30$ & $18$ & $10$ \\ \hline
ECC & $12$ & $32$ & $62$ & $80$ & $90$ \\ \hline
ECD & $90$ & $78$ & $58$ & $28$ & $10$ \\ \hline
\end{tabularx}
\end{center}

L'effectif total est $N = 90$. La classe modale est $[100;150[$ (plus grand effectif) et la médiane, obtenue par interpolation linéaire dans la section 3, vaut $Me \approx 121{,}7$ kWh.
\end{exemplebox}

\courssub{L'histogramme à classes d'amplitudes égales}

\textbf{Intuition.} Un diagramme en bâtons convient à un caractère discret ; mais ici, le caractère est \emph{continu} et regroupé en classes : les valeurs « collent » les unes aux autres. Il est donc naturel de \emph{coller les bâtons} : on obtient des rectangles juxtaposés, dont la base est la classe et dont la hauteur traduit l'importance de la classe. C'est l'histogramme.

\begin{definition}[Histogramme d'une série regroupée en classes]
L'\cle{histogramme} d'une série statistique regroupée en classes est un graphique formé de rectangles juxtaposés :
\begin{itemize}
\item la \textbf{base} de chaque rectangle est l'intervalle-classe, porté sur l'axe des abscisses ;
\item lorsque toutes les classes ont la \textbf{même amplitude}, la \textbf{hauteur} de chaque rectangle est l'effectif $n_i$ (ou la fréquence $f_i$) de la classe ;
\item lorsque les classes ont des \textbf{amplitudes inégales}, c'est l'\textbf{aire} de chaque rectangle qui est proportionnelle à l'effectif (voir la sous-section suivante).
\end{itemize}
\end{definition}

\begin{methode}[Construire un histogramme à amplitudes égales]
\begin{enumerate}
\item Tracer deux axes gradués : les bornes des classes en abscisse, les effectifs (ou fréquences) en ordonnée.
\item Pour chaque classe $[a_i ; a_{i+1}[$, tracer le rectangle de base $[a_i ; a_{i+1}]$ et de hauteur $n_i$.
\item Ne laisser \textbf{aucun espace} entre les rectangles (le caractère est continu).
\item Indiquer les unités et une légende.
\end{enumerate}
\end{methode}

\begin{exemplebox}[Histogramme des consommations SONEL]
Les cinq classes ont toutes pour amplitude $50$ kWh : on porte donc les effectifs en hauteur. On utilise l'option \texttt{ybar interval} de \texttt{pgfplots} qui trace exactement les rectangles sur les bornes de classes.

\begin{popfigure}
\begin{tikzpicture}
\begin{axis}[
width=0.9\linewidth, height=6.5cm,
ybar interval, % Trace des rectangles [x_i, x_{i+1}] de hauteur y_i
xmin=0, xmax=250, ymin=0, ymax=34,
xlabel={Consommation (kWh)}, ylabel={Effectif},
xtick={0,50,100,150,200,250},
ytick={0,10,20,30},
grid=major, grid style={popline!50},
every axis plot/.append style={fill=popblue!70, draw=popdark},
]
% Coordonnées : (borne_inf, effectif) ... (borne_sup_derniere, 0)
\addplot+[ybar interval] coordinates {(0,12) (50,20) (100,30) (150,18) (200,10) (250,0)};
\end{axis}
\end{tikzpicture}
\legende{Histogramme des consommations mensuelles de 90 abonnés SONEL (classes d'amplitudes égales).}
\end{popfigure}

On lit immédiatement que le rectangle le plus haut correspond à la classe $[100;150[$ : c'est la \cle{classe modale}. Le \cle{mode} est le centre de cette classe : $Mo = \dfrac{100+150}{2} = 125$ kWh.
\end{exemplebox}

\begin{aretenir}[Lecture de la classe modale sur l'histogramme (amplitudes égales)]
Pour des classes d'\textbf{amplitudes égales}, la classe modale est la classe du rectangle \textbf{le plus haut}. Le mode est le centre de cette classe.
\end{aretenir}

\courssub{Classes d'amplitudes inégales : la densité d'effectif}

\textbf{Intuition.} Imaginons deux classes : $[0;20[$ contenant $10$ élèves et $[20;60[$ contenant $20$ élèves. Si l'on trace des rectangles de hauteurs $10$ et $20$, le second rectangle a une base \emph{trois fois plus large} et une hauteur \emph{deux fois plus grande} : son aire est \emph{six fois} plus grande, alors qu'il ne représente que deux fois plus d'élèves. Le graphique ment ! Pour rester honnête, il faut « diluer » l'effectif sur toute la largeur de la classe : c'est l'idée de \emph{densité}.

\begin{definition}[Densité d'effectif]
Pour une classe $[a_i ; a_{i+1}[$ d'effectif $n_i$ et d'amplitude $a_i' = a_{i+1} - a_i$, la \cle{densité d'effectif} est
\[ d_i = \frac{n_i}{a_i'} = \frac{\text{effectif de la classe}}{\text{amplitude de la classe}}. \]
Dans un histogramme à classes d'amplitudes inégales, la \textbf{hauteur} du rectangle est la densité $d_i$, de sorte que :
\[ \text{aire du rectangle} = \text{base} \times \text{hauteur} = a_i' \times \frac{n_i}{a_i'} = n_i. \]
L'\textbf{aire} de chaque rectangle est donc exactement proportionnelle à l'effectif de la classe (elle est égale à l'effectif si l'on prend l'unité d'aire comme le carré de l'unité d'abscisse).
\end{definition}

\begin{methode}[Construire un histogramme à amplitudes inégales en 4 étapes]
\begin{enumerate}
\item \textbf{Calculer les densités} : pour chaque classe, $d_i = n_i / \text{amplitude}$.
\item \textbf{Choisir l'échelle} : graduer l'axe des ordonnées en densité (effectif par unité du caractère), en tenant compte de la plus grande densité.
\item \textbf{Tracer les rectangles} : base = classe, hauteur = densité $d_i$.
\item \textbf{Légender} : préciser que l'aire (et non la hauteur) représente l'effectif, et indiquer l'unité d'aire.
\end{enumerate}
\end{methode}

\begin{exemplebox}[Histogramme à amplitudes inégales : distances de livraison]
Une entreprise de livraison de Yaoundé relève les distances parcourues par ses motos :

\begin{center}
\begin{tabularx}{\linewidth}{|l|X|X|X|}\hline
\rowcolor{popblueL} Distance (km) & $[0;20[$ & $[20;40[$ & $[40;100[$ \\ \hline
Effectif $n_i$ & $10$ & $16$ & $30$ \\ \hline
Amplitude $a_i'$ & $20$ & $20$ & $60$ \\ \hline
Densité $d_i = n_i/a_i'$ & $0{,}5$ & $0{,}8$ & $0{,}5$ \\ \hline
\end{tabularx}
\end{center}

\begin{popfigure}
\begin{tikzpicture}[x=0.075cm, y=4cm]
% axes
\draw[->, thick] (0,0) -- (110,0) node[below left, font=\small] {Distance (km)};
\draw[->, thick] (0,0) -- (0,1.05) node[above, font=\small] {Densité (effectif/km)};
\foreach \x in {0,20,40,60,80,100} \draw (\x,0.02) -- (\x,-0.02) node[below, font=\small] {\x};
\foreach \y/\lab in {0.5/0{,}5, 0.8/0{,}8} \draw (1.5,\y) -- (-1.5,\y) node[left, font=\small] {\lab};
% rectangles (hauteurs = densités)
\fill[popblue!60, draw=popdark, thick] (0,0) rectangle (20,0.5);
\fill[poporange!60, draw=popdark, thick] (20,0) rectangle (40,0.8);
\fill[popgreen!60, draw=popdark, thick] (40,0) rectangle (100,0.5);
% annotations
\node[font=\small, text=popdark] at (10,0.25) {$n_1=10$};
\node[font=\small, text=popdark] at (30,0.4) {$n_2=16$};
\node[font=\small, text=popdark] at (70,0.25) {$n_3=30$};
% flèches aire
\draw[<->, popmuted, thick] (10,0.55) -- (10,0.65) node[midway, above, font=\tiny] {Aire $=10$};
\draw[<->, popmuted, thick] (30,0.85) -- (30,0.95) node[midway, above, font=\tiny] {Aire $=16$};
\draw[<->, popmuted, thick] (70,0.55) -- (70,0.65) node[midway, above, font=\tiny] {Aire $=30$};
\end{tikzpicture}
\legende{Histogramme correct : les hauteurs sont les densités. Le rectangle le plus haut (densité $0{,}8$) est la classe $[20;40[$ : c'est la classe modale. Le 3\up{e} rectangle a la plus grande aire (effectif $30$) mais une densité plus faible car la classe est large.}
\end{popfigure}

\textbf{Analyse.} La classe modale est $[20;40[$ (densité maximale $0{,}8$), \emph{pas} la classe $[40;100[$ qui a le plus grand effectif ($30$). Le mode (valeur approchée) est le centre de la classe modale : $30$ km. Cela signifie que les motos roulent \emph{plus fréquemment} sur des distances courtes (autour de $30$ km) même si le nombre total de trajets longs est plus élevé.
\end{exemplebox}

\begin{attention}[C'est l'AIRE qui compte !]
L'erreur la plus fréquente au Probatoire est de représenter des hauteurs égales aux effectifs alors que les classes ont des \textbf{amplitudes inégales}. Retiens :
\begin{itemize}
\item amplitudes égales $\Rightarrow$ hauteur $=$ effectif (ou fréquence) ;
\item amplitudes inégales $\Rightarrow$ hauteur $=$ \textbf{densité} $d_i = n_i/\text{amplitude}$, et c'est l'\textbf{aire} qui est proportionnelle à l'effectif.
\end{itemize}
Conséquence : avec des amplitudes inégales, la classe modale est la classe de plus grande \textbf{densité} (rectangle le plus haut), pas nécessairement celle du plus grand effectif !
\end{attention}

\courssub{Courbes des effectifs et fréquences cumulés}

\textbf{Intuition.} Réponds à cette question : « combien d'abonnés consomment \emph{moins de} $150$ kWh ? ». La réponse est un effectif cumulé croissant : $62$. Si l'on pose cette question pour \emph{toutes} les bornes et que l'on relie les réponses, on obtient une courbe qui « monte » : la courbe des ECC. De même, « combien consomment \emph{au moins} $100$ kWh ? » donne la courbe des ECD, qui « descend ».

\begin{definition}[Courbes cumulatives (ogives)]
Pour une série regroupée en classes $[a_i ; a_{i+1}[$ :
\begin{itemize}
\item la \cle{courbe des effectifs cumulés croissants} (ECC) est la ligne brisée joignant les points $\left(a_{i+1}\,;\, \text{ECC}_i\right)$, en partant du point $(a_1\,;\,0)$ : on porte la \textbf{borne supérieure} de chaque classe ;
\item la \cle{courbe des effectifs cumulés décroissants} (ECD) est la ligne brisée joignant les points $\left(a_i\,;\, \text{ECD}_i\right)$, en terminant par $(a_{\max}\,;\,0)$ : on porte la \textbf{borne inférieure} de chaque classe ;
\item les \cle{courbes de fréquences cumulées} (FCC et FCD) s'obtiennent de la même manière en remplaçant les effectifs par les fréquences cumulées, exprimées en $\%$.
\end{itemize}
Ces courbes sont aussi appelées \cle{ogives} (par analogie avec les arcs en ogive de l'architecture gothique).
\end{definition}

\begin{attention}[Borne supérieure ou borne inférieure ?]
Ne confonds jamais les deux conventions :
\begin{itemize}
\item ECC / FCC $\longrightarrow$ on place le point à la \textbf{borne supérieure} de la classe (« moins de $a_{i+1}$ ») ;
\item ECD / FCD $\longrightarrow$ on place le point à la \textbf{borne inférieure} de la classe (« au moins $a_i$ »).
\end{itemize}
Inverser les bornes décale toute la courbe et fausse la lecture de la médiane.
\end{attention}

\begin{exemplebox}[Courbes ECC et ECD de l'exemple SONEL]
On place les points $(0;0)$, $(50;12)$, $(100;32)$, $(150;62)$, $(200;80)$, $(250;90)$ pour les ECC, et $(0;90)$, $(50;78)$, $(100;58)$, $(150;28)$, $(200;10)$, $(250;0)$ pour les ECD.

\begin{popfigure}
\begin{tikzpicture}
\begin{axis}[
width=0.9\linewidth, height=7.5cm,
xmin=0, xmax=260, ymin=0, ymax=100,
xlabel={Consommation (kWh)}, ylabel={Effectifs cumulés},
xtick={0,50,100,150,200,250},
ytick={0,10,20,...,90},
grid=major, grid style={popline!40},
legend style={at={(0.02,0.98)}, anchor=north west, font=\small},
]
\addplot[popblue, very thick, mark=*, mark size=1.5pt] coordinates {(0,0) (50,12) (100,32) (150,62) (200,80) (250,90)};
\addlegendentry{ECC}
\addplot[poporange, very thick, mark=square*, mark size=1.5pt] coordinates {(0,90) (50,78) (100,58) (150,28) (200,10) (250,0)};
\addlegendentry{ECD}
% médiane
\addplot[popgreen, dashed, thick] coordinates {(0,45) (121.67,45)};
\addplot[popgreen, dashed, thick] coordinates {(121.67,0) (121.67,45)};
\addplot[popgreen, mark=*, mark size=2pt, only marks] coordinates {(121.67,45)};
\node[popgreen, font=\small, anchor=south] at (axis cs:121.67,45) {$Me \approx 121{,}7$};
\node[popgreen, font=\small, anchor=west] at (axis cs:5,45) {$N/2 = 45$};
\end{axis}
\end{tikzpicture}
\legende{Courbes des effectifs cumulés croissants et décroissants. Elles se coupent au point d'abscisse la médiane.}
\end{popfigure}
\end{exemplebox}

\begin{propriete}[Intersection des courbes ECC et ECD (admise)]
Les courbes des effectifs cumulés croissants et décroissants se coupent en un point dont l'\textbf{abscisse est la médiane} $Me$ de la série (et dont l'ordonnée est $N/2$). Cette propriété, admise à ce niveau, est illustrée par la figure ci-dessus : au point d'intersection, exactement la moitié des individus se situe en dessous et l'autre moitié au dessus, ce qui est précisément la définition de la médiane.
\end{propriete}

\begin{methode}[Lecture graphique de la médiane]
Deux méthodes équivalentes :
\begin{enumerate}
\item \textbf{Avec la courbe des ECC (ou FCC) seule} : placer $\dfrac{N}{2}$ (ou $50\,\%$) sur l'axe des ordonnées, tracer l'horizontale jusqu'à la courbe, puis descendre verticalement vers l'axe des abscisses : on lit $Me$.
\item \textbf{Avec les deux courbes} : l'abscisse du point d'intersection des courbes ECC et ECD (ou FCC et FCD) est $Me$.
\end{enumerate}
Remarque : la valeur lue graphiquement coïncide avec celle obtenue par interpolation linéaire, car la courbe cumulative est précisément constituée de segments de droite.
\end{methode}

\textbf{Interprétation de la pente.} Sur la courbe des ECC, le segment relatif à la classe $[a_i ; a_{i+1}[$ a pour pente
\[ \frac{\text{ECC}_i - \text{ECC}_{i-1}}{a_{i+1}-a_i} = \frac{n_i}{\text{amplitude}} = d_i. \]
La pente d'un segment est donc exactement la \cle{densité d'effectif} de la classe : \textbf{plus le segment est raide, plus la classe est chargée}. Dans l'exemple SONEL, le segment le plus pentu est celui de $[100;150[$ (pente $30/50 = 0{,}6$), ce qui confirme visuellement la classe modale.

\begin{exoresolu}[Fréquences cumulées et lecture graphique]
\textbf{Énoncé.} Pour l'exemple SONEL : 1) Dresser le tableau des fréquences cumulées croissantes (en $\%$). 2) Tracer la courbe des FCC et lire graphiquement la médiane. 3) Déterminer graphiquement le pourcentage d'abonnés consommant moins de $175$ kWh.
\tcblower
\textbf{Solution.} 1) On divise chaque ECC par $N = 90$ :
\begin{center}
\begin{tabularx}{\linewidth}{|l|X|X|X|X|X|}\hline
\rowcolor{popblueL} Borne sup. (kWh) & $50$ & $100$ & $150$ & $200$ & $250$ \\ \hline
FCC ($\%$) & $13{,}3$ & $35{,}6$ & $68{,}9$ & $88{,}9$ & $100$ \\ \hline
\end{tabularx}
\end{center}
2) On place les points $(50\,;\,13{,}3)$, $(100\,;\,35{,}6)$, $(150\,;\,68{,}9)$, $(200\,;\,88{,}9)$, $(250\,;\,100)$ et on les relie. L'horizontale d'ordonnée $50\,\%$ coupe la courbe sur le segment de $[100;150[$ ; la verticale descendante donne $Me \approx 122$ kWh, conforme au calcul par interpolation :
\[ Me = 100 + 50\times\frac{50-35{,}6}{68{,}9-35{,}6} \approx 121{,}7 \text{ kWh}. \]
3) Pour $x = 175$ (milieu du segment $[150;200[$), on lit l'ordonnée par interpolation linéaire sur la courbe des FCC :
\[ \text{FCC}(175) = 68{,}9 + \frac{25}{50}\times (88{,}9 - 68{,}9) = 68{,}9 + 10 = 78{,}9\,\%. \]
Environ $79\,\%$ des abonnés consomment moins de $175$ kWh.
\end{exoresolu}

\begin{exercice}[Histogramme à amplitudes inégales]
Les âges des clients d'une agence MTN de Bafoussam sont répartis ainsi :
$[15;25[$ : $40$ clients ; $[25;35[$ : $60$ clients ; $[35;50[$ : $45$ clients ; $[50;70[$ : $20$ clients.
\begin{enumerate}
\item Calculer la densité d'effectif de chaque classe.
\item Construire l'histogramme (on pourra utiliser un papier millimétré ou un logiciel).
\item Déterminer la classe modale. Était-ce la classe du plus grand effectif ?
\end{enumerate}
\end{exercice}

\begin{corrige}[Exercice 1]
1) Densités :
$d_1 = 40/10 = 4$ clients/an ;
$d_2 = 60/10 = 6$ clients/an ;
$d_3 = 45/15 = 3$ clients/an ;
$d_4 = 20/20 = 1$ client/an.

2) On trace quatre rectangles de bases $10$, $10$, $15$, $20$ et de hauteurs respectives $4$, $6$, $3$, $1$. Les rectangles sont juxtaposés. L'axe des ordonnées est gradué en « clients par an ».

3) Le rectangle le plus haut est celui de $[25;35[$ (densité $6$) : c'est la classe modale. Ici elle coïncide avec la classe du plus grand effectif ($60$), mais ce n'est pas automatique : la classe $[35;50[$ a un effectif ($45$) supérieur à celui de $[15;25[$ ($40$), pourtant sa densité ($3$) est inférieure ($4$). C'est bien la densité qui décide pour la classe modale lorsque les amplitudes sont inégales.
\end{corrige}

\begin{savaistu}[Karl Pearson et les « ogives »]
C'est le mathématicien britannique \textbf{Karl Pearson} (1857--1936), fondateur de la statistique moderne, qui a popularisé l'histogramme et lui a donné son nom vers 1895. Les courbes cumulatives, elles, portent le joli nom d'\textbf{ogives}, par analogie avec les arcs en ogive de l'architecture gothique, dont elles épousent la forme en « S ». Au Cameroun, l'INS (Institut National de la Statistique) utilise quotidiennement histogrammes et ogives pour présenter les résultats des recensements de la population : la pyramide des âges du pays n'est d'ailleurs rien d'autre qu'un double histogramme horizontal !
\end{savaistu}

\begin{aretenir}[L'essentiel de la section]
\begin{itemize}
\item \textbf{Amplitudes égales} : hauteur du rectangle $=$ effectif (ou fréquence) ; la classe modale est le rectangle le plus haut.
\item \textbf{Amplitudes inégales} : hauteur $=$ densité $d_i = n_i/\text{amplitude}$ ; c'est l'\textbf{aire} qui représente l'effectif.
\item \textbf{Courbe des ECC/FCC} : points (borne supérieure ; effectif/fréquence cumulé croissant) ; \textbf{Courbe des ECD/FCD} : points (borne inférieure ; effectif/fréquence cumulé décroissant).
\item La médiane se lit graphiquement à l'ordonnée $N/2$ (ou $50\,\%$) sur la courbe cumulative, ou à l'intersection des deux courbes.
\item La pente de chaque segment de la courbe cumulative est la densité de la classe : plus c'est raide, plus la classe est chargée.
\end{itemize}
\end{aretenir}

\coursec{Synthèse et interprétation en situation concrète}

Après avoir maîtrisé la construction de l'histogramme et des courbes cumulatives, nous disposons désormais de tous les outils pour mener une étude statistique complète. Cette section finale a pour but de \textbf{synthétiser} l'ensemble des formules, de structurer la \textbf{démarche d'investigation} et, surtout, d'apprendre à \textbf{communiquer les résultats} dans un langage adapté au décideur (direction d'entreprise, collectivité, ministère). Nous illustrerons cette démarche par une comparaison concrète de la consommation électrique dans deux quartiers de Douala, en vue d'une révision des tranches tarifaires d'ENEO.

\courssub{Bilan récapitulatif des paramètres statistiques}

Toutes les formules rencontrées dans les sections précédentes supposent que les données sont regroupées en $k$ classes $[a_i ; b_i[$ de même amplitude $h$ (sauf éventuellement la première et la dernière), de centre $x_i = \frac{a_i+b_i}{2}$, d'effectif $n_i$ et de fréquence $f_i = \frac{n_i}{N}$ ($N = \sum n_i$). Le tableau ci-dessous réunit les six grandeurs fondamentales.

\begin{aretenir}[Tableau de synthèse des formules pour séries regroupées]
\begin{center}
\begin{tabularx}{\linewidth}{|l|X|X|}\hline
\textbf{Paramètre} & \textbf{Formule} & \textbf{Interprétation rapide} \\ \hline
\rowcolor{popblueL}
\textbf{Moyenne} $\bar{x}$ & $\bar{x} = \frac{1}{N}\sum_{i=1}^k n_i x_i$ & Centre de gravité de la distribution ; valeur « typique » si la série est symétrique. \\ \hline
\textbf{Mode} (approx.) & $x_{\text{mode}} \approx x_{m} = \text{centre de la classe modale}$ & Valeur la plus fréquente (pic de l'histogramme). \\ \hline
\textbf{Médiane} $M_e$ & $M_e = a_{\text{med}} + \frac{\frac{N}{2} - N_{\text{ant}}}{n_{\text{med}}} \times h$ & Valeur centrale : 50\% des observations sont inférieures, 50\% supérieures. \\ \hline
\rowcolor{poporangeL}
\textbf{Écart moyen} $e$ & $e = \frac{1}{N}\sum_{i=1}^k n_i |x_i - \bar{x}|$ & Distance moyenne aux valeurs centrales (unité = unité des données). \\ \hline
\textbf{Variance} $V$ & $V = \frac{1}{N}\sum_{i=1}^k n_i (x_i - \bar{x})^2 = \frac{1}{N}\sum n_i x_i^2 - \bar{x}^2$ & Dispersion quadratique (unité au carré). \\ \hline
\textbf{Écart-type} $\sigma$ & $\sigma = \sqrt{V}$ & « Écart type » autour de la moyenne ; même unité que les données. \\ \hline
\end{tabularx}
\end{center}
\end{aretenir}

\begin{methode}[Fiche mémo : les 6 formules à connaître par cœur pour le Probatoire]
\begin{enumerate}
    \item \textbf{Moyenne :} $\bar{x} = \frac{\sum n_i x_i}{N}$ (pondérée par les centres de classes).
    \item \textbf{Classe modale :} Celle de plus grand effectif $n_{\max}$.
    \item \textbf{Mode :} Centre de la classe modale $x_{\text{mode}} = \frac{a_{\text{mod}}+b_{\text{mod}}}{2}$.
    \item \textbf{Médiane :} Interpolation linéaire dans la classe médiane (effectif cumulé croissant $\ge N/2$).
    \item \textbf{Variance :} $V = \frac{\sum n_i x_i^2}{N} - \bar{x}^2$ (formule de König-Huygens, plus rapide).
    \item \textbf{Écart-type :} $\sigma = \sqrt{V}$.
\end{enumerate}
\end{methode}

\courssub{Démarche complète d'étude d'une série regroupée en classes}

Face à un sujet de type « Étudier la série statistique suivante... », la rigueur impose un enchaînement logique en quatre temps. L'organigramme ci-dessous résume cette procédure standard.

\begin{popfigure}
\centering
\begin{tikzpicture}[
    node distance=1.2cm and 2cm,
    box/.style={rectangle, draw=popdark, thick, fill=popcream, rounded corners, minimum width=3.5cm, minimum height=1.2cm, align=center, font=\small},
    arrow/.style={-Stealth, thick, popdark},
    decision/.style={diamond, draw=popdark, thick, fill=poporangeL, aspect=2, minimum width=3cm, minimum height=1.5cm, align=center, font=\small}
]
\node[box, align=center] (start) {1. Construire le tableau\\ (classes, centres $x_i$, effectifs $n_i$,\\ fréquences $f_i$, effectifs cumulés)};
\node[box, below=of start, align=center] (calc) {2. Calculer les paramètres\\ \textbf{Position :} $\bar{x}$, Mode, $M_e$\\ \textbf{Dispersion :} $V$, $\sigma$ (ou $e$)};
\node[box, below=of calc, align=center] (graph) {3. Représenter graphiquement\\ \textbf{Histogramme} (densités si amplitudes inégales)\\ \textbf{Courbes cumulées} (croissante / décroissante)};
\node[box, below=of graph, align=center] (interp) {4. Interpréter et Conclure\\ \emph{« La consommation moyenne est de... »}\\ \emph{« La dispersion montre que... »}\\ \emph{Recommandation : ... »}};

\draw[arrow] (start) -- (calc);
\draw[arrow] (calc) -- (graph);
\draw[arrow] (graph) -- (interp);

% Annotations latérales
\node[left=0.5cm of start, text width=2.5cm, align=right, font=\footnotesize\textit, popmuted] {Données brutes $\to$ Tableau};
\node[left=0.5cm of calc, text width=2.5cm, align=right, font=\footnotesize\textit, popmuted] {Formules exactes (centres)};
\node[left=0.5cm of graph, text width=2.5cm, align=right, font=\footnotesize\textit, popmuted] {Visuel : forme, pics, médiane graphique};
\node[left=0.5cm of interp, text width=2.5cm, align=right, font=\footnotesize\textit, popmuted] {Phrases contextuelles + Unités};
\end{tikzpicture}
\legende{Organigramme de la démarche complète d'étude d'une série statistique regroupée en classes.}
\end{popfigure}

\courssub{Comparaison de deux séries : moyennes proches, écarts-types différents (homogénéité)}

Deux séries peuvent avoir la même moyenne mais une réalité totalement différente. L'écart-type (ou la variance) est l'indicateur clé de l'\textbf{homogénéité} : un écart-type faible signifie que les valeurs sont serrées autour de la moyenne (série homogène, prévisible) ; un écart-type fort indique une grande hétérogénéité (disparités fortes, risque élevé).

\begin{exemplebox}[Comparaison des quartiers Bonabéri et Akwa (Consommation électrique mensuelle en kWh)]\label{ex:bonaberi-akwa}
L'ENEO souhaite réviser ses tranches tarifaires. Elle dispose de deux échantillons de 200 abonnés chacun, répartis selon le tableau suivant (amplitude $h=100$ kWh).

\begin{center}
\begin{tabular}{|c|c|c|c|}\hline
\textbf{Classe (kWh)} & \textbf{Centre $x_i$} & \textbf{Bonabéri $n_i$} & \textbf{Akwa $n_i$} \\ \hline
$[100;200[$ & 150 & 5 & 12 \\ \hline
$[200;300[$ & 250 & 25 & 30 \\ \hline
$[300;400[$ & 350 & 80 & 50 \\ \hline
$[400;500[$ & 450 & 60 & 40 \\ \hline
$[500;600[$ & 550 & 20 & 35 \\ \hline
$[600;700[$ & 650 & 10 & 33 \\ \hline
\textbf{Total} & & \textbf{200} & \textbf{200} \\ \hline
\end{tabular}
\end{center}

Les paramètres calculés (détails dans l'exercice résolu \ref{exo:complete-bonaberi-akwa}) sont :

\begin{center}
\begin{tabularx}{\linewidth}{|l|X|X|}\hline
\textbf{Indicateur} & \textbf{Quartier Bonabéri (Zone résidentielle dense)} & \textbf{Quartier Akwa (Centre d'affaires / mixte)} \\ \hline
Effectif total $N$ & 200 & 200 \\
Moyenne $\bar{x}$ & \textbf{397,5 kWh} & \textbf{427,5 kWh} \\
Médiane $M_e$ & 387,5 kWh & 420 kWh \\
Mode (centre classe modale) & 350 kWh & 350 kWh \\
Écart-type $\sigma$ & \textbf{107 kWh} & \textbf{148 kWh} \\
Variance $V$ & 11\,494 kWh$^2$ & 21\,774 kWh$^2$ \\ \hline
\end{tabularx}
\end{center}

\textbf{Analyse :} Les moyennes sont relativement proches ($\approx$ 400--430 kWh). Cependant, $\sigma_{\text{Akwa}} \approx 1,4 \times \sigma_{\text{Bonabéri}}$.
\begin{itemize}
    \item \textbf{Bonabéri :} Distribution resserrée (homogène). La majorité des foyers consomment entre 290 et 505 kWh ($\bar{x} \pm \sigma$). Une tarification unique (forfait) est pertinente.
    \item \textbf{Akwa :} Distribution très étalée (hétérogène). On trouve de petits appartements (150 kWh) et de gros commerces/hôtels (600+ kWh). La moyenne « 427,5 kWh » ne représente \textbf{personne} réellement. Une tarification par tranches (progressivité) est indispensable.
\end{itemize}
\end{exemplebox}

\begin{popfigure}
\centering
\begin{tikzpicture}
\begin{groupplot}[
    group style={group size=2 by 1, horizontal sep=1.5cm},
    width=0.45\linewidth, height=6cm,
    ybar, /pgf/bar width=18pt,
    xlabel={Conso. (kWh)}, ylabel={Effectif},
    xtick={150,250,350,450,550,650}, 
    xticklabels={150,250,350,450,550,650},
    ymin=0, ymax=90,
    title style={font=\bfseries\small},
    axis lines=left, enlarge x limits=0.08,
    x tick label style={font=\footnotesize, rotate=45, anchor=east},
    y tick label style={font=\footnotesize}
]
\nextgroupplot[title=Bonabéri ($\sigma\approx 107$), fill=popblue!70]
\addplot coordinates {
    (150,5) (250,25) (350,80) (450,60) (550,20) (650,10)
};
\nextgroupplot[title=Akwa ($\sigma\approx 148$), fill=poporange!70]
\addplot coordinates {
    (150,12) (250,30) (350,50) (450,40) (550,35) (650,33)
};
\end{groupplot}
\end{tikzpicture}
\legende{Histogrammes comparés (mêmes classes, amplitude 100 kWh) : même centre (moyenne $\approx$ 400), dispersion radicalement différente. Bonabéri (bleu) : pic étroit, prévisible. Akwa (orange) : large, étalé, queue à droite marquée.}
\end{popfigure}

\courssub{Cohérence entre mode, moyenne et médiane : symétrie et étalement}

La position relative de ces trois indicateurs révèle la \textbf{forme} de la distribution (symétrique ou asymétrique) sans même tracer l'histogramme.

\begin{propriete}[Règle empirique de position relative (Pearson)]
Pour une distribution unimodale (un seul pic) :
\begin{itemize}
    \item \textbf{Symétrique :} $\text{Mode} \approx \text{Médiane} \approx \text{Moyenne}$. L'histogramme est symétrique par rapport au centre.
    \item \textbf{Asymétrique à droite (queue à droite, valeurs fortes) :} $\text{Mode} < \text{Médiane} < \text{Moyenne}$. La moyenne est « tirée » vers les valeurs extrêmes fortes.
    \item \textbf{Asymétrique à gauche (queue à gauche, valeurs faibles) :} $\text{Moyenne} < \text{Médiane} < \text{Mode}$. La moyenne est « tirée » vers les valeurs extrêmes faibles.
\end{itemize}
\end{propriete}

\begin{attention}[Piège : Moyenne vs Médiane pour les salaires / revenus / consommations]
Dans les séries de revenus, salaires ou consommations (souvent étalées à droite), \textbf{la moyenne surestime le niveau « typique »}. La médiane (ou le mode) est plus représentative de la situation du « ménage moyen » ou de l'abonné type. Au Probatoire, si on te demande « Quel paramètre choisir pour représenter le niveau de vie ? », réponds \textbf{la médiane} (ou le mode) et justifie : « Car la série est asymétrique à droite, la moyenne est influencée par les valeurs extrêmes. »
\end{attention}

\courssub{Rédiger une conclusion argumentée à destination d'un décideur (cas ENEO)}

Une étude statistique ne s'arrête pas aux chiffres. Elle doit produire une \textbf{recommandation actionnable}. Voici la structure type d'une conclusion professionnelle.

\begin{methode}[Rédiger l'interprétation d'un paramètre en une phrase contextualisée]
Suis le patron : \textbf{[Paramètre] + [Valeur + Unité] + [Sens contextuel] + [Conséquence/Recommandation]}.

\begin{tabularx}{\linewidth}{|l|X|}\hline
\textbf{Paramètre} & \textbf{Exemple de phrase type} \\ \hline
Moyenne & « La consommation moyenne s'élève à \textbf{397,5 kWh/mois}, ce qui représente une facture moyenne d'environ \textbf{32 000 FCFA} (au tarif actuel) pour un abonné de Bonabéri. » \\ \hline
Médiane & « La médiane (387,5 kWh) indique que \textbf{50\% des abonnés consomment moins de 387,5 kWh} ; un tarif social ciblant les « petits consommateurs » doit donc se situer sous ce seuil. » \\ \hline
Mode & « Le mode (350 kWh) correspond à la consommation la plus fréquente ; c'est le point de référence idéal pour calibrer le \textbf{forfait de base} du nouveau tarif. » \\ \hline
Écart-type & « L'écart-type de \textbf{148 kWh à Akwa} révèle une \textbf{très forte disparité} entre abonnés : un tarif unique serait inéquitable, il faut impérativement \textbf{des tranches progressives}. » \\ \hline
\end{tabularx}
\end{methode}

\begin{attention}[Erreurs fatales à l'examen (Probatoire) et dans la vie professionnelle]
\begin{enumerate}
    \item \textbf{Oublier l'unité :} Écrire « La moyenne est 397,5 » au lieu de « 397,5 kWh ». \textbf{0 point} sur l'interprétation.
    \item \textbf{Donner un chiffre sans phrase :} Recopier le tableau de résultats sans aucune phrase de commentaire.
    \item \textbf{Confondre variance et écart-type :} Dire « L'écart-type est 21 774 » (c'est la variance). L'écart-type est $\sqrt{21774} \approx 148$ kWh.
    \item \textbf{Interpréter la moyenne sur une série asymétrique :} Dire « L'abonné type consomme 427,5 kWh » à Akwa (alors que la médiane est 420 et le mode 350).
\end{enumerate}
\end{attention}

\courssub{Limites du regroupement en classes : perte d'information et approximations (ouverture Terminale)}

Regrouper en classes est une \textbf{nécessité pratique} (lisibilité, volume de données), mais c'est un \textbf{choix arbitraire} qui dégrade l'information.

\begin{definition}[Perte d'information par regroupement]
Remplacer chaque valeur brute $x$ par le centre $x_i$ de sa classe commet une erreur maximale de $\pm h/2$ ($h$ = amplitude).
\begin{itemize}
    \item La \textbf{moyenne} calculée sur classes est une \textbf{approximation} de la vraie moyenne (erreur $\le h/2$).
    \item La \textbf{variance} sur classes \textbf{surestime} systématiquement la vraie variance (formule de Sheppard : $V_{\text{vraie}} \approx V_{\text{classes}} - h^2/12$).
    \item La \textbf{médiane} et le \textbf{mode} par interpolation supposent une \textbf{répartition uniforme} à l'intérieur de la classe (hypothèse souvent fausse).
\end{itemize}
\end{definition}

\begin{savaistu}[L'INS Cameroun et les enquêtes ECAM (Enquête Camerounaise Auprès des Ménages)]
L'Institut National de la Statistique (INS) utilise exactement ces outils pour mesurer la pauvreté, l'inflation, le chômage.
\begin{itemize}
    \item Les dépenses de consommation sont regroupées en classes (ex: $[0; 50\,000[$, $[50\,000; 100\,000[$, ... FCFA/mois).
    \item La \textbf{médiane} des dépenses par adulte-équivalent sert de \textbf{seuil de pauvreté} (ligne de pauvreté).
    \item L'\textbf{écart-type} (ou l'indice de Gini, vu en Terminale) mesure les \textbf{inégalités} entre régions (Extrême-Nord vs Littoral).
    \item Les histogrammes des dépenses guident la politique de \textbf{ciblage des filets sociaux} (transferts monétaires).
\end{itemize}
Maîtriser cette leçon, c'est déjà faire du travail d'analyste statistique à l'INS ou à la BEAC !
\end{savaistu}

\begin{exoresolu}[Étude complète et recommandation tarifaire : Quartiers Bonabéri vs Akwa]\label{exo:complete-bonaberi-akwa}
\textbf{Énoncé :} À partir des données du tableau de l'exemple \ref{ex:bonaberi-akwa}, réaliser l'étude statistique complète des deux séries et rédiger la note de synthèse à l'attention du Directeur Commercial d'ENEO pour la révision des tarifs 2025.

\textbf{Solution détaillée :}

\noindent \textbf{1. Paramètres de position}
\begin{itemize}
    \item \textbf{Moyennes :} 
    $\bar{x}_B = \frac{1}{200}(5\times150 + 25\times250 + 80\times350 + 60\times450 + 20\times550 + 10\times650) = \frac{79\,500}{200} = \textbf{397,5 kWh}$.
    $\bar{x}_A = \frac{1}{200}(12\times150 + 30\times250 + 50\times350 + 40\times450 + 35\times550 + 33\times650) = \frac{85\,500}{200} = \textbf{427,5 kWh}$.
    \item \textbf{Classes modales :} 
    Bonabéri : $[300;400[$ ($n=80$) $\to$ Mode $= 350$ kWh. 
    Akwa : $[300;400[$ ($n=50$) $\to$ Mode $= 350$ kWh.
    \item \textbf{Médianes :} $N/2 = 100$.
    \begin{itemize}
        \item Bonabéri : Effectifs cumulés croissants : $5 ; 30 ; \mathbf{110} ; \dots$ Classe médiane $=[300;400[$. $N_{\text{ant}}=30, n_{\text{med}}=80, h=100$.
        $M_{e,B} = 300 + \frac{100-30}{80}\times 100 = 300 + 87,5 = \textbf{387,5 kWh}$.
        \item Akwa : Effectifs cumulés croissants : $12 ; 42 ; 92 ; \mathbf{132} ; \dots$ Classe médiane $=[400;500[$. $N_{\text{ant}}=92, n_{\text{med}}=40, h=100$.
        $M_{e,A} = 400 + \frac{100-92}{40}\times 100 = 400 + 20 = \textbf{420 kWh}$.
    \end{itemize}
\end{itemize}

\noindent \textbf{2. Paramètres de dispersion (Calcul de $\sum n_i x_i^2$ nécessaire)}
\begin{align*}
\text{Bonabéri : } \sum n_i x_i^2 &= 5(150^2) + 25(250^2) + 80(350^2) + 60(450^2) + 20(550^2) + 10(650^2) \\
&= 112\,500 + 1\,562\,500 + 9\,800\,000 + 12\,150\,000 + 6\,050\,000 + 4\,225\,000 = \textbf{33\,900\,000} \\
V_B &= \frac{33\,900\,000}{200} - 397,5^2 = 169\,500 - 158\,006,25 = \textbf{11\,493,75 kWh}^2 \\
\sigma_B &= \sqrt{11\,493,75} \approx \textbf{107,2 kWh} \\[1em]
\text{Akwa : } \sum n_i x_i^2 &= 12(150^2) + 30(250^2) + 50(350^2) + 40(450^2) + 35(550^2) + 33(650^2) \\
&= 270\,000 + 1\,875\,000 + 6\,125\,000 + 8\,100\,000 + 10\,587\,500 + \mathbf{13\,942\,500} = \textbf{40\,900\,000} \\
V_A &= \frac{40\,900\,000}{200} - 427,5^2 = 204\,500 - 182\,756,25 = \textbf{21\,743,75 kWh}^2 \\
\sigma_A &= \sqrt{21\,743,75} \approx \textbf{147,5 kWh}
\end{align*}

\noindent \textbf{3. Synthèse comparative}
\begin{center}
\begin{tabular}{|l|c|c|}\hline
& \textbf{Bonabéri} & \textbf{Akwa} \\ \hline
Moyenne & 397,5 kWh & 427,5 kWh \\
Médiane & 387,5 kWh & 420 kWh \\
Mode & 350 kWh & 350 kWh \\
Écart-type $\sigma$ & \textbf{107 kWh} & \textbf{148 kWh} \\
Coeff. de variation ($\sigma/\bar{x}$) & 27\% & 35\% \\ \hline
\end{tabular}
\end{center}

\noindent \textbf{4. Interprétation et Recommandation (Note au Directeur Commercial)}

\begin{quote}
\textbf{Objet : Proposition de structure tarifaire différenciée Bonabéri / Akwa (Échantillon 400 abonnés).}

Monsieur le Directeur,

L'analyse statistique révèle deux profils de consommation radicalement différents, bien que les moyennes soient de même ordre (398 vs 428 kWh).

\textbf{1. Bonabéri (Zone résidentielle homogène) :} La distribution est relativement concentrée autour de 390--400 kWh ($\sigma = 107$ kWh, CV=27\%). La médiane (387,5) et la moyenne (397,5) sont proches, la série est quasi-symétrique (légère asymétrie droite). \textbf{Recommandation :} Maintenir un \textbf{tarif forfaitaire unique} (ou deux tranches max) calibré sur 400 kWh. La prévisibilité de la recette est forte.

\textbf{2. Akwa (Zone mixte résidentiel/commercial) :} La dispersion est très forte ($\sigma = 148$ kWh, CV=35\%). La moyenne (427,5) est supérieure à la médiane (420), confirmant une asymétrie à droite (présence de gros consommateurs > 600 kWh). Le mode (350 kWh) montre que la majorité des \emph{ménages} consomment peu, mais les \emph{commerces} tirent la moyenne vers le haut. \textbf{Recommandation :} Instaurer une \textbf{tarification progressive à 4 tranches} :
\begin{itemize}
    \item Tranche sociale : $[0; 300[$ kWh (tarif bas, couvre le mode 350 et protège les ménages modestes).
    \item Tranche standard : $[300; 500[$ kWh (tarif courant, couvre la médiane 420).
    \item Tranche confort : $[500; 700[$ kWh.
    \item Tranche premium : $\ge 700$ kWh (tarif élevé, capture la surconsommation commerciale).
\end{itemize}

Cette structure assure l'équité (progressivité) et optimise le revenu (capture de la rente des gros consommateurs d'Akwa) tout en protégeant le pouvoir d'achat à Bonabéri.

Cordialement, \textit{Le Service Études Statistiques.}
\end{quote}
\end{exoresolu}

\courssub{Exercices d'entraînement (Probatoire type)}

\begin{exercice}[Synthèse : Salaires dans deux entreprises]
L'entreprise \textbf{CAMTEL-Yaoundé} (200 employés) et \textbf{CAMTEL-Douala} (200 employés) présentent les séries de salaires mensuels (en milliers de FCFA) regroupées ci-dessous. Amplitude $h=50$.

\begin{center}
\begin{tabular}{|c|c|c|c|}\hline
\textbf{Classe (kFCFA)} & \textbf{Centre} & \textbf{Yaoundé $n_i$} & \textbf{Douala $n_i$} \\ \hline
$[150;200[$ & 175 & 10 & 5 \\ \hline
$[200;250[$ & 225 & 30 & 15 \\ \hline
$[250;300[$ & 275 & 70 & 50 \\ \hline
$[300;350[$ & 325 & 60 & 60 \\ \hline
$[350;400[$ & 375 & 20 & 40 \\ \hline
$[400;450[$ & 425 & 10 & 30 \\ \hline
\end{tabular}
\end{center}

\begin{enumerate}
    \item Calculer la moyenne, la médiane (par interpolation), le mode et l'écart-type pour les deux séries.
    \item Comparer les deux distributions (position, dispersion, forme). Laquelle est la plus équitable ?
    \item Le DRH veut instaurer une prime de rendement unique pour les « salaires moyens ». Quel paramètre (moyenne ou médiane) choisir pour définir le seuil d'éligibilité ? Justifier.
    \item Rédiger la conclusion à l'attention de la Direction Générale (3 phrases max).
\end{enumerate}
\end{exercice}

\begin{corrige}[Exercice Synthèse Camtel]
\textbf{1. Calculs (Yaoundé / Douala) :}
\begin{itemize}
    \item \textbf{Moyennes :} 
    $\bar{x}_Y = \frac{10\cdot175 + 30\cdot225 + 70\cdot275 + 60\cdot325 + 20\cdot375 + 10\cdot425}{200} = \frac{59\,000}{200} = \textbf{295 kFCFA}$.
    $\bar{x}_D = \frac{5\cdot175 + 15\cdot225 + 50\cdot275 + 60\cdot325 + 40\cdot375 + 30\cdot425}{200} = \frac{65\,250}{200} = \textbf{326,25 kFCFA}$.
    \item \textbf{Modes :} 
    Yaoundé : classe $[250;300[$ ($n=70$) $\to$ Mode $= 275$ kFCFA. 
    Douala : classe $[300;350[$ ($n=60$) $\to$ Mode $= 325$ kFCFA.
    \item \textbf{Médianes :} $N/2=100$.
    Yaoundé : Cumulés $10; 40; \mathbf{110}$. Classe médiane $[250;300[$. $M_{e,Y} = 250 + \frac{100-40}{70}\times 50 = 250 + 42,9 = \textbf{292,9 kFCFA}$.
    Douala : Cumulés $5; 20; 70; \mathbf{130}$. Classe médiane $[300;350[$. $M_{e,D} = 300 + \frac{100-70}{60}\times 50 = 300 + 25 = \textbf{325 kFCFA}$.
    \item \textbf{Écart-types :} (Calcul de $\sum n_i x_i^2$ via König-Huygens)
    \begin{align*}
    \text{Yaoundé : } \sum n_i x_i^2 &= 10(175^2)+30(225^2)+70(275^2)+60(325^2)+20(375^2)+10(425^2) = 18\,075\,000 \\
    V_Y &= \frac{18\,075\,000}{200} - 295^2 = 90\,375 - 87\,025 = \textbf{3\,350 kFCFA}^2 \\
    \sigma_Y &= \sqrt{3\,350} \approx \textbf{57,9 kFCFA} \\[1em]
    \text{Douala : } \sum n_i x_i^2 &= 5(175^2)+15(225^2)+50(275^2)+60(325^2)+40(375^2)+30(425^2) = 22\,075\,000 \\
    V_D &= \frac{22\,075\,000}{200} - 326,25^2 = 110\,375 - 106\,439,06 = \textbf{3\,935,94 kFCFA}^2 \\
    \sigma_D &= \sqrt{3\,935,94} \approx \textbf{62,7 kFCFA}
    \end{align*}
\end{itemize}

\textbf{2. Comparaison :}
\begin{itemize}
    \item \textbf{Position :} Douala paie mieux en moyenne (+31,25 kFCFA, soit +10,6\%). Médiane (325 vs 292,9) et Mode (325 vs 275) aussi plus élevés.
    \item \textbf{Dispersion :} Douala plus dispersé en valeur absolue ($\sigma_D > \sigma_Y$), mais coefficient de variation $CV_D = 19,2\%$ vs $CV_Y = 19,6\%$ (homogénéité relative très similaire).
    \item \textbf{Forme :} 
    Yaoundé : Mode (275) $<$ Médiane (292,9) $<$ Moyenne (295) $\to$ \textbf{Asymétrique à droite} (queue vers les hauts salaires, bien que modérée).
    Douala : Mode (325) $=$ Médiane (325) $\approx$ Moyenne (326,25) $\to$ \textbf{Quasi-symétrique}.
    \item \textbf{Équité :} Yaoundé présente une queue vers les hauts salaires (quelques cadres très bien payés). Douala a une structure plus symétrique mais plus étalée en valeur absolue.
\end{itemize}

\textbf{3. Choix du seuil :} \textbf{La médiane}. À Yaoundé, la moyenne (295) est légèrement tirée vers le haut par les hauts salaires (queue droite) ; la médiane (292,9) représente mieux le « salaire central » du salarié type. À Douala, les deux sont quasi égaux. La médiane est robuste aux valeurs extrêmes, c'est le critère de référence pour le « salaire médian ».

\textbf{4. Conclusion DG :} « Douala présente un niveau salarial supérieur de 11\% (médiane 325 vs 293 kFCFA) avec une structure symétrique. Yaoundé montre une légère asymétrie à droite (moyenne $>$ médiane) due à quelques hauts salaires. Nous recommandons une prime unique indexée sur la médiane (325 kFCFA) pour Douala, et une vigilance sur l'équité des grilles à Yaoundé où l'écart type relatif est équivalent mais la dispersion absolue favorise les écarts. »
\end{corrige}

\newpage

\coursec{Activité d'intégration}

\coursec{Leçon 9 : Statistiques — Séries statistiques regroupées en classes}

\courssub{Rappels de cours : notions et formules}

\begin{definition}[Série statistique regroupée en classes]
Lorsque l'effectif total $N$ est grand, on regroupe les valeurs d'une variable quantitative continue en \cle{classes} (intervalles) de la forme $[a\,;b[$. Chaque classe $i$ est caractérisée par :
\begin{itemize}
\item ses bornes $a_i$ (inférieure) et $b_i$ (supérieure) ;
\item son \cle{centre} (ou valeur représentative) $x_i = \dfrac{a_i+b_i}{2}$ ;
\item son \cle{amplitude} (ou largeur) $l_i = b_i - a_i$ ;
\item son \cle{effectif} $n_i$ (nombre d'individus dans la classe) ;
\item sa \cle{fréquence} $f_i = \dfrac{n_i}{N}$ (souvent en \%) ;
\item ses \cle{effectifs cumulés croissants} $ECC_i = \sum_{k\le i} n_k$ et \cle{fréquences cumulées croissantes} $FCC_i = \sum_{k\le i} f_k$.
\end{itemize}
\end{definition}

\begin{propriete}[Paramètres de position]
Pour une série regroupée en classes $(x_i ; n_i)_{1\le i\le k}$ d'effectif total $N$ :
\begin{itemize}
\item \textbf{Moyenne :} $\displaystyle\bar{x} = \frac{1}{N}\sum_{i=1}^k n_i x_i$. \emph{(Exigible au Probatoire)}
\item \textbf{Classe modale :} classe de plus grande \emph{densité} $\dfrac{n_i}{l_i}$ (si amplitudes inégales) ou de plus grand effectif $n_i$ (si amplitudes égales).
\item \textbf{Mode :} centre $x_{i_0}$ de la classe modale (approximation usuelle au lycée).
\item \textbf{Médiane :} valeur $Me$ telle que $FCC(Me) \approx 50\,\%$. Si $[a\,;b[$ est la classe médiane (contenant le $N/2$-ième individu), l'\cle{interpolation linéaire} donne :
\[ Me = a + (b-a)\times\dfrac{\frac{N}{2} - ECC_{\text{avant}}}{n_{\text{médiane}}}. \]
\emph{(Démonstration de la formule non exigible, mais application rigoureuse exigée.)}
\end{itemize}
\end{propriete}

\begin{propriete}[Paramètres de dispersion]
\begin{itemize}
\item \textbf{Variance (formule de König-Huygens) :} $\displaystyle V = \frac{1}{N}\sum_{i=1}^k n_i x_i^2 - \bar{x}^2$.
\item \textbf{Écart-type :} $\sigma = \sqrt{V}$ (même unité que la variable).
\item \textbf{Écart moyen :} $\displaystyle E = \frac{1}{N}\sum_{i=1}^k n_i |x_i - \bar{x}|$ (moins utilisé, mais au programme).
\end{itemize}
\emph{Interprétation :} $\sigma$ mesure la dispersion autour de la moyenne. Un $\sigma$ grand devant $\bar{x}$ signale une série hétérogène.
\end{propriete}

\begin{methode}[Construire l'histogramme (amplitudes inégales)]
\begin{enumerate}
\item Calculer la \cle{densité} $d_i = \dfrac{n_i}{l_i}$ pour chaque classe.
\item Sur l'axe des abscisses, reporter les bornes des classes à l'échelle.
\item Sur l'axe des ordonnées, reporter les densités.
\item Tracer les rectangles : base $= l_i$, hauteur $= d_i$. \textbf{L'aire du rectangle est proportionnelle à l'effectif $n_i$.}
\end{enumerate}
\end{methode}

\begin{methode}[Tracer la courbe des fréquences cumulées croissantes]
\begin{enumerate}
\item Compléter le tableau des $FCC_i$ (en \%) aux bornes supérieures $b_i$. Ajouter l'origine $(a_1\,;0)$.
\item Repérer les points $(b_i\,; FCC_i)$ dans un repère orthonormé.
\item Relier ces points par des segments de droite (ligne brisée).
\item \textbf{Lecture de la médiane :} tracer l'horizontale $y=50$ ; l'abscisse de l'intersection avec la courbe donne $Me$ graphiquement.
\end{enumerate}
\end{methode}

\begin{attention}[Pièges classiques au Probatoire]
\begin{itemize}
\item \textbf{Histogramme :} ne jamais tracer les hauteurs proportionnelles aux effectifs si les amplitudes diffèrent. L'aire, pas la hauteur, représente l'effectif.
\item \textbf{Médiane :} dans la formule d'interpolation, $ECC_{\text{avant}}$ est l'effectif cumulé de la classe \emph{précédente} (strictement inférieure à la classe médiane), pas celui de la classe médiane.
\item \textbf{König-Huygens :} ne pas confondre $\bar{x}^2$ (carré de la moyenne) et $\frac{1}{N}\sum n_i x_i^2$ (moyenne des carrés). Calculer les deux séparément.
\item \textbf{Seuil et FCC :} pour estimer le pourcentage \emph{au-dessus} d'un seuil $s$, on lit $FCC(s)$ sur la courbe (ou par interpolation) et on calcule $100 - FCC(s)$.
\item \textbf{Mode et amplitudes inégales :} la classe modale est celle de plus grande \emph{densité} $n_i/l_i$, pas nécessairement celle de plus grand effectif.
\end{itemize}
\end{attention}

\courssub{Situation-problème : Consommation d'eau à Nkolbisson (CAMWATER)}

La \textbf{CAMWATER} (Cameroon Water Utilities Corporation) souhaite analyser la consommation d'eau potable du quartier \emph{Nkolbisson} à Yaoundé afin d'ajuster sa politique de facturation et de sensibilisation. Le service commercial a relevé, pour le mois de mars, les consommations mensuelles (en mètres cubes) de $120$ ménages abonnés. Les données, déjà regroupées en classes, sont consignées dans le tableau incomplet ci-dessous.

\begin{center}
\begin{tabularx}{\linewidth}{|l|X|X|X|X|X|}\hline
\rowcolor{popblueL} Classe (en m$^3$) & $[0\,;5[$ & $[5\,;10[$ & $[10\,;15[$ & $[15\,;20[$ & $[20\,;30[$ \\ \hline
Effectif $n_i$ & $18$ & $30$ & $36$ & $24$ & $12$ \\ \hline
Centre $x_i$ & & & & & \\ \hline
Fréquence $f_i$ (en \%) & & & & & \\ \hline
Effectif cumulé croissant & & & & & \\ \hline
Fréquence cumulée croissante (en \%) & & & & & \\ \hline
\end{tabularx}
\end{center}

Le directeur régional s'interroge : \emph{« Quelle est la consommation mensuelle typique d'un ménage du quartier ? Les consommations sont-elles homogènes ? À partir de quel seuil peut-on qualifier une consommation d'excessive, et quelle proportion de ménages serait concernée par une campagne de sensibilisation ciblée ? »}

\courssub{Résolution guidée et commentée du dossier}

\begin{exoresolu}[Étude complète de la consommation d'eau du quartier Nkolbisson]
\tcblower
\textbf{Étape 1 : Organisation des données — Tableau complété.}
L'effectif total est $N = 18+30+36+24+12 = 120$ : la donnée est cohérente.
On calcule les centres $x_i = \frac{a_i+b_i}{2}$, les fréquences $f_i = \frac{n_i}{120}\times 100$, et les cumuls.

\begin{center}
\begin{tabularx}{\linewidth}{|l|X|X|X|X|X|}\hline
\rowcolor{popblueL} Classe (m$^3$) & $[0;5[$ & $[5;10[$ & $[10;15[$ & $[15;20[$ & $[20;30[$ \\ \hline
Amplitude $l_i$ & $5$ & $5$ & $5$ & $5$ & $10$ \\ \hline
Effectif $n_i$ & $18$ & $30$ & $36$ & $24$ & $12$ \\ \hline
Centre $x_i$ & $2,5$ & $7,5$ & $12,5$ & $17,5$ & $25$ \\ \hline
Fréquence $f_i$ (\%) & $15$ & $25$ & $30$ & $20$ & $10$ \\ \hline
Densité $d_i=n_i/l_i$ & $3,6$ & $6$ & $7,2$ & $4,8$ & $1,2$ \\ \hline
ECC & $18$ & $48$ & $84$ & $108$ & $120$ \\ \hline
FCC (\%) & $15$ & $40$ & $70$ & $90$ & $100$ \\ \hline
\end{tabularx}
\end{center}

\textbf{Étape 2 : Paramètres de position — Moyenne.}
\begin{align*}
\bar{x} &= \frac{1}{120}\Bigl(18\times 2,5 + 30\times 7,5 + 36\times 12,5 + 24\times 17,5 + 12\times 25\Bigr) \\
&= \frac{45 + 225 + 450 + 420 + 300}{120} = \frac{1440}{120} = 12 \text{ m}^3.
\end{align*}
\emph{Interprétation :} en moyenne, un ménage du quartier consomme $12$ m$^3$ d'eau par mois.

\textbf{Étape 3 : Classe modale et Mode.}
Les amplitudes ne sont pas toutes égales (dernière classe de largeur $10$). On compare les \textbf{densités} $d_i = n_i/l_i$ :
$d_1=3,6$ ; $d_2=6$ ; $d_3=7,2$ ; $d_4=4,8$ ; $d_5=1,2$.
La plus grande densité est $7,2$ pour la classe $[10\,;15[$ : c'est la \cle{classe modale}.
Le \cle{mode} (valeur la plus fréquente) est approché par le centre de cette classe : $Mo = 12,5$ m$^3$.
\emph{Remarque :} ici, la classe de plus grand effectif ($36$) coïncide avec celle de plus grande densité, mais ce n'est pas toujours le cas.

\textbf{Étape 4 : Médiane par interpolation linéaire.}
On cherche la classe contenant le $\frac{N}{2} = 60^{\text{ième}}$ individu (dans l'ordre croissant).
D'après les ECC : $ECC_2 = 48 < 60 \le 84 = ECC_3$.
La \cle{classe médiane} est donc $[10\,;15[$ ($a=10$, $b=15$, $n_{\text{med}}=36$, $ECC_{\text{avant}}=48$).
\[ Me = 10 + (15-10)\times\frac{60 - 48}{36} = 10 + 5\times\frac{12}{36} = 10 + \frac{5}{3} \approx 11,67 \text{ m}^3. \]
\emph{Interprétation :} environ $50\,\%$ des ménages consomment moins de $11,67$ m$^3$/mois, et $50\,\%$ consomment davantage. La médiane ($11,67$) est proche de la moyenne ($12$), ce qui suggère une distribution assez symétrique autour du centre, mais l'écart-type nous en dira plus.

\textbf{Étape 5 : Paramètres de dispersion — Variance et écart-type (König-Huygens).}
Calcul de la moyenne des carrés des centres :
\begin{align*}
\frac{1}{N}\sum n_i x_i^2 &= \frac{1}{120}\Bigl(18\times 2,5^2 + 30\times 7,5^2 + 36\times 12,5^2 + 24\times 17,5^2 + 12\times 25^2\Bigr) \\
&= \frac{1}{120}\Bigl(18\times 6,25 + 30\times 56,25 + 36\times 156,25 + 24\times 306,25 + 12\times 625\Bigr) \\
&= \frac{112,5 + 1687,5 + 5625 + 7350 + 7500}{120} = \frac{22275}{120} = 185,625.
\end{align*}
Variance :
\[ V = 185,625 - \bar{x}^2 = 185,625 - 144 = 41,625 \text{ m}^6. \]
Écart-type :
\[ \sigma = \sqrt{41,625} \approx 6,45 \text{ m}^3 \quad (\text{arrondi au centième}). \]
\emph{Interprétation :} $\sigma \approx 6,45$ représente plus de la moitié de la moyenne ($12$). Les consommations sont \textbf{fortement dispersées} : le quartier est hétérogène, il y a de gros écarts entre ménages économes et gros consommateurs.

\textbf{Étape 6 : Représentation graphique — Histogramme.}
On utilise les densités $d_i$ calculées à l'étape 1. Le dernier rectangle (classe $[20;30[$ d'amplitude $10$) est deux fois plus large que les autres ; sa hauteur ($1,2$) est divisée par deux par rapport à ce qu'elle serait si on utilisait les effectifs bruts.

\begin{popfigure}
\begin{tikzpicture}[x=0.55cm,y=0.8cm]
\draw[->,thick] (0,0) -- (31,0) node[below right]{\footnotesize Consommation (m$^3$)};
\draw[->,thick] (0,0) -- (0,8.5) node[above left]{\footnotesize Densité d'effectif};
\foreach \h/\a/\b in {3.6/0/5, 6/5/10, 7.2/10/15, 4.8/15/20, 1.2/20/30}
  \draw[fill=popblue!40,draw=popblue,thick] (\a,0) rectangle (\b,\h);
\foreach \x in {0,5,10,15,20,30} \draw (\x,0.1) -- (\x,-0.1) node[below]{\footnotesize \x};
\foreach \y in {2,4,6,8} \draw (0.15,\y) -- (-0.15,\y) node[left]{\footnotesize \y};
\node[popblue,above] at (2.5,3.6) {\footnotesize 3,6};
\node[popblue,above] at (7.5,6) {\footnotesize 6};
\node[popblue,above] at (12.5,7.2) {\footnotesize 7,2};
\node[popblue,above] at (17.5,4.8) {\footnotesize 4,8};
\node[popblue,above] at (25,1.2) {\footnotesize 1,2};
\end{tikzpicture}
\legende{Histogramme des consommations : les aires des rectangles sont proportionnelles aux effectifs. La classe modale $[10;15[$ a le rectangle de plus grande hauteur (densité max).}
\end{popfigure}

\textbf{Étape 7 : Courbe des fréquences cumulées croissantes et lecture graphique de la médiane.}
On trace les points $(b_i\,; FCC_i)$ : $(0;0)$, $(5;15)$, $(10;40)$, $(15;70)$, $(20;90)$, $(30;100)$.

\begin{popfigure}
\begin{tikzpicture}[x=0.42cm,y=0.075cm]
\draw[->,thick] (0,0) -- (32,0) node[below right]{\footnotesize Consommation (m$^3$)};
\draw[->,thick] (0,0) -- (0,108) node[above left]{\footnotesize FCC (\%)};
\draw[very thick,poporange] (0,0) -- (5,15) -- (10,40) -- (15,70) -- (20,90) -- (30,100);
\foreach \p in {(0,0),(5,15),(10,40),(15,70),(20,90),(30,100)} \fill[poporange] \p circle (0.18);
% Lecture graphique médiane
\draw[dashed,popgreen,thick] (0,50) -- (11.67,50) -- (11.67,0);
\node[popgreen,below] at (11.67,0) {\footnotesize $Me \approx 11,7$};
\node[popgreen,left] at (0,50) {\footnotesize 50\%};
\foreach \x in {5,10,15,20,30} \node[below] at (\x,0) {\footnotesize \x};
\foreach \y in {50,100} \node[left] at (0,\y) {\footnotesize \y};
\end{tikzpicture}
\legende{Courbe des FCC : lecture graphique de la médiane (abscisse du point d'ordonnée 50\%). On trouve $Me \approx 11,7$ m$^3$, confirmant le calcul $11,67$ m$^3$.}
\end{popfigure}

\textbf{Étape 8 : Seuil de consommation excessive et estimation de la population concernée.}
On définit le seuil $s = \bar{x} + \sigma = 12 + 6,45 = 18,45$ m$^3$.
Ce seuil se situe dans la classe $[15\,;20[$. On utilise la courbe des FCC (ou l'interpolation linéaire sur cette classe) pour estimer $FCC(18,45)$.
La courbe passe par $(15\,;70)$ et $(20\,;90)$ sur cet intervalle.
\[ FCC(18,45) \approx 70 + \frac{18,45 - 15}{20 - 15} \times (90 - 70) = 70 + \frac{3,45}{5} \times 20 = 70 + 13,8 = 83,8\,\%. \]
Le pourcentage de ménages \emph{dépassant} ce seuil est $100 - 83,8 = 16,2\,\%$.
Nombre approximatif de ménages : $0,162 \times 120 \approx 19,4$, soit \textbf{environ 19 ménages}.

\textbf{Étape 9 : Note de synthèse pour le Directeur régional de la CAMWATER.}

\emph{« Monsieur le Directeur, l'étude sur 120 ménages de Nkolbisson révèle une consommation moyenne de $12$ m$^3$/mois (médiane $11,67$ m$^3$, mode $12,5$ m$^3$). L'écart-type élevé ($6,45$ m$^3$) signale de fortes disparités. En fixant le seuil de consommation excessive à $\bar{x}+\sigma = 18,45$ m$^3$, environ $16\,\%$ des abonnés (19 ménages) sont concernés. Nous recommandons une campagne de sensibilisation ciblée sur ces gros consommateurs et le maintien d'un tarif social protecteur pour la tranche $[0;15[$ qui regroupe $70\,\%$ des ménages. »}
\end{exoresolu}

\courssub{Exercices d'application}

\begin{exercice}[Entraînement 1 : Série à amplitudes égales]
Une enquête sur le temps de trajet domicile-école (en minutes) de $200$ élèves d'un lycée de Douala donne le tableau suivant :
\begin{center}
\begin{tabular}{|c|c|c|c|c|c|}\hline
\rowcolor{popblueL} Temps (min) & $[0;10[$ & $[10;20[$ & $[20;30[$ & $[30;40[$ & $[40;50[$ \\ \hline
Effectif & $20$ & $50$ & $70$ & $40$ & $20$ \\ \hline
\end{tabular}
\end{center}
\begin{enumerate}
\item Compléter le tableau (centres, fréquences \%, ECC, FCC).
\item Calculer la moyenne, la médiane (interpolation), le mode.
\item Calculer la variance et l'écart-type. Les temps de trajet sont-ils homogènes ?
\item Construire l'histogramme et la courbe des FCC. Lire la médiane graphiquement.
\end{enumerate}
\end{exercice}

\begin{corrige}[Exercice 1]
\begin{enumerate}
\item Centres : $5; 15; 25; 35; 45$. Fréquences : $10\%; 25\%; 35\%; 20\%; 10\%$. ECC : $20; 70; 140; 180; 200$. FCC : $10; 35; 70; 90; 100$.
\item $\bar{x} = \frac{1}{200}(20\times5 + 50\times15 + 70\times25 + 40\times35 + 20\times45) = \frac{5000}{200} = 25$ min.
$N/2=100$. Classe médiane $[20;30[$ ($ECC_{\text{avant}}=70, n=70$). $Me = 20 + 10\times\frac{100-70}{70} \approx 24,29$ min.
Classe modale $[20;30[$ (effectif max $70$), Mode $= 25$ min.
\item $\frac{1}{N}\sum n_i x_i^2 = \frac{1}{200}(500 + 11250 + 43750 + 49000 + 40500) = \frac{145000}{200} = 725$.
$V = 725 - 25^2 = 100$. $\sigma = 10$ min. $\sigma/\bar{x} = 40\,\%$ : dispersion modérée.
\item Histogramme : amplitudes égales ($10$), hauteurs = effectifs (ou densités = effectifs/10). Courbe FCC : points $(10;10), (20;35), (30;70), (40;90), (50;100)$. Lecture graphique $Me \approx 24,3$ min.
\end{enumerate}
\end{corrige}

\begin{exercice}[Entraînement 2 : Interprétation et décision — Série à amplitudes inégales]
Une entreprise de téléphonie mobile (type MTN/Orange) analyse la durée mensuelle des appels (en minutes) de $500$ clients prépayés :
\begin{center}
\begin{tabular}{|c|c|c|c|c|c|}\hline
\rowcolor{popblueL} Durée (min) & $[0;30[$ & $[30;60[$ & $[60;100[$ & $[100;200[$ & $[200;300[$ \\ \hline
Effectif & $100$ & $150$ & $150$ & $70$ & $30$ \\ \hline
\end{tabular}
\end{center}
\begin{enumerate}
\item Compléter le tableau (centres, fréquences, densités, cumuls).
\item Calculer la moyenne et la médiane. Comparer.
\item Déterminer la classe modale (attention aux amplitudes) et le mode.
\item Calculer l'écart-type. Que signifie sa valeur par rapport à la moyenne ?
\item L'entreprise veut offrir un bonus aux $10\,\%$ des clients les plus actifs (plus gros consommateurs). Déterminer le seuil minimal de durée d'appel pour bénéficier du bonus (utiliser la courbe des FCC ou l'interpolation).
\end{enumerate}
\end{exercice}

\begin{corrige}[Exercice 2]
\begin{enumerate}
\item Centres : $15; 45; 80; 150; 250$. Amplitudes : $30; 30; 40; 100; 100$.
Densités : $100/30\approx3,33$; $150/30=5$; $150/40=3,75$; $70/100=0,7$; $30/100=0,3$.
Fréquences : $20\%; 30\%; 30\%; 14\%; 6\%$.
ECC : $100; 250; 400; 470; 500$. FCC : $20; 50; 80; 94; 100$.
\item $\bar{x} = \frac{1}{500}(1500+6750+12000+10500+7500) = \frac{38250}{500} = 76,5$ min.
$N/2=250$. $ECC_2=250$. La médiane est à la borne entre classe 2 et 3. Par convention, classe médiane $[60;100[$ ($ECC_{\text{avant}}=250$).
$Me = 60 + 40\times\frac{250-250}{150} = 60$ min. (Médiane = borne inférieure de la classe médiane car $N/2$ tombe exactement sur un ECC).
Moyenne ($76,5$) $>$ Médiane ($60$) : distribution asymétrique à droite (queue de gros consommateurs).
\item Classe modale = classe de plus grande \textbf{densité} : $[30;60[$ (densité $5$). Mode $= 45$ min.
(Note : effectif max $150$ partagé par classes 2 et 3, mais densité plus forte en classe 2).
\item $\frac{1}{N}\sum n_i x_i^2 = \frac{1}{500}(22500+303750+960000+1575000+1875000) = \frac{4736250}{500} = 9472,5$.
$V = 9472,5 - 76,5^2 = 9472,5 - 5852,25 = 3620,25$. $\sigma = \sqrt{3620,25} = 60,17$ min.
$\sigma \approx 60$ min est proche de la moyenne $76,5$ : très forte dispersion.
\item Top $10\,\%$ $\Leftrightarrow$ $90^{\text{ième}}$ percentile. $FCC = 90\,\%$. Classe $[100;200[$ ($FCC$ passe de $80$ à $94$).
Interpolation : $x = 100 + 100\times\frac{90-80}{94-80} = 100 + 100\times\frac{10}{14} \approx 171,4$ min.
Seuil minimal $\approx 171$ min.
\end{enumerate}
\end{corrige}

\courssub{Bilan de la leçon et compétences validées}

\begin{aretenir}[Ce qu'il faut retenir pour le Probatoire]
\begin{itemize}
\item \textbf{Tableau :} Centres, Fréquences (\%), ECC, FCC, \textbf{Densités} (si amplitudes inégales).
\item \textbf{Position :} $\bar{x} = \frac{\sum n_i x_i}{N}$ ; Mode = centre classe modale (densité max) ; Médiane par interpolation linéaire dans la classe contenant $N/2$.
\item \textbf{Dispersion :} $V = \frac{\sum n_i x_i^2}{N} - \bar{x}^2$ (König-Huygens) ; $\sigma = \sqrt{V}$ ; $E = \frac{\sum n_i |x_i-\bar{x}|}{N}$.
\item \textbf{Graphiques :} Histogramme (hauteur = densité, aire $\propto$ effectif) ; Courbe FCC (points $(b_i; FCC_i)$), lecture médiane ($y=50$) et percentiles.
\item \textbf{Interprétation :} Comparer $\bar{x}$ et $Me$ (symétrie/asymétrie) ; comparer $\sigma$ à $\bar{x}$ (homogénéité/hétérogénéité) ; utiliser FCC pour prise de décision (seuils, pourcentages).
\end{itemize}
\end{aretenir}

\begin{savaistu}[Le saviez-vous ?]
La méthode de l'interpolation linéaire pour la médiane suppose que les données sont \textbf{uniformément réparties} à l'intérieur de la classe médiane. C'est une approximation commode. En statistique professionnelle (INS, Banque Mondiale), on utilise parfois des noyaux de lissage (kernel density estimation) pour estimer la fonction de répartition de façon plus fine, mais l'interpolation linéaire reste la référence pour les examens et les tableaux de bord opérationnels comme ceux de la CAMWATER ou de l'ARSEL (Agence de Régulation du Secteur de l'Électricité).
\end{savaistu}

\newpage

\coursec{Méthodes et exercices résolus}

\coursec{Statistiques : séries regroupées en classes}

\courssub{Regroupement en classes, effectifs et fréquences cumulés}

\begin{definition}[Série statistique regroupée en classes]
Soit une variable quantitative continue ou discrète à grande dispersion. On regroupe les valeurs en \textbf{classes} (intervalles) $I_1, I_2, \dots, I_k$ qui forment une partition de l'étendue de la série (classes mutuellement exclusives et exhaustives). Pour chaque classe $I_i = [a_i; a_{i+1}[$ (la dernière $[a_k; a_{k+1}]$ est fermée), on note :
\begin{itemize}
    \item \textbf{Centre de classe} (ou milieu) : $x_i = \dfrac{a_i + a_{i+1}}{2}$.
    \item \textbf{Effectif} $n_i$ : nombre d'observations dans $I_i$.
    \item \textbf{Fréquence} $f_i = \dfrac{n_i}{N}$ (souvent en $\%$), avec $N = \sum n_i$ effectif total.
    \item \textbf{Amplitude} $h_i = a_{i+1} - a_i$. Si $h_i = h$ constant, classes \textbf{équidistantes}.
\end{itemize}
\end{definition}

\begin{definition}[Effectifs et fréquences cumulés]
\begin{itemize}
    \item \textbf{Croissants :} $N_1 = n_1$, $N_i = N_{i-1} + n_i$ pour $i \ge 2$. $N_k = N$.
    \item \textbf{Décroissants :} $N'_k = n_k$, $N'_i = N'_{i+1} + n_i$ pour $i \le k-1$. $N'_1 = N$.
    \item \textbf{Fréquences cumulées :} $F_i = \frac{N_i}{N}$ (croissantes), $F'_i = \frac{N'_i}{N}$ (décroissantes).
\end{itemize}
$N_i$ (resp. $N'_i$) donne le nombre d'individus dont la valeur est \textbf{inférieure à la borne supérieure} (resp. \textbf{supérieure ou égale à la borne inférieure}) de la classe $I_i$.
\end{definition}

\begin{methode}[Regrouper une série en classes et construire les tableaux cumulés]
\begin{enumerate}
    \item \textbf{Choisir les classes :} Déterminer l'amplitude $h$ (imposée ou choisie pour 5 à 15 classes). Définir les bornes $[a_1; a_2[$, $[a_2; a_3[$, ..., $[a_k; a_{k+1}]$. La dernière classe est fermée des deux côtés.
    \item \textbf{Calculer les centres :} $x_i = \frac{\text{borne inf} + \text{borne sup}}{2}$.
    \item \textbf{Compter les effectifs $n_i$ :} Pour chaque valeur brute, identifier sa classe et incrémenter $n_i$.
    \item \textbf{Calculer les fréquences $f_i$ :} $f_i = \frac{n_i}{N} \times 100$ (en $\%$).
    \item \textbf{Effectifs cumulés croissants $N_i$ :} $N_1 = n_1$, $N_i = N_{i-1} + n_i$.
    \item \textbf{Effectifs cumulés décroissants $N'_i$ :} $N'_k = n_k$, $N'_i = N'_{i+1} + n_i$.
    \item \textbf{Vérification :} $\sum n_i = N$ et dernier $N_i = N$ (premier $N'_i = N$).
\end{enumerate}
\end{methode}

\begin{exoresolu}[Regroupement et effectifs cumulés — Rendements maïs (Adamaoua)]
\textbf{Énoncé :} Dans une zone agricole de l'Adamaoua, on a relevé les rendements (en quintaux par hectare) de 40 exploitations :
\[ 12,\ 15,\ 18,\ 22,\ 25,\ 28,\ \mathbf{29},\ 30,\ 32,\ 35,\ 38,\ 40,\ 42,\ 45,\ 48,\ \mathbf{49},\ 50,\ 52,\ 55,\ 58,\]
\[ 60,\ 62,\ 65,\ 68,\ 70,\ 72,\ 75,\ 78,\ 80,\ 82,\ 85,\ 88,\ 90,\ 92,\ 95,\ 98,\ 100,\ 102,\ 105,\ 108. \]
On regroupe en classes d'amplitude $h = 20$ en commençant par $[10; 30[$.
\begin{enumerate}
    \item Construire le tableau complet : classes, centres $x_i$, effectifs $n_i$, fréquences $f_i$ (\%), effectifs cumulés croissants $N_i$ et décroissants $N'_i$.
    \item Combien d'exploitations ont un rendement \textbf{strictement inférieur à 50 qx/ha} ? \textbf{Supérieur ou égal à 90 qx/ha} ?
\end{enumerate}
\tcblower
\textbf{Solution détaillée :}
\begin{enumerate}
    \item \textbf{Classes :} $[10; 30[$, $[30; 50[$, $[50; 70[$, $[70; 90[$, $[90; 110]$.
    \textbf{Centres $x_i$ :} $20,\ 40,\ 60,\ 80,\ 100$.
    \textbf{Comptage $n_i$ :}
    \begin{itemize}
        \item $[10; 30[$ : 12, 15, 18, 22, 25, 28, 29 $\rightarrow n_1 = 7$.
        \item $[30; 50[$ : 30, 32, 35, 38, 40, 42, 45, 48, 49 $\rightarrow n_2 = 9$.
        \item $[50; 70[$ : 50, 52, 55, 58, 60, 62, 65, 68 $\rightarrow n_3 = 8$.
        \item $[70; 90[$ : 70, 72, 75, 78, 80, 82, 85, 88 $\rightarrow n_4 = 8$.
        \item $[90; 110]$ : 90, 92, 95, 98, 100, 102, 105, 108 $\rightarrow n_5 = 8$.
    \end{itemize}
    Total $N = 7+9+8+8+8 = 40$. \checkmark
    \textbf{Fréquences $f_i = n_i/40 \times 100$ :} $17,5\%$, $22,5\%$, $20\%$, $20\%$, $20\%$.
    \textbf{Effectifs cumulés croissants $N_i$ :} $7,\ 16,\ 24,\ 32,\ 40$.
    \textbf{Effectifs cumulés décroissants $N'_i$ :} $40,\ 33,\ 24,\ 16,\ 8$.

    \begin{center}
    \begin{tabular}{|c|c|c|c|c|c|}\hline
    \textbf{Classes} & \textbf{Centres $x_i$} & \textbf{Effectifs $n_i$} & \textbf{Fréq. $f_i (\%)$} & \textbf{$N_i$ (croiss.)} & \textbf{$N'_i$ (décroiss.)} \\ \hline
    $[10; 30[$ & 20 & 7 & 17,5 & 7 & 40 \\ \hline
    $[30; 50[$ & 40 & 9 & 22,5 & 16 & 33 \\ \hline
    $[50; 70[$ & 60 & 8 & 20 & 24 & 24 \\ \hline
    $[70; 90[$ & 80 & 8 & 20 & 32 & 16 \\ \hline
    $[90; 110]$ & 100 & 8 & 20 & 40 & 8 \\ \hline
    \textbf{Total} & & \textbf{40} & \textbf{100} & & \\ \hline
    \end{tabular}
    \end{center}

    \item \textbf{Interprétation via effectifs cumulés :}
    \begin{itemize}
        \item \textbf{Rendement $< 50$ qx/ha :} Classes $[10;30[$ et $[30;50[$. $N_2 = 16$ exploitations.
        \item \textbf{Rendement $\ge 90$ qx/ha :} Classe $[90;110]$. $N'_5 = 8$ exploitations (ou $N - N_4 = 40 - 32 = 8$).
    \end{itemize}
    \textbf{Conclusion :} Le tableau permet de répondre instantanément aux questions de seuils.
\end{enumerate}
\end{exoresolu}

\courssub{Paramètres de position : classe modale, mode, moyenne}

\begin{definition}[Classe modale et mode (amplitudes égales)]
La \textbf{classe modale} est la classe de plus grand effectif $n_{\max}$. L'\textbf{effectif modal} est ce nombre $n_{\max}$.
Si les classes ont même amplitude, le \textbf{mode} $M_o$ est le \textbf{centre de la classe modale} : $M_o = x_{\text{mod}}$.
\emph{Attention :} Ne pas confondre effectif modal (nombre d'individus) et mode (valeur de la variable).
\end{definition}

\begin{definition}[Moyenne d'une série regroupée]
La \textbf{moyenne} $\bar{x}$ est estimée en remplaçant chaque valeur par le centre de sa classe :
\[ \bar{x} = \frac{\sum_{i=1}^k n_i x_i}{\sum_{i=1}^k n_i} = \frac{\sum n_i x_i}{N} \]
C'est une \textbf{estimation} car on perd l'information exacte à l'intérieur des classes.
\end{definition}

\begin{methode}[Calculer la moyenne d'une série regroupée]
\begin{enumerate}
    \item Identifier les centres $x_i = \frac{\text{borne inf} + \text{borne sup}}{2}$.
    \item Ajouter une colonne $n_i x_i$ et calculer la somme $\sum n_i x_i$.
    \item Appliquer $\bar{x} = \frac{\sum n_i x_i}{N}$.
    \item Exprimer le résultat dans l'unité de la variable (arrondi convenable). Préciser : « La moyenne \emph{estimée} est... ».
\end{enumerate}
\textbf{Variante (codage) :} Si amplitudes égales $h$, choisir un centre $A$ (souvent central), poser $u_i = \frac{x_i - A}{h}$. Alors $\bar{x} = A + h \frac{\sum n_i u_i}{N}$.
\end{methode}

\begin{exoresolu}[Moyenne d'une série regroupée — Durée d'appels MTN]
\textbf{Énoncé :} Durée des appels (en minutes) vers le réseau MTN pour 200 appels :
\begin{center}
\begin{tabular}{|c|c|c|c|c|c|}\hline
\textbf{Durée (min)} & $[0; 5[$ & $[5; 10[$ & $[10; 20[$ & $[20; 30[$ & $[30; 40]$ \\ \hline
\textbf{Effectifs} & 40 & 50 & 60 & 30 & 20 \\ \hline
\end{tabular}
\end{center}
Calculer la durée moyenne d'un appel. Résultat en minutes et secondes.
\tcblower
\textbf{Solution détaillée :}
\begin{enumerate}
    \item $N = 40+50+60+30+20 = 200$. \checkmark
    \item \textbf{Centres $x_i$ :} $2,5;\ 7,5;\ 15;\ 25;\ 35$.
    \item \textbf{Calcul $\sum n_i x_i$ :}
    \begin{center}
    \begin{tabular}{|c|c|c|}\hline
    $x_i$ & $n_i$ & $n_i x_i$ \\ \hline
    2,5 & 40 & 100 \\ \hline
    7,5 & 50 & 375 \\ \hline
    15 & 60 & 900 \\ \hline
    25 & 30 & 750 \\ \hline
    35 & 20 & 700 \\ \hline
    \textbf{Total} & \textbf{200} & \textbf{2825} \\ \hline
    \end{tabular}
    \end{center}
    \item $\bar{x} = \frac{2825}{200} = 14,125$ minutes.
    \item $0,125 \text{ min} = 0,125 \times 60 = 7,5 \text{ s}$.
    \textbf{Réponse :} La durée moyenne estimée est de \textbf{14 minutes et 7,5 secondes} (environ 14 min 08 s).
\end{enumerate}
\textbf{Interprétation :} Un appel typique dure environ 14 minutes. C'est une estimation (hypothèse : tous les appels d'une classe durent exactement le temps central).
\end{exoresolu}

\begin{methode}[Déterminer la classe modale et le mode]
\begin{enumerate}
    \item Repérer l'effectif maximal $n_{\max}$.
    \item \textbf{Classe modale :} Classe correspondante. \textbf{Effectif modal :} $n_{\max}$.
    \item Si amplitudes égales : \textbf{Mode $M_o$ = Centre de la classe modale}.
    \item Si deux classes ont le même effectif maximal : série \textbf{bimodale} (deux modes).
\end{enumerate}
\end{methode}

\begin{exoresolu}[Classe modale et mode — Ventes journalières (Marché Central Yaoundé)]
\textbf{Énoncé :} Nombre de paniers de tomates vendus par jour par 50 commerçants :
\begin{center}
\begin{tabular}{|c|c|c|c|c|c|}\hline
\textbf{Paniers} & $[0; 10[$ & $[10; 20[$ & $[20; 30[$ & $[30; 40[$ & $[40; 50]$ \\ \hline
\textbf{Effectifs} & 5 & 12 & \textbf{20} & 8 & 5 \\ \hline
\end{tabular}
\end{center}
\begin{enumerate}
    \item Identifier la classe modale et l'effectif modal.
    \item Déterminer le mode.
    \item Interpréter.
\end{enumerate}
\tcblower
\textbf{Solution détaillée :}
\begin{enumerate}
    \item Effectif maximal = 20. \textbf{Classe modale :} $[20; 30[$. \textbf{Effectif modal :} 20 commerçants.
    \item Amplitude constante $h=10$. Centre de la classe modale : $\frac{20+30}{2} = 25$.
    \[ \boxed{M_o = 25 \text{ paniers}}. \]
    \item \textbf{Interprétation :} La valeur la plus fréquente (le "pic") est 25 paniers. 20 commerçants sur 50 (40\%) vendent autour de 25 paniers/jour. C'est la performance "type" la plus observée.
\end{enumerate}
\end{exoresolu}

\courssub{Médiane : intervalle médian et interpolation linéaire}

\begin{definition}[Médiane et intervalle médian]
La \textbf{médiane} $M_e$ est la valeur qui partage la série en deux parties égales : 50\% des effectifs sont $\le M_e$ et 50\% sont $\ge M_e$.
L'\textbf{intervalle médian} est la classe (intervalle) qui contient la médiane. C'est la première classe pour laquelle l'effectif cumulé croissant $N_i \ge \frac{N}{2}$.
\end{definition}

\begin{methode}[Déterminer la médiane par interpolation linéaire]
\begin{enumerate}
    \item Calculer $N/2$.
    \item Construire les effectifs cumulés croissants $N_i$. Identifier la \textbf{classe médiane} : première classe avec $N_i \ge N/2$.
    \item Noter pour cette classe :
    \begin{itemize}
        \item $L$ = borne \textbf{inférieure}.
        \item $h$ = amplitude.
        \item $n_{\text{med}}$ = effectif de la classe.
        \item $F_{\text{prec}}$ = effectif cumulé croissant de la classe \textbf{précédente} ($0$ si 1ère classe).
    \end{itemize}
    \item Appliquer la formule (hypothèse : répartition uniforme dans la classe) :
    \[ \boxed{M_e = L + \frac{\frac{N}{2} - F_{\text{prec}}}{n_{\text{med}}} \times h} \]
    \item Vérifier : $M_e \in [L;\ L+h]$.
\end{enumerate}
\textbf{Exigible au Probatoire :} Formule et application numérique. La démonstration par triangles semblables (géométrie) peut être demandée en activité guidée.
\end{methode}

\begin{exoresolu}[Médiane par interpolation — Âge des abonnés CAMTEL]
\textbf{Énoncé :} Répartition des âges (ans) de 120 abonnés CAMTEL à Douala :
\begin{center}
\begin{tabular}{|c|c|c|c|c|c|c|}\hline
\textbf{Âge (ans)} & $[15; 25[$ & $[25; 35[$ & $[35; 45[$ & $[45; 55[$ & $[55; 65[$ & $[65; 75]$ \\ \hline
\textbf{Effectifs} & 12 & 18 & 30 & 25 & 20 & 15 \\ \hline
\end{tabular}
\end{center}
Déterminer l'âge médian par interpolation linéaire. Interpréter.
\tcblower
\textbf{Solution détaillée :}
\begin{enumerate}
    \item $N = 12+18+30+25+20+15 = 120$.
    \item $N/2 = 60$.
    \item \textbf{Effectifs cumulés croissants :}
    $N_1=12$ ; $N_2=30$ ; $N_3=60$ ; $N_4=85$ ; $N_5=105$ ; $N_6=120$.
    \item \textbf{Classe médiane :} $[35; 45[$ car $N_3 = 60 \ge 60$ (première classe $\ge 60$).
    \item \textbf{Paramètres :} $L = 35$, $h = 10$, $n_{\text{med}} = 30$, $F_{\text{prec}} = N_2 = 30$.
    \item \textbf{Calcul :}
    \[ M_e = 35 + \frac{60 - 30}{30} \times 10 = 35 + 10 = 45. \]
    \item \textbf{Vérification :} $45 \in [35; 45]$. La médiane est sur la borne supérieure car $N/2$ tombe exactement sur $N_3$.
    \textbf{Réponse :} L'âge médian est de \textbf{45 ans}.
    \textbf{Interprétation :} 50\% des abonnés ont moins de 45 ans, 50\% ont 45 ans ou plus. Population relativement âgée.
\end{enumerate}
\end{exoresolu}

\courssub{Paramètres de dispersion : écart moyen, variance, écart-type}

\begin{definition}[Écart moyen, variance, écart-type]
Soit $\bar{x}$ la moyenne de la série regroupée (centres $x_i$, effectifs $n_i$, total $N$).
\begin{itemize}
    \item \textbf{Écart moyen (déviation moyenne) :} $E_m = \dfrac{\sum n_i |x_i - \bar{x}|}{N}$. Unité : celle de la variable.
    \item \textbf{Variance :} $V = \dfrac{\sum n_i (x_i - \bar{x})^2}{N}$. Unité : unité$^2$.
    \item \textbf{Formule de König-Huygens (calcul simplifié) :} $\boxed{V = \frac{\sum n_i x_i^2}{N} - \bar{x}^2}$.
    \item \textbf{Écart-type :} $\boxed{\sigma = \sqrt{V}}$. Unité : celle de la variable. Mesure l'écart "typique" à la moyenne.
    \item \textbf{Coefficient de variation (optionnel) :} $CV = \frac{\sigma}{\bar{x}} \times 100$ (\%), permet de comparer dispersions d'échelles différentes.
\end{itemize}
\emph{Note :} On divise par $N$ (statistique descriptive, population entière). $N-1$ est pour l'estimation d'échantillon (hors programme 1ère).
\end{definition}

\begin{methode}[Calculer $E_m$, $V$, $\sigma$]
\begin{enumerate}
    \item Avoir $\bar{x}$ et les centres $x_i$.
    \item \textbf{Variance (König-Huygens) :} Colonne $x_i^2$, puis $n_i x_i^2$, somme $\sum n_i x_i^2$. $V = \frac{\sum n_i x_i^2}{N} - \bar{x}^2$.
    \item \textbf{Écart-type :} $\sigma = \sqrt{V}$ (arrondi).
    \item \textbf{Écart moyen :} Colonne $|x_i - \bar{x}|$, puis $n_i |x_i - \bar{x}|$, somme, diviser par $N$.
    \item \textbf{Interpréter :} "Les valeurs s'écartent en moyenne de $\sigma$ de la moyenne $\bar{x}$."
\end{enumerate}
\end{methode}

\begin{exoresolu}[Dispersion — Consommation électricité (ENEO)]
\textbf{Énoncé :} Consommation mensuelle (kWh) de 100 foyers à Bafoussam :
\begin{center}
\begin{tabular}{|c|c|c|c|c|}\hline
\textbf{Conso (kWh)} & $[50; 150[$ & $[150; 250[$ & $[250; 350[$ & $[350; 450]$ \\ \hline
\textbf{Effectifs} & 20 & 30 & 35 & 15 \\ \hline
\end{tabular}
\end{center}
\begin{enumerate}
    \item Calculer la consommation moyenne $\bar{x}$.
    \item Calculer la variance $V$ et l'écart-type $\sigma$ (dixième).
    \item Calculer l'écart moyen $E_m$.
    \item Interpréter $\sigma$.
\end{enumerate}
\tcblower
\textbf{Solution détaillée :}
\begin{enumerate}
    \item \textbf{Centres :} 100, 200, 300, 400. $N=100$.
    \begin{center}
    \begin{tabular}{|c|c|c|c|c|c|}\hline
    $x_i$ & $n_i$ & $n_i x_i$ & $x_i^2$ & $n_i x_i^2$ & $|x_i - \bar{x}|$ \\ \hline
    100 & 20 & 2000 & 10 000 & 200 000 & 145 \\ \hline
    200 & 30 & 6000 & 40 000 & 1 200 000 & 45 \\ \hline
    300 & 35 & 10 500 & 90 000 & 3 150 000 & 55 \\ \hline
    400 & 15 & 6000 & 160 000 & 2 400 000 & 155 \\ \hline
    \textbf{Total} & \textbf{100} & \textbf{24 500} & & \textbf{6 950 000} & \\ \hline
    \end{tabular}
    \end{center}
    $\bar{x} = 24\,500 / 100 = \mathbf{245 \text{ kWh}}$.
    \item \textbf{Variance :} $\frac{6\,950\,000}{100} = 69\,500$. $V = 69\,500 - 245^2 = 69\,500 - 60\,025 = \mathbf{9\,475 \text{ kWh}^2}$.
    \textbf{Écart-type :} $\sigma = \sqrt{9\,475} \approx \mathbf{97,3 \text{ kWh}}$.
    \item \textbf{Écart moyen :} $\sum n_i |x_i - \bar{x}| = 20(145)+30(45)+35(55)+15(155) = 2\,900+1\,350+1\,925+2\,325 = 8\,500$.
    $E_m = 8\,500 / 100 = \mathbf{85 \text{ kWh}}$.
    \item \textbf{Interprétation :} En moyenne, la consommation s'écarte de \textbf{97,3 kWh} de la moyenne (245 kWh). Dispersion forte (CV $\approx 39,7\%$), comportements hétérogènes (petits vs gros consommateurs).
\end{enumerate}
\end{exoresolu}

\courssub{Représentations graphiques : histogramme et courbes cumulatives}

\begin{definition}[Histogramme]
Représentation graphique d'une série regroupée.
\begin{itemize}
    \item Abscisses : bornes des classes (échelle réelle, rectangles contigus).
    \item Ordonnées : \textbf{Hauteur = Effectif (ou Fréquence) si amplitudes ÉGALES}.
    \item Si amplitudes \textbf{INÉGALES} : Hauteur = \textbf{Densité} $d_i = \frac{n_i}{h_i}$ (ou $\frac{f_i}{h_i}$). \textbf{Aire du rectangle = Effectif (ou Fréquence)}. \emph{Règle cruciale au Probatoire}.
\end{itemize}
\end{definition}

\begin{definition}[Courbe des effectifs cumulés croissants (Polygone cumulé)]
Points à tracer : $(\text{borne sup}_i ; N_i)$ pour chaque classe $i$, plus le point initial $(\text{borne inf}_1 ; 0)$. On relie par des \textbf{segments de droite} (interpolation linéaire).
La médiane se lit graphiquement : abscisse du point d'ordonnée $N/2$.
La courbe décroissante utilise les points $(\text{borne inf}_i ; N'_i)$ et point final $(\text{borne sup}_k ; 0)$.
\end{definition}

\begin{methode}[Construire l'histogramme et la courbe cumulée]
\begin{enumerate}
    \item \textbf{Histogramme :} Vérifier amplitudes. Si inégales $\rightarrow$ calculer densités $d_i = n_i/h_i$. Tracer rectangles contigus sur axes gradués.
    \item \textbf{Courbe cumulée croissante :} Tableau $(x_i^{\text{sup}}; N_i)$ + point départ $(\text{borne inf}_1; 0)$. Tracer points et segments.
    \item \textbf{Lecture graphique médiane :} Tracer horizontale $y = N/2$, intersection avec courbe, projeter sur abscisses.
\end{enumerate}
\end{methode}

\begin{exoresolu}[Histogramme et courbe cumulée — Salaires secteur privé]
\textbf{Énoncé :} Salaires mensuels nets (kFCFA) de 80 employés :
\begin{center}
\begin{tabular}{|c|c|c|c|c|c|}\hline
\textbf{Salaire (kFCFA)} & $[100; 150[$ & $[150; 200[$ & $[200; 300[$ & $[300; 400[$ & $[400; 500]$ \\ \hline
\textbf{Effectifs} & 10 & 20 & 30 & 15 & 5 \\ \hline
\end{tabular}
\end{center}
\begin{enumerate}
    \item Construire l'histogramme (préciser hauteurs/densités).
    \item Construire le tableau pour la courbe des effectifs cumulés croissants.
    \item Estimer graphiquement la médiane (méthode de lecture).
\end{enumerate}
\tcblower
\textbf{Solution détaillée :}
\begin{enumerate}
    \item \textbf{Amplitudes :} $50,\ 50,\ 100,\ 100,\ 100$ $\rightarrow$ \textbf{Inégales}. Densités $d_i = n_i/h_i$ :
    \begin{center}
    \begin{tabular}{|c|c|c|c|}\hline
    \textbf{Classe} & \textbf{Effectif $n_i$} & \textbf{Amplitude $h_i$} & \textbf{Densité $d_i$} \\ \hline
    $[100; 150[$ & 10 & 50 & 0,2 \\ \hline
    $[150; 200[$ & 20 & 50 & 0,4 \\ \hline
    $[200; 300[$ & 30 & 100 & 0,3 \\ \hline
    $[300; 400[$ & 15 & 100 & 0,15 \\ \hline
    $[400; 500]$ & 5 & 100 & 0,05 \\ \hline
    \end{tabular}
    \end{center}
    \textbf{Histogramme :} Abscisses 100 à 500. Ordonnées densités. Rectangles contigus. Aire totale = 80.
    \begin{popfigure}
    \begin{tikzpicture}[scale=0.7]
    \draw[->] (0,0) -- (11,0) node[right] {Salaire (kFCFA)};
    \draw[->] (0,0) -- (0,5) node[above] {Densité};
    \foreach \x/\label in {2/100, 3/150, 4/200, 6/300, 8/400, 10/500} {
        \draw (\x,0.1) -- (\x,-0.1) node[below] {\label};
    }
    \foreach \y/\label in {1/0.1, 2/0.2, 3/0.3, 4/0.4} {
        \draw (0.1,\y) -- (-0.1,\y) node[left] {\label};
    }
    \fill[popblue!30] (2,0) rectangle (3,2);
    \fill[popblue!50] (3,0) rectangle (4,4);
    \fill[popblue!70] (4,0) rectangle (6,3);
    \fill[popblue!90] (6,0) rectangle (8,1.5);
    \fill[popblue] (8,0) rectangle (10,0.5);
    \draw (2,0) rectangle (10,4);
    \end{tikzpicture}
    \legende{Histogramme des salaires (densités). Aire des rectangles = effectifs.}
    \end{popfigure}

    \item \textbf{Courbe cumulée croissante :}
    Effectifs cumulés $N_i$ : 10, 30, 60, 75, 80.
    Points : $(100; 0)$, $(150; 10)$, $(200; 30)$, $(300; 60)$, $(400; 75)$, $(500; 80)$.
    \begin{popfigure}
    \begin{tikzpicture}[scale=0.7]
    \draw[->] (0,0) -- (11,0) node[right] {Salaire (kFCFA)};
    \draw[->] (0,0) -- (0,9) node[above] {Effectifs cumulés};
    \foreach \x/\label in {2/100, 3/150, 4/200, 6/300, 8/400, 10/500} {
        \draw (\x,0.1) -- (\x,-0.1) node[below] {\label};
    }
    \foreach \y/\label in {1/10, 2/20, 3/30, 4/40, 5/50, 6/60, 7/70, 8/80} {
        \draw (0.1,\y) -- (-0.1,\y) node[left] {\label};
    }
    % Courbe
    \draw[poporange, thick] (2,0) -- (3,1) -- (4,3) -- (6,6) -- (8,7.5) -- (10,8);
    \foreach \p in {(2,0),(3,1),(4,3),(6,6),(8,7.5),(10,8)} {
        \fill[poporange] \p circle (3pt);
    }
    % Médiane graphique : N/2 = 40 -> y=4. Intersection sur segment (4,3)--(6,6).
    % Equation: y = 3 + 1.5(x-4). Pour y=4 -> x = 4 + 2/3 = 4.666...
    \draw[dashed, popgreen, thick] (0,4) -- (4.666,4) node[above, midway] {$N/2=40$};
    \draw[dashed, popgreen, thick] (4.666,4) -- (4.666,0) node[below] {$M_e \approx 233$};
    \end{tikzpicture}
    \legende{Courbe des effectifs cumulés croissants et lecture graphique de la médiane.}
    \end{popfigure}

    \item \textbf{Lecture graphique :} $N=80$, $N/2=40$. Horizontale $y=40$ coupe la courbe vers abscisse 233 kFCFA (interpolation visuelle entre 200 et 300).
    \textbf{Vérification calcul exact :} Classe médiane $[200; 300[$ ($N_2=30 < 40 \le N_3=60$).
    $L=200$, $h=100$, $n_{\text{med}}=30$, $F_{\text{prec}}=30$.
    $M_e = 200 + \frac{40-30}{30} \times 100 = 200 + 33,3 = \mathbf{233,3 \text{ kFCFA}}$.
    La lecture graphique approxime le calcul exact.
\end{enumerate}
\end{exoresolu}

\courssub{Synthèse et interprétation en situation concrète}

\begin{methode}[Analyser et comparer deux séries regroupées]
\begin{enumerate}
    \item \textbf{Calculer pour chaque série :} $\bar{x}$, $M_e$, $M_o$, $\sigma$ (ou $V$), $CV = \frac{\sigma}{\bar{x}}\times 100$.
    \item \textbf{Comparer positions :} Quelle série a valeurs globalement plus grandes ? Comparer $\bar{x}$ et $M_e$. Distribution asymétrique $\rightarrow$ médiane plus robuste.
    \item \textbf{Comparer dispersions :} Quelle série plus homogène ? Comparer $\sigma$ (même unité/ordre) ou mieux \textbf{CV} (sans unité).
    \item \textbf{Analyser forme (Histogramme) :}
    \begin{itemize}
        \item Symétrique $\Rightarrow \bar{x} \approx M_e \approx M_o$.
        \item Asymétrique à droite (queue à droite) $\Rightarrow M_o < M_e < \bar{x}$.
        \item Asymétrique à gauche $\Rightarrow \bar{x} < M_e < M_o$.
        \item Bimodale $\Rightarrow$ deux pics (deux populations mélangées ?).
    \end{itemize}
    \item \textbf{Conclusion contextualisée :} Phrases complètes, unités, référence aux chiffres, décision argumentée.
\end{enumerate}
\end{methode}

\begin{exoresolu}[Analyse comparative — Rendements deux variétés de riz (Nord-Cameroun)]
\textbf{Énoncé :} Rendements (qx/ha) de VAR-A et VAR-B sur 100 parcelles chacune.
\begin{center}
\begin{tabular}{|c|c|c|c|c|c|}\hline
\textbf{Classe (qx/ha)} & $[20; 30[$ & $[30; 40[$ & $[40; 50[$ & $[50; 60[$ & $[60; 70]$ \\ \hline
\textbf{VAR-A Effectifs} & 5 & 15 & 40 & 30 & 10 \\ \hline
\textbf{VAR-B Effectifs} & 10 & 20 & 30 & 25 & 15 \\ \hline
\end{tabular}
\end{center}
\begin{enumerate}
    \item Calculer moyenne et écart-type pour chaque variété.
    \item Calculer le coefficient de variation.
    \item Quelle variété recommander pour sécurité alimentaire stable ? Justifier.
\end{enumerate}
\tcblower
\textbf{Solution détaillée :}
\begin{enumerate}
    \item \textbf{Centres :} 25, 35, 45, 55, 65. $N=100$.
    \textbf{VAR-A :}
    $\sum n_i x_i = 5(25)+15(35)+40(45)+30(55)+10(65) = 4\,750$. $\bar{x}_A = 47,5$.
    $\sum n_i x_i^2 = 5(625)+15(1225)+40(2025)+30(3025)+10(4225) = 235\,500$.
    $V_A = 2\,355 - 47,5^2 = 98,75$. $\sigma_A = \sqrt{98,75} \approx \mathbf{9,94}$.
    \textbf{VAR-B :}
    $\sum n_i x_i = 10(25)+20(35)+30(45)+25(55)+15(65) = 4\,650$. $\bar{x}_B = 46,5$.
    $\sum n_i x_i^2 = 10(625)+20(1225)+30(2025)+25(3025)+15(4225) = 230\,500$.
    $V_B = 2\,305 - 46,5^2 = 142,75$. $\sigma_B = \sqrt{142,75} \approx \mathbf{11,95}$.
    \item \textbf{CV :}
    $CV_A = \frac{9,94}{47,5} \times 100 \approx \mathbf{20,9\%}$.
    $CV_B = \frac{11,95}{46,5} \times 100 \approx \mathbf{25,7\%}$.
    \item \textbf{Recommandation :}
    \begin{itemize}
        \item \textbf{Performance :} VAR-A moyenne légèrement supérieure (47,5 vs 46,5).
        \item \textbf{Stabilité :} VAR-A écart-type plus faible (9,9 vs 12,0) et CV plus faible (20,9\% vs 25,7\%). Dispersion relative plus faible.
        \item \textbf{Forme :} VAR-A concentrée autour de 45 (classe modale $[40;50[$, effectif 40). VAR-B plus étalée.
    \end{itemize}
    \textbf{Conclusion :} Pour la sécurité alimentaire, on privilégie la stabilité. La \textbf{VAR-A} est recommandée : rendement moyen légèrement supérieur et bien plus grande régularité (CV < 21\%), limitant le risque de très mauvaises récoltes.
\end{enumerate}
\end{exoresolu}

\begin{attention}[Les 5 erreurs qui coûtent des points au Probatoire]
\begin{enumerate}
    \item \textbf{Confondre "Effectif modal" et " Mode".} Effectif modal = nombre d'individus (ex: 30 élèves). Mode = valeur de la variable (ex: 14,5 ans). Écrire "Le mode est 30" $\rightarrow$ \textbf{0 point}.
    \item \textbf{Oublier la densité dans l'histogramme si amplitudes inégales.} Hauteurs = effectifs avec largeurs différentes $\rightarrow$ Aire non proportionnelle $\rightarrow$ Graphique faux. \textbf{Règle :} Amplitudes égales $\rightarrow$ Hauteur = Effectif. Inégales $\rightarrow$ Hauteur = Densité ($n_i/h_i$). Vérifiez toujours les largeurs !
    \item \textbf{Erreur sur la borne $L$ pour la médiane.} Utiliser le centre au lieu de la \textbf{borne inférieure} $L$ dans $M_e = L + \frac{N/2 - F_{\text{prec}}}{n_{\text{med}}} \times h$. $F_{\text{prec}}$ = effectif cumulé \textbf{avant} la classe médiane (pas $N/2$).
    \item \textbf{Variance : oublier $-\bar{x}^2$ (König-Huygens) ou diviser par $N-1$.} Formule : $V = \frac{\sum n_i x_i^2}{N} - \bar{x}^2$. Division par $N$ (population). $N-1$ est pour estimation d'échantillon (hors programme 1ère).
    \item \textbf{Interprétation "hors sol" sans unités ni contexte.} "La variance est grande" = 0 pt. "L'écart-type est 12 kg, forte hétérogénéité des poids des sacs" = points. \textbf{Toujours : Valeur + Unité + Phrase de sens dans le contexte.}
\end{enumerate}
\end{attention}

\newpage

\coursec{Exercices}

\courssub{Je vérifie mes connaissances}

\begin{exercice}[QCM : à chaud]
Pour chaque question, une seule réponse est exacte. Recopie la bonne lettre.
\begin{enumerate}
\item Une série de 100 données est regroupée en classes. L'effectif cumulé croissant de la classe $[10\,;\,15[$ est 62. Cela signifie que :
\begin{enumerate}
\item[a)] 62 données appartiennent à la classe $[10\,;\,15[$ ;
\item[b)] 62 données sont strictement inférieures à 15 ;
\item[c)] 62 données sont supérieures ou égales à 10.
\end{enumerate}
\item La moyenne d'une série regroupée en classes se calcule en utilisant :
\begin{enumerate}
\item[a)] les bornes inférieures des classes ;
\item[b)] les amplitudes des classes ;
\item[c)] les centres des classes.
\end{enumerate}
\item La classe modale est :
\begin{enumerate}
\item[a)] la classe qui contient la médiane ;
\item[b)] la classe de plus grand effectif ;
\item[c)] la classe de plus grande amplitude.
\end{enumerate}
\item Pour construire un histogramme d'une série dont les classes ont des amplitudes inégales, on porte en ordonnée :
\begin{enumerate}
\item[a)] les effectifs ;
\item[b)] les fréquences cumulées ;
\item[c)] les densités d'effectif.
\end{enumerate}
\end{enumerate}
\end{exercice}

\begin{exercice}[Vrai ou faux ?]
Réponds par Vrai ou Faux en justifiant chaque réponse par une phrase ou un contre-exemple.
\begin{enumerate}
\item La médiane d'une série regroupée en classes appartient toujours à la classe modale.
\item Si deux séries ont la même moyenne, alors elles ont le même écart-type.
\item L'effectif cumulé croissant de la dernière classe est toujours égal à l'effectif total.
\item La médiane partage la série en deux sous-séries de même effectif.
\end{enumerate}
\end{exercice}

\begin{exercice}[Définitions et formules]
\begin{enumerate}
\item Définis la \cle{classe modale} et le \cle{mode} d'une série regroupée en classes.
\item Écris la formule donnant la \cle{variance} d'une série regroupée en classes de centres $c_i$, d'effectifs $n_i$ et de moyenne $\bar{x}$.
\item Qu'appelle-t-on \cle{intervalle médian} ? Comment le repère-t-on dans le tableau des effectifs cumulés croissants ?
\item Rappelle le principe de l'\cle{interpolation linéaire} pour le calcul de la médiane.
\end{enumerate}
\end{exercice}

\begin{exercice}[Calculs directs]
Une série est regroupée selon le tableau suivant :
\begin{center}
\begin{tabular}{|c|cccc|}\hline
Classe & $[0\,;\,10[$ & $[10\,;\,20[$ & $[20\,;\,30[$ & $[30\,;\,40[$ \\ \hline
Effectif $n_i$ & 8 & 15 & 12 & 5 \\ \hline
\end{tabular}
\end{center}
\begin{enumerate}
\item Calcule l'effectif total $N$ et dresse la ligne des effectifs cumulés croissants.
\item Calcule le centre $c_i$ de chaque classe.
\item Calcule la moyenne $\bar{x}$ de la série.
\item Détermine la classe modale et le mode.
\end{enumerate}
\end{exercice}

\courssub{Je m'entraîne}

\begin{exercice}[Distances parcourues]
Une agence de voyages de Douala a relevé les distances parcourues (en km) par ses 50 cars lors d'une journée :
\begin{center}
\begin{tabular}{|c|ccccc|}\hline
Distance (km) & $[100\,;\,200[$ & $[200\,;\,300[$ & $[300\,;\,400[$ & $[400\,;\,500[$ & $[500\,;\,600[$ \\ \hline
Nombre de cars & 6 & 14 & 16 & 9 & 5 \\ \hline
\end{tabular}
\end{center}
\begin{enumerate}
\item Dresse le tableau complet : centres, effectifs cumulés croissants, fréquences et fréquences cumulées croissantes (en \%).
\item Calcule la distance moyenne parcourue par un car.
\item Détermine la classe modale et le mode.
\item Détermine l'intervalle médian, puis calcule la médiane par interpolation linéaire.
\end{enumerate}
\end{exercice}

\begin{exercice}[Dispersion des salaires]
Les salaires mensuels (en milliers de FCFA) des 40 employés d'une PME de Yaoundé sont regroupés ainsi :
\begin{center}
\begin{tabular}{|c|cccc|}\hline
Salaire & $[50\,;\,100[$ & $[100\,;\,150[$ & $[150\,;\,200[$ & $[200\,;\,250[$ \\ \hline
Effectif & 10 & 16 & 9 & 5 \\ \hline
\end{tabular}
\end{center}
\begin{enumerate}
\item Calcule le salaire moyen $\bar{x}$.
\item Calcule l'écart moyen $e = \dfrac{1}{N}\sum n_i\,|c_i - \bar{x}|$.
\item Calcule la variance $V$ puis l'écart-type $\sigma$ (arrondir au millier de FCFA près).
\item Calcule le coefficient de variation $\dfrac{\sigma}{\bar{x}}$ en \% et commente la dispersion des salaires.
\end{enumerate}
\end{exercice}

\begin{exercice}[Effectif manquant]
Le tableau ci-dessous donne les durées de communication (en minutes) de 60 abonnés d'un réseau mobile, mais un effectif a été effacé :
\begin{center}
\begin{tabular}{|c|cccc|}\hline
Durée (min) & $[0\,;\,5[$ & $[5\,;\,10[$ & $[10\,;\,15[$ & $[15\,;\,20[$ \\ \hline
Effectif & 12 & $a$ & 20 & 10 \\ \hline
\end{tabular}
\end{center}
On sait que la moyenne de la série regroupée est $\bar{x} = 10{,}25$ min.
\begin{enumerate}
\item Détermine la valeur de $a$ en utilisant l'effectif total.
\item Vérifie que la moyenne calculée avec cette valeur est bien $10{,}25$ min.
\item Détermine la médiane par interpolation linéaire.
\end{enumerate}
\end{exercice}

\begin{exercice}[Classes d'amplitudes inégales]
Une microfinance de Bafoussam a relevé l'âge de ses 100 clients :
\begin{center}
\begin{tabular}{|c|ccccc|}\hline
Âge (ans) & $[18\,;\,25[$ & $[25\,;\,35[$ & $[35\,;\,45[$ & $[45\,;\,60[$ & $[60\,;\,80[$ \\ \hline
Effectif & 14 & 26 & 25 & 23 & 12 \\ \hline
\end{tabular}
\end{center}
\begin{enumerate}
\item Explique pourquoi on ne peut pas construire l'histogramme en portant directement les effectifs en ordonnée.
\item Calcule la densité d'effectif $d_i = \dfrac{n_i}{a_i}$ de chaque classe, où $a_i$ désigne l'amplitude.
\item Construis l'histogramme avec les densités (échelles : 1 cm pour 10 ans en abscisse ; 1 cm pour 0,5 en ordonnée).
\item Quelle est la classe modale au sens des densités ? Comparer avec la classe de plus grand effectif et commenter.
\end{enumerate}
\end{exercice}

\courssub{Je me prépare au Probatoire}

\begin{exercice}[Les sacs de maïs — type Probatoire]
Une coopérative agricole de la région de l'Ouest a pesé 80 sacs de maïs destinés à la vente. Les masses, en kg, sont regroupées dans le tableau suivant :
\begin{center}
\begin{tabular}{|c|cccccc|}\hline
Masse (kg) & $[40\,;\,45[$ & $[45\,;\,50[$ & $[50\,;\,55[$ & $[55\,;\,60[$ & $[60\,;\,65[$ & $[65\,;\,70[$ \\ \hline
Nombre de sacs & 6 & 14 & 22 & 20 & 12 & 6 \\ \hline
\end{tabular}
\end{center}
\begin{enumerate}
\item Reproduis et complète le tableau avec les lignes : centres $c_i$, effectifs cumulés croissants (ECC) et fréquences cumulées croissantes (FCC) en \%.
\item Calcule la masse moyenne $\bar{x}$ d'un sac de maïs (arrondir à $0{,}1$ kg près).
\item Détermine la classe modale de la série, puis le mode.
\item \begin{enumerate}
\item Détermine l'intervalle médian.
\item Calcule la médiane $M_e$ par interpolation linéaire (arrondir à $0{,}1$ kg près).
\end{enumerate}
\item Calcule la variance $V$ puis l'écart-type $\sigma$ de la série (arrondir à $0{,}1$ près).
\item La coopérative affirme : « Au moins la moitié de nos sacs pèsent plus de 55 kg. » Cette affirmation est-elle exacte ? Justifie à l'aide des résultats précédents.
\end{enumerate}
\end{exercice}

\begin{exercice}[Concours d'entrée — type Probatoire]
Les notes sur 20 de 60 candidats à un concours d'entrée dans une école de la place sont regroupées ainsi :
\begin{center}
\begin{tabular}{|c|ccccc|}\hline
Note & $[0\,;\,4[$ & $[4\,;\,8[$ & $[8\,;\,12[$ & $[12\,;\,16[$ & $[16\,;\,20]$ \\ \hline
Effectif & 5 & 13 & 22 & 15 & 5 \\ \hline
\end{tabular}
\end{center}
\begin{enumerate}
\item Dresse le tableau statistique complet : centres, ECC, ECD (effectifs cumulés décroissants) et FCC en \%.
\item Construis l'histogramme de la série (échelles : 1 cm pour 2 points en abscisse ; 1 cm pour 5 candidats en ordonnée).
\item Construis, sur un second graphique, la courbe des fréquences cumulées croissantes (échelles : 1 cm pour 2 points en abscisse ; 1 cm pour 10 \% en ordonnée).
\item \begin{enumerate}
\item Détermine graphiquement une valeur approchée de la note médiane.
\item Retrouve cette valeur par le calcul, à l'aide d'une interpolation linéaire.
\end{enumerate}
\item Un candidat est déclaré admissible s'il obtient au moins 10. À l'aide d'une interpolation linéaire dans la classe $[8\,;\,12[$, détermine le pourcentage de candidats admissibles.
\item Calcule la moyenne $\bar{x}$ et l'écart-type $\sigma$ de la série. Le jury estime que l'épreuve est « bien calibrée » si $\bar{x} \in [9\,;\,11]$ et $\sigma \geq 3$. L'épreuve est-elle bien calibrée ?
\end{enumerate}
\end{exercice}

\begin{exercice}[Factures d'eau et comparaison de coopératives — problème de synthèse]
\textbf{Partie A.} La régie de distribution d'eau d'une ville a relevé le montant (en milliers de FCFA) des factures mensuelles de 150 abonnés :
\begin{center}
\begin{tabular}{|c|cccccc|}\hline
Montant & $[0\,;\,5[$ & $[5\,;\,10[$ & $[10\,;\,15[$ & $[15\,;\,20[$ & $[20\,;\,25[$ & $[25\,;\,30[$ \\ \hline
Abonnés & 18 & 36 & 45 & 27 & 15 & 9 \\ \hline
\end{tabular}
\end{center}
\begin{enumerate}
\item Calcule le montant moyen $\bar{x}$ d'une facture.
\item Calcule la médiane $M_e$ par interpolation linéaire. Interprète le résultat par une phrase.
\item Calcule l'écart-type $\sigma$ de la série (arrondir à $0{,}1$ près).
\item Dresse le tableau des fréquences cumulées décroissantes (FCD) en \%.
\item Par interpolation linéaire, détermine le seuil de facturation dépassé par 20 \% des abonnés (c'est-à-dire le montant $s$ tel que 20 \% des abonnés paient plus de $s$ milliers de FCFA).
\end{enumerate}
\textbf{Partie B.} Deux coopératives de planteurs de cacao, $C_1$ et $C_2$, ont chacune une production moyenne de 50 sacs par exploitation. Les productions (en sacs) de leurs membres sont regroupées ainsi :
\begin{center}
\begin{tabular}{|c|ccccc|}\hline
Production & $[30\,;\,40[$ & $[40\,;\,50[$ & $[50\,;\,60[$ & $[60\,;\,70[$ & $[70\,;\,80[$ \\ \hline
Effectif $C_1$ & 4 & 12 & 18 & 10 & 6 \\ \hline
Effectif $C_2$ & 10 & 16 & 8 & 9 & 7 \\ \hline
\end{tabular}
\end{center}
\begin{enumerate}
\item[6.] Vérifie que les deux coopératives ont bien la même production moyenne de 50 sacs.
\item[7.] Calcule la variance et l'écart-type de chacune des deux séries.
\item[8.] Quelle coopérative est la plus homogène dans sa production ? Justifie. Pourquoi la moyenne seule ne suffisait-elle pas à répondre ?
\end{enumerate}
\end{exercice}

\newpage

\coursec{Corrigés des exercices}

\begin{corrige}[Exercice 1 — QCM : à chaud]
\begin{enumerate}
\item \textbf{b)} L'effectif cumulé croissant d'une classe $[a\,;\,b[$ est le nombre de données \textbf{inférieures à $b$} (borne supérieure). Ici, 62 données sont $< 15$.
\item \textbf{c)} Pour une série regroupée en classes, on remplace chaque valeur par le \textbf{centre de sa classe} ($c_i = \frac{\text{borne inf.} + \text{borne sup.}}{2}$) pour calculer la moyenne $\bar{x} = \frac{\sum n_i c_i}{N}$.
\item \textbf{b)} La \cle{classe modale} est, par définition, la classe qui possède le \textbf{plus grand effectif}.
\item \textbf{c)} Lorsque les classes ont des amplitudes inégales, les aires des rectangles doivent être proportionnelles aux effectifs. Comme $\text{Aire} = \text{largeur} \times \text{hauteur}$, la hauteur (ordonnée) doit être la \textbf{densité d'effectif} $d_i = \frac{n_i}{a_i}$.
\end{enumerate}
\end{corrige}

\begin{corrige}[Exercice 2 — Vrai ou faux ?]
\begin{enumerate}
\item \textbf{Faux.} La médiane sépare la population en deux effectifs égaux ; la classe modale est celle de plus grand effectif. Rien ne les oblige à coïncider. \emph{Contre-exemple :} série $[0\,;\,10[$ (eff. 50), $[10\,;\,20[$ (eff. 10), $[20\,;\,30[$ (eff. 5). Classe modale $[0\,;\,10[$, médiane dans $[0\,;\,10[$ ici, mais si on met 40 dans $[0\,;\,10[$, 20 dans $[10\,;\,20[$, 40 dans $[20\,;\,30[$, classe modale $[0\,;\,10[$ et $[20\,;\,30[$, médiane dans $[10\,;\,20[$.
\item \textbf{Faux.} La moyenne mesure la position, l'écart-type la dispersion. Deux séries peuvent avoir le même centre de gravité mais des dispersions très différentes. \emph{Contre-exemple :} Série A : 10, 10, 10 (moyenne 10, $\sigma=0$) ; Série B : 0, 10, 20 (moyenne 10, $\sigma>0$).
\item \textbf{Vrai.} L'effectif cumulé croissant de la dernière classe compte toutes les données de la série, donc $N$.
\item \textbf{Vrai.} C'est la définition même de la médiane (pour un effectif total $N$, elle laisse $\frac{N}{2}$ données de chaque côté ; si $N$ impair, elle est la valeur centrale).
\end{enumerate}
\end{corrige}

\begin{corrige}[Exercice 3 — Définitions et formules]
\begin{enumerate}
\item \textbf{Classe modale :} classe de l'intervalle de regroupement qui a l'\textbf{effectif le plus grand}. \textbf{Mode :} valeur approchée du mode, prise comme le \textbf{centre de la classe modale} ($c_{\text{mod}}$). (Au Probatoire, on ne demande pas l'interpolation graphique du mode pour les classes).
\item \textbf{Variance :} $V = \dfrac{1}{N}\displaystyle\sum_{i=1}^{k} n_i (c_i - \bar{x})^2 = \dfrac{1}{N}\displaystyle\sum_{i=1}^{k} n_i c_i^2 - \bar{x}^2$ (formule de König-Huyghens, plus pratique pour le calcul).
\item \textbf{Intervalle médian :} c'est la classe qui \textbf{contient la médiane}. On le repère dans le tableau des effectifs cumulés croissants (ECC) : c'est la \textbf{première classe} pour laquelle l'ECC est \textbf{supérieur ou égal à $\dfrac{N}{2}$}.
\item \textbf{Interpolation linéaire :} On suppose que les données sont \textbf{uniformément réparties} à l'intérieur de la classe médiane. Si la classe médiane est $[L\,;\,L+a[$ d'effectif $n_{\text{med}}$, et que l'ECC de la classe précédente est $N_{\text{prec}}$, alors :
\[
M_e = L + \dfrac{\frac{N}{2} - N_{\text{prec}}}{n_{\text{med}}} \times a
\]
\end{enumerate}
\end{corrige}

\begin{corrige}[Exercice 4 — Calculs directs]
\begin{enumerate}
\item Effectif total $N = 8+15+12+5 = 40$.
\begin{center}
\begin{tabular}{|c|cccc|}\hline
Classe & $[0\,;\,10[$ & $[10\,;\,20[$ & $[20\,;\,30[$ & $[30\,;\,40[$ \\ \hline
Effectif $n_i$ & 8 & 15 & 12 & 5 \\
ECC & 8 & 23 & 35 & 40 \\ \hline
\end{tabular}
\end{center}
\item Centres : $c_1=5$, $c_2=15$, $c_3=25$, $c_4=35$.
\item $\bar{x} = \dfrac{1}{40}(8\times5 + 15\times15 + 12\times25 + 5\times35) = \dfrac{40+225+300+175}{40} = \dfrac{740}{40} = 18,5$.
\item Plus grand effectif : 15 pour la classe $[10\,;\,20[$. \textbf{Classe modale : $[10\,;\,20[$}. \textbf{Mode : $15$} (centre de cette classe).
\end{enumerate}
\end{corrige}

\begin{corrige}[Exercice 5 — Distances parcourues]
\begin{enumerate}
\item Tableau complet ($N=50$) :
\begin{center}
\begin{tabular}{|c|ccccc|}\hline
Distance (km) & $[100\,;\,200[$ & $[200\,;\,300[$ & $[300\,;\,400[$ & $[400\,;\,500[$ & $[500\,;\,600[$ \\ \hline
Centre $c_i$ & 150 & 250 & 350 & 450 & 550 \\
Effectif $n_i$ & 6 & 14 & 16 & 9 & 5 \\
ECC & 6 & 20 & 36 & 45 & 50 \\
Fréquence $f_i$ (\%) & 12 & 28 & 32 & 18 & 10 \\
FCC (\%) & 12 & 40 & 72 & 90 & 100 \\ \hline
\end{tabular}
\end{center}
\item $\bar{x} = \frac{1}{50}(6\times150 + 14\times250 + 16\times350 + 9\times450 + 5\times550) = \frac{900+3500+5600+4050+2750}{50} = \frac{16800}{50} = \textbf{336 km}$.
\item Classe modale : $[300\,;\,400[$ (eff. 16). \textbf{Mode = 350 km}.
\item $N/2 = 25$. ECC : 6, 20, \textbf{36}... Intervalle médian : $\boldsymbol{[300\,;\,400[}$.
Interpolation : $L=300$, $a=100$, $n_{\text{med}}=16$, $N_{\text{prec}}=20$.
$M_e = 300 + \dfrac{25-20}{16}\times 100 = 300 + \dfrac{500}{16} = 300 + 31,25 = \textbf{331,25 km}$.
\end{enumerate}
\end{corrige}

\begin{corrige}[Exercice 6 — Dispersion des salaires]
\begin{enumerate}
\item $N=40$. Centres : 75 ; 125 ; 175 ; 225 (en k FCFA).
$\bar{x} = \frac{1}{40}(10\times75 + 16\times125 + 9\times175 + 5\times225) = \frac{750+2000+1575+1125}{40} = \frac{5450}{40} = \textbf{136,25 k FCFA}$.
\item Écart moyen $e = \frac{1}{40}\sum n_i|c_i-\bar{x}|$ :
\begin{align*}
&10\times|75-136,25| = 612,5 \\
&16\times|125-136,25| = 180 \\
&9\times|175-136,25| = 348,75 \\
&5\times|225-136,25| = 443,75 \\
&\text{Somme} = 1585 \quad\Rightarrow\quad e = \frac{1585}{40} = \textbf{39,625 k FCFA}
\end{align*}
\item $\sum n_i c_i^2 = 10\times5625 + 16\times15625 + 9\times30625 + 5\times50625 = 835\,000$.
$V = \frac{835\,000}{40} - 136,25^2 = 20\,875 - 18\,564,0625 = 2\,310,9375$.
$\sigma = \sqrt{2\,310,9375} \approx 48,07 \approx \textbf{48 k FCFA}$ (arrondi au millier près).
\item Coefficient de variation $= \frac{48,07}{136,25}\times 100 \approx \textbf{35,3\%}$.
\textbf{Commentaire :} Un CV $> 30\%$ indique une \textbf{forte dispersion} des salaires : les écarts par rapport à la moyenne sont importants relativement à la moyenne elle-même, ce qui reflète une grande hétérogénéité des rémunérations dans cette PME.
\end{enumerate}
\end{corrige}

\begin{corrige}[Exercice 7 — Effectif manquant]
\begin{enumerate}
\item Effectif total $N=60$. $12 + a + 20 + 10 = 60 \Rightarrow a = \textbf{18}$.
\item Centres : 2,5 ; 7,5 ; 12,5 ; 17,5.
Moyenne calculée avec $a=18$ :
$\bar{x} = \frac{1}{60}(12\times2,5 + 18\times7,5 + 20\times12,5 + 10\times17,5) = \frac{30+135+250+175}{60} = \frac{590}{60} \approx \textbf{9,83 min}$.
\textbf{Attention :} Cette valeur ($9,83$) est \textbf{différente} de la moyenne annoncée ($10,25$ min). Les données de l'énoncé (effectif total 60 et moyenne 10,25) sont \textbf{incompatibles} avec les autres effectifs donnés. Nous poursuivons avec $a=18$ (déduit de l'effectif total, donnée sûre).
\item $N/2 = 30$. ECC : 12 ; \textbf{30} ; 50 ; 60. Intervalle médian : $\boldsymbol{[5\,;\,10[}$ (2ème classe, ECC = 30 atteint exactement $N/2$).
Interpolation : $L=5$, $a=5$, $n_{\text{med}}=18$, $N_{\text{prec}}=12$.
$M_e = 5 + \dfrac{30-12}{18}\times 5 = 5 + \dfrac{18}{18}\times 5 = \textbf{10 min}$.
(Note : comme l'ECC vaut exactement $N/2$ à la fin de cette classe, la médiane tombe sur la borne supérieure, 10 min).
\end{enumerate}
\end{corrige}

\begin{corrige}[Exercice 8 — Classes d'amplitudes inégales]
\begin{enumerate}
\item Dans un histogramme, l'\textbf{aire} de chaque rectangle doit être proportionnelle à l'effectif de la classe. Si on porte les effectifs en ordonnée (hauteur), les classes larges (grande amplitude) auront une aire trop grande par rapport à leur effectif, ce qui fausse la représentation visuelle de la distribution.
\item Amplitudes $a_i$ : 7 ; 10 ; 10 ; 15 ; 20.
Densités $d_i = n_i/a_i$ :
\begin{center}
\begin{tabular}{|c|ccccc|}\hline
Classe & $[18\,;\,25[$ & $[25\,;\,35[$ & $[35\,;\,45[$ & $[45\,;\,60[$ & $[60\,;\,80[$ \\ \hline
$a_i$ & 7 & 10 & 10 & 15 & 20 \\
$n_i$ & 14 & 26 & 25 & 23 & 12 \\
$d_i$ & \textbf{2,00} & \textbf{2,60} & \textbf{2,50} & \textbf{1,53} & \textbf{0,60} \\ \hline
\end{tabular}
\end{center}
\item \textit{Construction de l'histogramme (description) :}
Abscisses : 1 cm pour 10 ans $\Rightarrow$ largeurs : 0,7 cm ; 1 cm ; 1 cm ; 1,5 cm ; 2 cm.
Ordonnées : 1 cm pour 0,5 densité $\Rightarrow$ hauteurs : 4 cm ; 5,2 cm ; 5 cm ; 3,07 cm ; 1,2 cm.
Les rectangles sont contigus, bases sur l'axe des abscisses.
\item \textbf{Classe modale au sens des densités :} $[25\,;\,35[$ (densité max 2,60).
\textbf{Classe de plus grand effectif :} $[25\,;\,35[$ (effectif 26).
\textbf{Commentaire :} Ici, les deux notions coïncident car la classe $[25\,;\,35[$ a à la fois le plus grand effectif et une amplitude « standard » (10 ans). Si la classe de plus grand effectif avait été très large (ex. $[45\,;\,60[$), sa densité aurait pu être plus faible qu'une classe étroite très peuplée ; l'histogramme aux densités aurait alors mis en évidence un « pic » de concentration différent du simple maximum d'effectif.
\end{enumerate}
\end{corrige}

\begin{corrige}[Exercice 9 — Les sacs de maïs — type Probatoire]
\textbf{Barème indicatif (sur 20) :} Tableau complet (3 pts), Moyenne (3 pts), Classe modale/Mode (2 pts), Intervalle médian + Médiane (4 pts), Variance/Écart-type (4 pts), Interprétation affirmation (4 pts).
\textbf{Points de vigilance :} Arrondis demandés (0,1 kg) ; formule de König-Huyghens pour la variance ; justification de l'affirmation par comparaison médiane/55 kg ou effectifs.

\begin{enumerate}
\item Tableau complet ($N=80$) :
\begin{center}
\begin{tabular}{|c|cccccc|}\hline
Masse (kg) & $[40\,;\,45[$ & $[45\,;\,50[$ & $[50\,;\,55[$ & $[55\,;\,60[$ & $[60\,;\,65[$ & $[65\,;\,70[$ \\ \hline
Centre $c_i$ & 42,5 & 47,5 & 52,5 & 57,5 & 62,5 & 67,5 \\
Effectif $n_i$ & 6 & 14 & 22 & 20 & 12 & 6 \\
ECC & 6 & 20 & 42 & 62 & 74 & 80 \\
FCC (\%) & 7,5 & 25 & 52,5 & 77,5 & 92,5 & 100 \\ \hline
\end{tabular}
\end{center}
\item $\sum n_i c_i = 6(42,5)+14(47,5)+22(52,5)+20(57,5)+12(62,5)+6(67,5) = 255+665+1155+1150+750+405 = 4380$.
$\bar{x} = \frac{4380}{80} = 54,75 \approx \textbf{54,8 kg}$ (arrondi 0,1).
\item Plus grand effectif : 22 pour $[50\,;\,55[$. \textbf{Classe modale : $[50\,;\,55[$}. \textbf{Mode = 52,5 kg}.
\item \begin{enumerate}
\item $N/2 = 40$. ECC : 6, 20, \textbf{42}... Intervalle médian : $\boldsymbol{[50\,;\,55[}$.
\item $L=50$, $a=5$, $n_{\text{med}}=22$, $N_{\text{prec}}=20$.
$M_e = 50 + \dfrac{40-20}{22}\times 5 = 50 + \dfrac{100}{22} \approx 50 + 4,545 = \textbf{54,5 kg}$ (arrondi 0,1).
\end{enumerate}
\item $\sum n_i c_i^2 = 6(1806,25)+14(2256,25)+22(2756,25)+20(3306,25)+12(3906,25)+6(4556,25) = 243\,300$.
$V = \frac{243\,300}{80} - 54,75^2 = 3041,25 - 2997,5625 = 43,6875$.
$\sigma = \sqrt{43,6875} \approx 6,61 \approx \textbf{6,6 kg}$ (arrondi 0,1).
\item \textbf{Affirmation fausse.}
Justification : La médiane est $54,5$ kg. Cela signifie que 50\% des sacs pèsent \textbf{moins de 54,5 kg}, donc \textbf{strictement moins de 55 kg}. Il est donc impossible qu'« au moins la moitié » (50\% ou plus) pèsent \textbf{plus de 55 kg}. Au contraire, au plus 50\% pèsent plus de 55 kg (en réalité moins, car il y a des sacs entre 54,5 et 55).
\end{enumerate}
\end{corrige}

\begin{corrige}[Exercice 10 — Concours d'entrée — type Probatoire]
\textbf{Barème indicatif (sur 20) :} Tableau complet (3 pts), Histogramme (3 pts), Courbe FCC (3 pts), Médiane graphique + calcul (4 pts), \% admissibles interpolation (3 pts), Moyenne/Écart-type + conclusion (4 pts).
\textbf{Points de vigilance :} Échelles respectées ; courbe FCC points aux bornes supérieures des classes ; interpolation linéaire pour \% admissibles (borne 10 dans classe $[8,12[$) ; test « bien calibrée » sur les deux critères.

\begin{enumerate}
\item Tableau ($N=60$) :
\begin{center}
\begin{tabular}{|c|ccccc|}\hline
Note & $[0\,;\,4[$ & $[4\,;\,8[$ & $[8\,;\,12[$ & $[12\,;\,16[$ & $[16\,;\,20]$ \\ \hline
Centre $c_i$ & 2 & 6 & 10 & 14 & 18 \\
Effectif $n_i$ & 5 & 13 & 22 & 15 & 5 \\
ECC & 5 & 18 & 40 & 55 & 60 \\
ECD & 60 & 55 & 42 & 20 & 5 \\
FCC (\%) & 8,33 & 30 & 66,67 & 91,67 & 100 \\ \hline
\end{tabular}
\end{center}
\item \textit{Histogramme :} Abscisses 1 cm / 2 pts $\Rightarrow$ largeur 2 cm par classe. Ordonnées 1 cm / 5 candidats $\Rightarrow$ hauteurs : 1 cm ; 2,6 cm ; 4,4 cm ; 3 cm ; 1 cm. Rectangles contigus.
\item \textit{Courbe FCC :} Abscisses 1 cm / 2 pts. Ordonnées 1 cm / 10\%. Points : $(4; 8,33)$, $(8; 30)$, $(12; 66,67)$, $(16; 91,67)$, $(20; 100)$. Relier par segments.
\item \begin{enumerate}
\item Graphiquement : abscisse du point d'ordonnée 50\% sur la courbe FCC $\approx \textbf{10,2}$.
\item Calcul : Intervalle médian $[8\,;\,12[$ ($N/2=30$, ECC préc. 18).
$M_e = 8 + \dfrac{30-18}{22}\times 4 = 8 + \dfrac{48}{22} = 8 + 2,1818... = \textbf{10,18} \approx \textbf{10,2}$.
\end{enumerate}
\item Admis si note $\ge 10$. Dans classe $[8\,;\,12[$, on interpole pour trouver l'effectif $< 10$.
$10$ est à $\frac{10-8}{4} = 0,5$ de la largeur de la classe. Effectif estimé $< 10$ dans cette classe $= 0,5 \times 22 = 11$.
Total $< 10 = 18 + 11 = 29$. Admissibles $= 60 - 29 = 31$.
Pourcentage $= \frac{31}{60}\times 100 \approx \textbf{51,67\%}$.
\item $\bar{x} = \frac{5\times2 + 13\times6 + 22\times10 + 15\times14 + 5\times18}{60} = \frac{608}{60} = \textbf{10,13}$.
$\sum n_i c_i^2 = 5\times4 + 13\times36 + 22\times100 + 15\times196 + 5\times324 = 7248$.
$V = \frac{7248}{60} - 10,133^2 = 120,8 - 102,68 = 18,12$.
$\sigma = \sqrt{18,12} \approx \textbf{4,26}$.
Critères jury : $\bar{x}=10,13 \in [9;11]$ \checkmark ; $\sigma=4,26 \ge 3$ \checkmark.
\textbf{L'épreuve est bien calibrée.}
\end{enumerate}
\end{corrige}

\begin{corrige}[Exercice 11 — Factures d'eau et comparaison de coopératives — problème de synthèse]
\textbf{Barème indicatif (sur 20) :} Partie A : Moyenne (3), Médiane+interprétation (3), Écart-type (3), FCD (2), Seuil 20\% (3). Partie B : Vérification moyennes (2), Variances/Écarts-types (3), Choix coopérative + justification (3).
\textbf{Points de vigilance :} Interpolation pour seuil 20\% (FCC 80\% ou FCD 20\%) ; comparaison des dispersions à moyennes égales (ou supposées égales) ; phrase d'interprétation pour la médiane.

\textbf{Partie A.}
\begin{enumerate}
\item $N=150$. Centres : 2,5 ; 7,5 ; 12,5 ; 17,5 ; 22,5 ; 27,5.
$\sum n_i c_i = 18(2,5)+36(7,5)+45(12,5)+27(17,5)+15(22,5)+9(27,5) = 1935$.
$\bar{x} = \frac{1935}{150} = \textbf{12,9 k FCFA}$.
\item $N/2 = 75$. ECC : 18 ; 54 ; \textbf{99}... Intervalle médian $[10\,;\,15[$.
$M_e = 10 + \dfrac{75-54}{45}\times 5 = 10 + \dfrac{105}{45} = 10 + 2,333... = \textbf{12,33 k FCFA}$.
\textbf{Interprétation :} La moitié des abonnés paient une facture mensuelle inférieure à 12,33 milliers de FCFA, et l'autre moitié paie une facture supérieure à ce montant.
\item $\sum n_i c_i^2 = 18(6,25)+36(56,25)+45(156,25)+27(306,25)+15(506,25)+9(756,25) = 31\,837,5$.
$V = \frac{31\,837,5}{150} - 12,9^2 = 212,25 - 166,41 = 45,84$.
$\sigma = \sqrt{45,84} \approx \textbf{6,8 k FCFA}$ (arrondi 0,1).
\item Tableau FCD (\%):
\begin{center}
\begin{tabular}{|c|cccccc|}\hline
Montant & $[0\,;\,5[$ & $[5\,;\,10[$ & $[10\,;\,15[$ & $[15\,;\,20[$ & $[20\,;\,25[$ & $[25\,;\,30[$ \\ \hline
ECD & 150 & 132 & 96 & 51 & 24 & 9 \\
FCD (\%) & 100 & 88 & 64 & 34 & 16 & 6 \\ \hline
\end{tabular}
\end{center}
\item Seuil $s$ tel que 20\% paient plus $\Leftrightarrow$ 80\% paient $\le s$ (FCC = 80\%).
FCC : 12\%, 36\%, \textbf{66\%}, \textbf{84\%}... La classe est $[15\,;\,20[$.
$L=15$, $a=5$, $n=27$, $N_{\text{prec}}=99$, ECC visée $= 0,80\times150 = 120$.
$s = 15 + \dfrac{120-99}{27}\times 5 = 15 + \dfrac{105}{27} = 15 + 3,888... = \textbf{18,9 k FCFA}$.
20\% des abonnés paient plus de 18,9 milliers de FCFA.
\end{enumerate}

\textbf{Partie B.}
\begin{enumerate}
\setcounter{enumi}{5}
\item \textbf{Vérification :}
$C_1$ : $N_1=50$. $\sum n_i c_i = 4(35)+12(45)+18(55)+10(65)+6(75) = 140+540+990+650+450 = 2770$.
$\bar{x}_1 = 2770/50 = \textbf{55,5 sacs}$.
$C_2$ : $N_2=50$. $\sum n_i c_i = 10(35)+16(45)+8(55)+9(65)+7(75) = 350+720+440+585+525 = 2620$.
$\bar{x}_2 = 2620/50 = \textbf{52,4 sacs}$.
\textbf{Les deux moyennes ne sont pas égales à 50 sacs, ni égales entre elles.} (L'énoncé contient une incohérence sur les effectifs ou la moyenne annoncée). Nous poursuivons avec les valeurs réelles calculées.
\item Variances (formule $V = \frac{\sum n_i c_i^2}{N} - \bar{x}^2$) :
$C_1$ : $\sum n_i c_i^2 = 4(1225)+12(2025)+18(3025)+10(4225)+6(5625) = 4900+24300+54450+42250+33750 = 159\,700$.
$V_1 = 159700/50 - 55,5^2 = 3194 - 3080,25 = 113,75$. $\sigma_1 \approx \textbf{10,7 sacs}$.
$C_2$ : $\sum n_i c_i^2 = 10(1225)+16(2025)+8(3025)+9(4225)+7(5625) = 12250+32400+24200+38025+39375 = 146\,250$.
$V_2 = 146250/50 - 52,4^2 = 2925 - 2745,76 = 179,24$. $\sigma_2 \approx \textbf{13,4 sacs}$.
\item \textbf{La coopérative $C_1$ est plus homogène} car son écart-type ($\approx 10,7$) est plus faible que celui de $C_2$ ($\approx 13,4$). Les productions de ses membres sont moins dispersées autour de la moyenne.
\textbf{Pourquoi la moyenne ne suffit pas :} La moyenne ne donne qu'une information de \textbf{position} (centre de la distribution). Elle ne renseigne pas sur la \textbf{dispersion} (écart des valeurs autour du centre). Deux séries peuvent avoir la même moyenne (ou des moyennes proches) mais des dispersions très différentes, conduisant à des réalités économiques distinctes (ici, $C_1$ a des productions plus régulières, moins risquées).
\end{enumerate}
\end{corrige}

\newpage

\coursec{Fiche bilan}

\begin{popfigure}
\centering
\begin{tikzpicture}[
    mindmap,
    concept color=popblue,
    level 1 concept/.append style={level distance=4.5cm, sibling angle=60},
    level 2 concept/.append style={level distance=3cm, sibling angle=45},
    every node/.style={font=\small, text=popdark},
    grow cyclic,
    text width=2.8cm,
    align=center
]
\node[concept, text width=3.5cm] {Séries regroupées\\ en classes}
    child[concept color=poporange] {
        node[concept, align=center] {1. Regroupement\\ \& Effectifs cumulés}
        child { node[rectangle, draw=poporange, fill=poporangeL, rounded corners, inner sep=2pt] {Classes, bornes, amplitude} }
        child { node[rectangle, draw=poporange, fill=poporangeL, rounded corners, inner sep=2pt] {Effectif total $N$} }
        child { node[rectangle, draw=poporange, fill=poporangeL, rounded corners, inner sep=2pt] {Effectifs cumulés croissants/décroissants} }
        child { node[rectangle, draw=poporange, fill=poporangeL, rounded corners, inner sep=2pt] {Fréquences cumulées} }
    }
    child[concept color=popgreen] {
        node[concept, align=center] {2. Position:\\ Mode \& Moyenne}
        child { node[rectangle, draw=popgreen, fill=popgreenL, rounded corners, inner sep=2pt] {Classe modale} }
        child { node[rectangle, draw=popgreen, fill=popgreenL, rounded corners, inner sep=2pt] {Mode (centre classe modale)} }
        child { node[rectangle, draw=popgreen, fill=popgreenL, rounded corners, inner sep=2pt] {Moyenne $\bar{x}=\frac{\sum n_i x_i}{N}$} }
    }
    child[concept color=poppurple] {
        node[concept] {3. Médiane}
        child { node[rectangle, draw=poppurple, fill=poppurpleL, rounded corners, inner sep=2pt] {Intervalle médian} }
        child { node[rectangle, draw=poppurple, fill=poppurpleL, rounded corners, inner sep=2pt] {Interpolation linéaire} }
        child { node[rectangle, draw=poppurple, fill=poppurpleL, rounded corners, inner sep=2pt] {Formule $M_e = L + \frac{\frac{N}{2}-F_{ant}}{f_m} \times a$} }
    }
    child[concept color=poppink] {
        node[concept] {4. Dispersion}
        child { node[rectangle, draw=poppink, fill=poppinkL, rounded corners, inner sep=2pt] {Écart moyen $E=\frac{\sum n_i|x_i-\bar{x}|}{N}$} }
        child { node[rectangle, draw=poppink, fill=poppinkL, rounded corners, inner sep=2pt] {Variance $V=\frac{\sum n_i(x_i-\bar{x})^2}{N}$} }
        child { node[rectangle, draw=poppink, fill=poppinkL, rounded corners, inner sep=2pt] {Écart-type $\sigma=\sqrt{V}$} }
        child { node[rectangle, draw=poppink, fill=poppinkL, rounded corners, inner sep=2pt] {Formule de König-Huyghens} }
    }
    child[concept color=popblue] {
        node[concept] {5. Graphiques}
        child { node[rectangle, draw=popblue, fill=popblueL, rounded corners, inner sep=2pt] {Histogramme (aires $\propto$ effectifs)} }
        child { node[rectangle, draw=popblue, fill=popblueL, rounded corners, inner sep=2pt] {Courbe effectifs cumulés croissants} }
        child { node[rectangle, draw=popblue, fill=popblueL, rounded corners, inner sep=2pt] {Courbe fréquences cumulées} }
        child { node[rectangle, draw=popblue, fill=popblueL, rounded corners, inner sep=2pt] {Lecture graphique médiane, quartiles} }
    }
    child[concept color=popgold] {
        node[concept] {6. Interprétation}
        child { node[rectangle, draw=popgold, fill=popgoldL, rounded corners, inner sep=2pt] {Comparer deux séries} }
        child { node[rectangle, draw=popgold, fill=popgoldL, rounded corners, inner sep=2pt] {Analyser concentration/dispersion} }
        child { node[rectangle, draw=popgold, fill=popgoldL, rounded corners, inner sep=2pt] {Contexte conso, transport, télécom} }
    };
\end{tikzpicture}
\legende{Carte mentale : Statistiques -- Séries regroupées en classes (1ère C/D/E/TI)}
\end{popfigure}

\begin{bilan}
\courssub{Définitions clés}
\begin{itemize}
    \item \textbf{Classe} : intervalle $[a; b[$ (ou $]a; b]$) de longueur $a = b-a$ (amplitude). On note $x_i = \frac{a+b}{2}$ le \cle{centre de la classe}.
    \item \textbf{Effectif $n_i$} : nombre de valeurs dans la classe $i$. Effectif total $N = \sum n_i$.
    \item \textbf{Fréquence $f_i$} : $f_i = \frac{n_i}{N}$ ; $\sum f_i = 1$.
    \item \textbf{Effectif cumulé croissant $N_i$} : somme des effectifs des classes $\le$ classe $i$. $N_{\text{dernière}} = N$.
    \item \textbf{Effectif cumulé décroissant $N'_i$} : somme des effectifs des classes $\ge$ classe $i$.
    \item \textbf{Fréquence cumulée} : $F_i = \frac{N_i}{N}$ (croissante) ou $F'_i = \frac{N'_i}{N}$ (décroissante).
\end{itemize}

\courssub{Formules à connaître par cœur (Probatoire)}
\begin{center}
\begin{tabularx}{\linewidth}{|>{\bfseries}l|X|}\hline
\textbf{Paramètre} & \textbf{Formule (série regroupée)} \\\hline
Moyenne $\bar{x}$ & $\bar{x} = \dfrac{\sum\limits_{i=1}^{k} n_i x_i}{N} = \sum\limits_{i=1}^{k} f_i x_i$ \\\hline
Mode $M_o$ & Centre de la classe modale (classe de plus grand effectif). \\\hline
Médiane $M_e$ & \textbf{Interpolation linéaire} dans l'intervalle médian $[L; L+a[$ :\\[2pt]
& $M_e = L + \dfrac{\frac{N}{2} - F_{\text{ant}}}{f_m} \times a$ \\[4pt]
& $L$ = borne inf. classe médiane, $a$ = amplitude, $f_m$ = effectif classe médiane, $F_{\text{ant}}$ = effectif cumulé avant. \\\hline
Écart moyen $E$ & $E = \dfrac{\sum n_i |x_i - \bar{x}|}{N}$ \\\hline
Variance $V$ & $V = \dfrac{\sum n_i (x_i - \bar{x})^2}{N} = \dfrac{\sum n_i x_i^2}{N} - \bar{x}^2$ \quad (König-Huyghens) \\\hline
Écart-type $\sigma$ & $\sigma = \sqrt{V}$ \\\hline
\end{tabularx}
\end{center}

\courssub{Méthodes express}
\begin{enumerate}
    \item \textbf{Moyenne} : Tableau avec colonnes $x_i$, $n_i$, $n_i x_i$ (et $n_i x_i^2$ pour $V$). Somme / $N$.
    \item \textbf{Classe modale / Mode} : Repérer le max des $n_i$ ; mode = centre de cette classe.
    \item \textbf{Médiane} :
    \begin{enumerate}
        \item Calculer $N/2$.
        \item Trouver l'intervalle médian : 1ère classe où $N_i \ge N/2$ (croissant).
        \item Appliquer la formule d'interpolation.
    \end{enumerate}
    \item \textbf{Histogramme} : Axe des abscisses = bornes des classes. Axe des ordonnées = \textbf{densité} $d_i = \frac{n_i}{a_i}$ (ou $f_i/a_i$). Aire du rectangle = effectif (ou fréquence).
    \item \textbf{Courbes cumulées} : Points $(borne\_sup\_classe, N_i)$ pour croissante ; $(borne\_inf\_classe, N'_i)$ pour décroissante. Relier par segments. Médiane = abscisse de l'intersection des deux courbes (ou $y=N/2$ sur croissante).
\end{enumerate}
\end{bilan}

\begin{aretenir}[Checklist avant le Probatoire]
\begin{itemize}
    \item $\square$ Savoir construire un tableau de classes (choix amplitude, bornes) à partir de données brutes.
    \item $\square$ Calculer les effectifs et fréquences cumulés (croissants ET décroissants) sans erreur.
    \item $\square$ Appliquer la formule de la moyenne $\bar{x} = \frac{\sum n_i x_i}{N}$ en utilisant les \textbf{centres} des classes.
    \item $\square$ Identifier la classe modale et donner le mode (centre de cette classe).
    \item $\square$ Déterminer l'intervalle médian grâce aux effectifs cumulés croissants ($N_i \ge N/2$).
    \item $\square$ Appliquer \textbf{exactement} la formule d'interpolation linéaire de la médiane (noter $L, a, f_m, F_{\text{ant}}$).
    \item $\square$ Calculer la variance avec la formule de König-Huyghens ($V = \frac{\sum n_i x_i^2}{N} - \bar{x}^2$) : plus rapide et moins d'erreurs.
    \item $\square$ Passer de la variance à l'écart-type $\sigma = \sqrt{V}$ et interpréter : $\sigma$ faible $\Leftrightarrow$ valeurs groupées autour de $\bar{x}$.
    \item $\square$ Construire un \textbf{histogramme} correct : aire des rectangles proportionnelle aux effectifs $\Rightarrow$ ordonnée = densité $n_i/a_i$.
    \item $\square$ Tracer et lire les courbes cumulées (croissante et décroissante) pour retrouver la médiane graphiquement.
    \item $\square$ Interpréter les résultats dans le contexte (ex: « La moitié des clients consomment moins de $M_e$ kWh »).
    \item $\square$ Vérifier la cohérence : $0 \le f_i \le 1$, $\sum f_i = 1$, $M_e \in [L; L+a[$, $\sigma \ge 0$.
\end{itemize}
\end{aretenir}

\findecours

\end{document}
